% =====================================================================
%  R. Nielsen, MATHEMATIK OG DIALEKTIK. En philosophisk Afhandling.
%  Kjøbenhavn: Gyldendal (F. Hegel), Thieles Bogtrykkeri, 1859.  VII + 148 pp.  Fraktur; formulas in antiqua.
%
%  TRANSCRIPTION -- GENERATED; DO NOT EDIT BY HAND.  Rebuild, chained (\&\&), from the repo root:
%    python3 ebook-method/mathematik-og-dialektik/draft.py --emit --force
%      this template + the scan's text layer (Google's OCR in the Google Books scan of the NYPL copy,
%      ~/bibliotek/Nielsen, Rasmus/Mathematik_og_Dialektik.pdf, see SCANS.tsv);
%    python3 ebook-method/collate/check.py mathematik-og-dialektik --witness ocr --apply
%      the decisions in decisions.tsv, taken against tesseract's independent reading of the images;
%    python3 ebook-method/mathematik-og-dialektik/oepass.py --apply
%      the ø pass (little to do: Google's layer reads ø);
%    python3 ebook-method/mathematik-og-dialektik/finish.py --apply
%      the page readers' formulas, footnotes, headings, letterspacing and antiqua; the page fixes
%      (pagefix/); the remaining OCR faults; the book's own Rettelser marked PRINTED AS IS.
%
%  WORD-ACCURACY (2026-10-06).  No e-book and no earlier transcription exist.  The draft (Google's
%  layer, good on the Fraktur, garbled in the formulas) was collated word by word against tesseract
%  (tessdata_best script/Fraktur; a third reading, deu_latf, voting where they differ): 4600
%  disagreements.  About 2100 were settled by rule (rules.py: each OCR's known glyph faults undone,
%  accepted only where the result is a word attested in at least two other books of the corpus).
%  Then every page of pp. 3-148 was read whole at the image by nine page readers (Sonnet,
%  page-reads.tsv), each giving for its pages: the open words (608 decided from these readings, a
%  few by hand), every formula in LaTeX, every footnote whole, headings, letterspacing, antiqua and
%  oddities.  finish.py sets the formulas and footnotes from these readings (904 formulas placed by
%  script).  Where it could not, 95 pages were rewritten whole by five text-only page-fix agents from
%  the readers' lines, keeping the Danish words (every word they changed is listed in
%  pagefix/CHANGES.tsv; a word-loss check, old against new, was read in full by the main session);
%  then two agents at the images settled 32 pages with doubtful spots and leftovers (note text
%  standing in the page text, OCR'd Greek, reader comments).  An unattested-word sweep over the
%  finished text found 32 more OCR faults (Fraktur k/f, s/f, Q/O, Ø/D: WORDFIX in finish.py).
%  The book's Rettelser (16 entries) were checked in the text: the misprints stand as printed and
%  are marked PRINTED AS IS (four that a reader or a rule had mended were put back).  Front matter
%  (title, motto, Forord, Indhold, Rettelser) read at the image by a reader; Forord pp. IV-V are
%  missing in the NYPL copy and were read by the main session from the second copy (Google Books
%  X3haAAAAcAAJ).
%  NOT ESTABLISHED: misprints that the OCRs and the readers all read as real words, outside the
%  Rettelser; whether the u for n in Tænkuing, Logikeu (p. 21), Auskuelse (p. 146), Ueudeligheden
%  (p. 147) is a turned letter in the print (the image readers say so; kept and marked); the
%  formulas rest on one reader each (no second reading), with the readers' own doubts in the ODD
%  lines of page-reads.tsv; letterspacing beyond what the readers reported.
%  PAGE MARKERS: 146 of 146 placed; check.py --pages: 129 exact, the other 17 begin with a formula
%  or follow a footnote running over from the page before (witness order), checked by hand.
%
%  CONVENTIONS.  Body in Fraktur; aa (not å); ø as printed.  Emphasis is letterspacing -> \emph{}.
%  Latin, French, English set in antiqua -> \textit{}.  Greek in Greek type (textalpha).  Formulas as
%  printed, in LaTeX ($x$ for a single variable); a misprint in a formula stays, with a comment.
%  Fraktur I/J are one sort: I for I-words, J for J-words.  The double hyphen inside a compound is
%  "-".  "% --- p. N ---" + \opage{N} at the first word of printed page N (Forord: roman).  Footnotes
%  1), 2) per page, as printed; a note running over to the next page is set whole at its mark, with
%  a % comment on the page it runs onto.  Folios, signature marks and the library's stamps are not
%  transcribed.  Print governs; misprints are kept and marked PRINTED AS IS.
% =====================================================================
\documentclass[12pt,a4paper,openany]{book}
\usepackage[utf8]{inputenc}
\usepackage[T1]{fontenc}
\usepackage{amsmath}
\usepackage{amssymb}
\usepackage{libertinus}
\usepackage{libertinust1math}
\usepackage{textalpha}   % polytonic Greek typed directly
\usepackage[danish]{babel}
\usepackage[protrusion=true,expansion=true]{microtype}
\usepackage{setspace}
\usepackage[margin=1.3in]{geometry}
\usepackage{fancyhdr}
\usepackage{hyperref}
\hypersetup{pdftitle={Mathematik og Dialektik}, pdfauthor={Rasmus Nielsen}, hidelinks}
\setlength{\headheight}{15pt}
\pagestyle{fancy}
\fancyhf{}
\fancyhead[LE,RO]{\thepage}
\fancyhead[C]{\textit{R. Nielsen: Mathematik og Dialektik (1859)}}
\renewcommand{\thefootnote}{\arabic{footnote})}
\setstretch{1.3}
\usepackage{marginnote}
\newcommand{\opage}[1]{\marginnote{\footnotesize\textit{[#1]}}}
\newcommand{\apage}[1]{\marginnote{\footnotesize\textit{[#1]}}}

\begin{document}

% --- front matter, read at the image 2026-10-06 (front-reads.txt; Forord pp. IV-V: see the header).
% PDF 9, title.  "En philosophisk Afhandling" and "Thieles Bogtrykkeri." are letterspaced.  The short
% rule under the title is drawn; the publisher's monogram seal and the library's marks are not transcribed.
\begin{titlepage}
\begin{center}
\vspace*{3em}
{\Huge Mathematik og Dialektik.}\\[1em]
\rule{4em}{0.4pt}\\[1.5em]
{\large \emph{En philosophisk Afhandling}}\\[1.5em]
af\\[1em]
{\large R. Nielsen,}\\[0.3em]
Professor i Philosophien.\\[4em]
{\large Kjøbenhavn.}\\[0.5em]
Forlagt af den Gyldendalske Boghandling (F. Hegel).\\[0.5em]
{\small \emph{Thieles Bogtrykkeri.}}\\[0.5em]
1859.
\end{center}
\end{titlepage}
% PDF 10, the motto, in small type lower on the page; the attribution set to the right.
\thispagestyle{empty}
\vspace*{14em}
{\small
\noindent Hvad der er Livsprincipet i Videnskaben, bør aldrig, af misforstaaet Hensyn
paa Populariteten, opoffres eller endog blot tilhylles; hvo som ikke har Evne og
Villie til at fatte de philosophiske Gjenstande i deres Sandhed, maa hellere reent
lade være at beskjæftige sig med dem.\par
\hfill J. L. \emph{Heiberg.}\par}
\clearpage
% ===== FORORD (pp. I-VII; PDF 11-15, pp. IV-V from the second copy) =====
% --- p. I ---  (Forord, PDF 11; no folio printed)
\begin{center}{\large Forord.}\end{center}

\opage{I}At tilveiebringe god Forstaaelse mellem Mathematik og
Philosophie er endda ikke saa ganske let, skjøndt man ifølge
Platos Vink næsten skulde mene, at Tingen gik som af sig selv.

Den platoniske Anskuelse, at Mathematiken væsentlig
maa betragtes som en Forskole for Dialektiken, er især udviklet
i Læren om Staten (6te og 7de Bog). Det Mathematiske,
der ligger paa Grændsen af det Synlige og det
Usynlige, egner sig just til at fremhjælpe Abstractionen. „De,
der beskjæftige sig med Geometrie og Arithmetik og andre
mathematiske Videnskaber, forudsætte i enhver Beviisførelse
visse Ting, som det Lige og det Ulige, Figurer, tre Slags
Vinkler og andet hermed Beslægtet. Sjælen er ved disse
Undersøgelser nødt til at betjene sig af Forudsætninger, ikke
for at komme til den første Grund, thi den kan ikke gaae ud
over disse Forudsætninger, men den bruger Billeder laante
fra de jordiske og sandselige Ting, og vælger iblandt disse
dem, der maae ansees for at have størst Tydelighed og maae agtes
høist“. „Er Talen saaledes om et Qvadrat og dets Diagonal,
føre de ikke Beviset for dets Skyld, som de tegne, men de
have Qvadratet selv i Tankerne“; de betjene sig ved alt
Andet, hvad de danne og tegne, ligeledes kun af Billeder,
„idet de tragte efter at erkjende hiint selv, som kun sees med
Forstanden“. Tallet og Eenheden høre endvidere til det i
Gjenstandene, der særlig opfordrer til Tænkning. „Thi ifald
% --- p. II ---
\opage{II}Eenheden i og for sig selv kunde tilstrækkelig opfattes ved
Synet eller ved en anden Sands, saa kunde den vel ikke
føre os til det Værende ...; men, naar den aldrig sees uden
med en Modsætning af den, saa at ingen Ting mere viser
sig som een, end som det Modsatte, da er vistnok en Dommer
nødvendig, der kan afgjøre Sagen, og Sjælen maa nødvendig
blive tvivlraadig, og sætte alle sine Tanker i Bevægelse for
igjennem en Undersøgelse at finde Svar paa det Spørgsmaal,
hvad dog Eenheden i og for sig er“. --- „Regne-
og Tællekunsten, der er nødvendig saavel for Krigeren, som for
Philosophen, egner sig derfor ganske til ved Lov at indføres
i Staten, og vi bør overtale dem, der skulle tage Deel i de
høieste Statsforretninger, at give sig af med Regnekunsten,
ikke paa den almindelige Maade, men saaledes, at de ved
den rene Tænkning komme til Anskuelse af Tallenes Natur,
og at lære den, ikke for som Kjøbmænd og Kræmmere at
bruge den ved Kjøb og Salg, men deels for at anvende den i
Krigen, og deels for at gjøre det lettere for Sjælen at hæve
sig fra det Forgjængelige og Foranderlige til det Sande og
Værende“.

Denne Platos Opfattelse af Mathematiken er jo vistnok
noget eensidig, thi Mathematiken er ikke for Krigens Skyld
og ikke for Philosophiens Skyld; den er oprindelig for sin
egen Skyld. Men dersom Mathematiken, fordi den er og
maa være en af Philosophien uafhængig Videnskab, nu ogsaa
vilde gjøre sig uafhængig af det Almeenvidenskabelige, og
derved løsrive sig fra den humane Aandsretning, vilde Eensidigheden
dog vistnok blive endnu større og derhos mere
betænkelig.

Hvorledes man skal drive det, naar man ret vil
drive det til at blive Stokmathematiker, har en fransk Forfatter
fra det 18de Aarhundrede (\textit{Saverien: Histoire du
Calcul des Infiniment-Petit}) angivet i træffende Charakteristik.
Talen er om Maclaurins Indsigelser imod at betragte det
uendelig Lille og det uendelig Store som Størrelser, eftersom
% --- p. III ---
\opage{III}Overgangen fra det Endelige til det Uendelige er aldeles
ubegribelig. „Imidlertid“ --- det er en Synsmaade, Forfatteren
tænker sig indført --- „ere jo dog de Conseqvenser,
man udleder af hine ubegribelige Forudsætninger, reelle
Sandheder. Og hvad gjør saa det til Sagen, at det opstillede
Princip er ubegribeligt, naar det kun leder os til at
opdage en Sandhed?“ Denne Synsmaade bliver dog afviist,
og det af følgende Grunde. „For det Første“, svarer Forfatteren,
„kunne vi, naar Principet er ubekjendt, jo ikke vide,
om ogsaa alle deraf følgende Conseqvenser lade sig indrømme:
vi gaae som i Blinde, og det turde hænde sig, at vi efter
møisommelig Anstrengelse tilsidst stødte paa Noget, som netop
udviste Principets Natur, d. v. s. paa Noget, vi ikke kunde
begribe. For det Andet“, spørger han: „er det dog
ikke en skammelig Nedværdigelse af Fornuften, at lade den
følge en Sammenkjædning af Led, en Række, hvis Begyndelse
og Ende, den ikke kan gjennemskue? Det er et stort Onde,
der fødes af den Slags Routine: vant til at operere uden
at indsee Grunden taber Aanden al Øvelse i at tænke, og
ender med at udføre Alt ganske maskinmæssig. \textit{Aussi voit-on
les jeunes Gens les plus incapables de refléxions devenir
facilement habiles Calculateurs ... Comme on croit que
cette aptitude est un fruit du génie on fait sonner bien
haut ce succès; et un enfant se trouve grand Homme,
sans qu’ il se soit encore avisé de penser. Les éloges
qu’on lui croit dûs et qu’on ne cesse de lui prodiguer, le
persuadent. En faut-il davantage pour faire éclore dans
lui l’orgueil, enfant chéri de l’ignorance, et pour le rendre
dans la suite un homme vain presque inutile à la Société
et toujours insupportable}“.

Vistnok er det endda ikke saa sjældent, at en „\textit{enfant
chéri de l’ignorance}“, en saadan uopdragen lille „\textit{grand
Homme}“ siden hen arbeider sig op til at blive noget Mere
end „\textit{habile Calculateur}“, til at blive Mathematiker, endog
dygtig Mathematiker og virksomt Medlem i Samfundet; men
% --- p. IV --- (pp. IV-V are missing in the NYPL copy scanned; read at the image of the second copy, Google Books X3haAAAAcAAJ, PDF 16-17)
\opage{IV}Indtryk fra Drengeaarene ere varige, og Mangel af god
Opdragelse ikke let at erstatte. Bliver hiin \textit{orgueil} i Saveriens
mørke Billede under begunstigende Omstændigheder
behørig udklækket til mathematisk Brøsighed, vil Samfundets
Dygtige dog bestandig have Intelligensen i nogen Raahed,
og da naturligviis, naar Leilighed gives, fare frem mod
Philosophien i megen Blæst.

For da at modstaae mathematiske Hvirvelvinde, naar
det blæser fra det raae Hjørne, og Blæsten tager til, maa
Philosophien see at sikkre sig nogle Grundstøtter, der ikke let
blæse omkuld, d. v. s. nogle Forudsætninger, som ingen billig
og sagkyndig Læser vil negte sin Anerkjendelse.

Den første Forudsætning er, at Mathematikeren, der
vil gjøre Indsigelser imod den Maade, hvorpaa en Philosopherende
har opfattet og anvendt mathematiske Formler og
Udtryk, uleiliger sig med at see saameget paa Tankegangen,
at han i slige Udtryk og Formler ikke indlægger Andet eller
Mere, end hvad den Philosopherende udtrykkelig af den Sammenhæng,
hvori de forekomme, sees at have villet lægge i dem.
Hvor let det er, naar man opfatter Sætninger, der ifølge
Sammenhængen kun gjælde med en bestemt Indskrænkning,
som om de vare sagte i reen Almindelighed, at gjøre endog
en Mathematikens Heros til Ignorant i Mathematiken, er
det da heller ikke vanskeligt at bevise. Det hedder i Platos
Timæus: „Men naar nu Alverdens Legeme maatte være
blevet en Flade uden Dybde, saa vilde een Midte have
været tilstrækkelig til at forbinde de to andre Dele og sig
med dem. Men da det skulde have det Fastes Beskaffenhed,
og det Faste aldrig forbindes ved een men ved to Midter,
saa har Gud imellem Ild og Jord sat Vand og Luft, og
saavidt muligt bragt dem i et saadant Forhold til hinanden,
at som Ilden forholder sig til Luften, saaledes forholder
Luften sig til Vandet, og som Luften til Vandet, saaledes
Vandet til Jorden“. Om denne Udvikling tilføies der saa i
C. J. Heises Anmærkning (til Oversættelsen, S. 28), at den
% --- p. V ---
\opage{V}er „grundet paa følgende to mathematiske Sætninger: 1) \emph{To
Flader eller saadannes Arealer kunne som Yderled
i en Proportion forbindes ved een Flade eller
eet Areal som Mellemled, og 2) To Legemer eller
sammes Voluminer kunne som Yderled ikke forbindes
ved eet Legeme eller eet Volumen, men
fordre to forskjellige saadanne som Mellemled}“.
Da man nu, hvad allerede Proklos har bemærket, dog alligevel
altid imellem to „Flader“ kan sætte to eller flere geometriske
Mellemled, og ligeledes mellem to „Legemer“ indskyde
eet geometrisk Mellemled: vil det jo --- paa „Kritikens“
nærværende Standpunkt --- være let, endog for en middelmaadig
Mathematiker temmelig let, at gribe Plato i „de
groveste Feil“, hvis Opgaven væsentlig er at gjøre Blæst.
Er Opgaven derimod ikke at gjøre Blæst, men at forstaae
Plato, maa her unegtelig gjøres adskillige Indskrænkninger,
idet her ifølge Sammenhængen maa forudsættes, at Plato
ikke har tænkt paa Problemernes Løsning ved Construction,
men ved Beregning, idet han har forudsat, at det søgte
Mellemled skal være rationalt, og de paagjældende Yderled
Qvadrater og Kuber af primiske Tal; men med disse Indskrænkninger
ere Sætningerne da ogsaa aldeles correcte.

Den næste Forudsætning er, at Mathematikeren, dersom
han i en philosophisk Afhandling her og der støder paa virkelige
Feil, Unøiagtigheder og Misforstaaelser, dog derfor ikke
strax frakjender den Philosopherende, der har kunnet begaae
slige Feil, al Forberedelse og Evne til at sætte sig ind
i en mathematisk Tankegang. Vistnok gjør en saadan
Umyndighedserklæring altid sin øieblikkelige Virkning, naar
Opgaven er at hovmesterere. Men, naar Opgaven ikke er
at hovmesterere, men at udfinde det rette Forhold mellem
Indsigt og Feiltagelse, vil den herved paaankede Fordømmelsesdom
jo langt bedre passe sig for en skrydende Rabulist, end
for den i Sandhed Overlegne, der virkelig seer, hvorpaa det
kommer an. Tør man fra Feil i Behandlingen af en
% --- p. VI ---
\opage{VI}mathematisk Formel eller fra Misforstaaelse af et enkelt
Udtryk øieblikkelig og uden videre slutte til Ukyndighed i det
Mathematiske: da blive her mange Ukyndige, da ere D’Alembert
og Euler jo ogsaa Ukyndige, ja, da er Cartesius en
Stymper, og Nutidens „strenge Kritik“ tør svinge Corporalstaven
over Newton og Leibnitz.

En tredie Forudsætning er, at den Mathematiker, der
vil indlade sig paa at bedømme det i Mathematiken Philosophiske,
allerførst maa gjøre sig den Uleilighed at sætte sig
nogenledes ind i hele den philosophiske Tankegang, hvoraf
Behandlingen særlig fremgaaer, og derhos vide at holde
Mathematikens og Philosophiens forskjellige Synsmaader
saaledes ud fra hinanden, at de behørig kunne sammenstilles,
og opklares ved hinanden indbyrdes. Vistnok er
det af Vigtighed, at Unøiagtigheder paavises og begaaede
Feil rettes; men der er Et, der i videnskabelig Forstand er
endnu vigtigere, og det er, at der udtrykkelig skjelnes imellem
væsentlige og uvæsentlige Feiltagelser, en Adskillelse, der i
Undersøgelsen angaaende Forholdet imellem Philosophie og
Mathematik er saameget mere uafviselig, som ingen videnskabelig
Dannet vil kunne negte, at Meget, naar det
sees fra Philosophiens Side, og tilmed i et Udkast, en
skizzeret Fremstilling, mere antydes end udvikles, kan være
aldeles uvæsentligt, medens det i en mathematisk og derhos
udførlig Lærebog vilde være endog meget væsentligt. Ved
enten slet ikke at agte paa denne Forskjel, eller bedre: ved
udtrykkelig at indrømme den i Ord, for siden under Bedømmelsen
selv at fornegte den i Gjerning, er det en let
Sag at forstørre uvæsentlige Feil, en ringe Kunst at gjøre
Philosophiens Myg til mathematiske Elephanter og imponere
Læseren ved mange Elephanter.

At slige Forudsætninger, som kun med Hensyn paa det
Videnskabelige ere ufravigelige, naturligviis ikke tør være
generende for Saadanne, som ellers fra deres eget Standpunkt
maatte føle sig tilskyndede til at betænke nærværende
% --- p. VII ---
\opage{VII}Afhandling med „en streng Kritik“, forstaaer sig vel af sig
selv. Er Ævret opgivet og Markleddet aabent, da maatte
den Forfatter i Sandhed være meget naiv, der vilde
mene, at nogle Sætninger, opstillede i et Forord, kunde
være ham enten Værn mod de Løse, eller Hegn mod de
Brølende, mod dem paa Literaturens Overdrev, som ville
stange; nei, de opstillede Sætninger skulle ingenlunde foreskrive
Pligter for Mathematikens Hornqvæg, de skulle kun
være forpligtende for Forfatteren selv og for dem, der med
ham ville noget Videnskabeligt.

Skulde en Bedømmelse af foreliggende Afhandling,
ved at holde sig til de opstillede Forudsætninger, enten
virkelig eller formeentlig kunne paavise betænkelige Feil
og Misgreb: erklærer Forfatteren sig herved pligtig til at
staae den navngivne Recensent, han høre nu forøvrigt blandt
de Yngre eller blandt de Ældre, blandt Autoriteter eller
Ikke-Autoriteter, paa enhver Maade til Regnskab, at Sagen
dog tilsidst maa kunne klares; hvorimod --- og dette bliver
herved ligeledes udtrykkelig fastsat --- ethvert Forsøg til
Kritik, som ved at gjøre Brud paa de nævnte Sætninger,
forvexler Væsentligt med Uvæsentligt, og forvirrer Synspunkterne,
bliver at afvise som ubeføiet, Forsøget komme for
Resten fra en lille eller fra en stor Autoritet, og byde, hvad
det vil, og være som det vil, med grovt Sand eller med
fiint Sand, med ringe Anseelse, eller med Brask og Bram.

\begin{flushleft}\hspace*{2em}Kjøbenhavn i December 1859.\end{flushleft}
\begin{flushright}R. Nielsen.\end{flushright}

\clearpage

% ===== TEXT (pp. 3--148) =====

% --- p. 3 ---
\setcounter{footnote}{0}%
\opage{3}Mathematik og Philosophie ere forskjelligartede Videnskaber;
de maae desaarsag hævde hver sin Selvstændighed. Mathematiken
er en exact, Philosophien en dialektisk Videnskab.
Som exact Videnskab har Mathematiken ingen Brug for
Philosophien; som dialektisk Videnskab har Philosophien derimod
Brug for Mathematiken; den exacte Videnskab forholder
sig til sig selv ved Udelukkelse af andre Videnskaber;
den dialektiske Videnskab forholder sig gjennem andre Videnskaber
til sig selv.\footnote[1]{Jvfr. Philos. Propædeut. i Grundtræk 1857, S. 99 § 22.}

Fra Erkjendelseslærens Standpunkt bliver Philosophiens
Forhold til sig selv gjennem det Mathematiske da i Almindelighed
at bestemme saaledes: Mathematiken fastsætter sine
Begreber i den exacte Form; de saaledes fastsatte Begreber
kunne ikke uden videre overføres i Philosophien; derimod
kan den philosophiske Dialektik ved at bringe de exacte
Forholdsbestemmelser til Bevidsthed vinde i Klarhed og
Skarphed, og saaledes opnaae den Punktlighed, som Tingenes
Erkjendelse nødvendig maa forudsætte. For at paavise,
hvorledes Erkjendelseslæren articulerer den dialektiske Begrebsudvikling
ved Momenter af den exacte Analyse, lade vi
i nærværende Afhandling Fremstillingen dreie sig om tre
Hovedpunkter: I. \emph{Functionsbegrebet}, II. \emph{Uendelighedsproblemet}
og III. Det Uendeliges Analyse.

% --- p. 4 ---
\setcounter{footnote}{0}%
\opage{4}
\begin{center}I.\end{center}
\begin{center}Functionsbegrebet.\end{center}

At Begrebet: „Function“ ikke udelukkende tilhører
Mathematiken, er bekjendt.\footnote[1]{Jvfr. Phil. og Mathem. S. 5—6.} Ligesom Anskuelserne, ifølge Kant,
beroe paa „Affectioner“, beroe Begreberne paa „Functioner“:
„Ich verstehe aber unter Function dieEinheit derHandlung verschiedene
Vorstellungen unter einer gemeinschaftlichen zu ordnen.“
(Kritik der r. V.) Den Handlingens Eenhed, hvorved Kant
charakteriserer Forstandsvirksomheden fra det almeenlogiske
Synspunkt, er, som man let seer, i væsentligt Slægtskab
med den Handlingens Eenhed, hvorved Forstanden opfatter
den mathematiske Relation mellem afhængige og uafhængige
Størrelser; kun at Forskjellen ikke oversees. Udenfor Mathematiken
falder, hvad Functionskategorien angaaer, Eftertrykket
paa den lovbestemte Handling, i Mathematiken derimod paa
det Afhængighedsforhold, hvorved Handlingen (Operationen)
bestemmes; udenfor Mathematiken opfattes Kategorien fornemmelig
fra den active, i Mathematiken, om man vil, fra
den passive eller dog neutrale Side. „En Størrelse, der
afhænger af een eller flere vilkaarlige Størrelser, hvormed
den altsaa kan tænkes forbunden ved en Ligning, kaldes en
Function af denne eller disse. Saaledes er t. Ex. $ax^{3}$ en
Function af $x$, idet $a$ tænkes uforandret, Logarithmen en
Function af Tallet, Sinus (som Linie) en Function af Radius
og Buen o. s. fr.“\footnote[2]{Chr. Jürgensen. Begyndelsesgrundene af den mathematiske Analyse. Kbhvn 1855. S. 18.}

Den særegne Betydning, Kategorien saaledes har erholdt
i Mathematiken, skal nu afgive Bidrag til en philosophisk
Udvikling af samme Kategorie, idet vi undersøge Begrebets
Reqvisiter, dets Forhold til Loven, og dets eiendommelige
Fremstilling.

% --- p. 5 ---
\setcounter{footnote}{0}%
\opage{5}
\begin{center}a) Functionens oprindelige Reqvisiter.\end{center}

Til eiendommelig Bestemmelse af den mathematiske
Function kræves:

\begin{center}α) en \emph{constant} og en \emph{variabel} Størrelse.\end{center}

I Udtrykket: $f(x)=ax$ er $a$ constant, $x$ variabel; i
Udtrykket $f(a)=ax$ er $x$ derimod constant og $a$ variabel. Kan
nu enhver Størrelse ifølge sit Begreb være snart constant,
snart variabel efter Omstændighederne; saa sees allerede heraf,
hvor let Mathematikens Axiom: „enhver Størrelse er lig sig
selv“, kan angribes af Dialektiken. At enhver endelig Størrelse
maa, forsaavidt den er constant, være sig selv lig: $a=a$, er uimodsigeligt; men hvorledes forholder det sig med
den variable Størrelse? Den variable Størrelse maa jo,
ifølge Begrebet, uophørlig kunne afvige fra en hvilkensom
helst given Værdi, og just derved i Ligheden med sig selv
blive sig selv ulig:
\[ x \overset{>}{\underset{<}{=}} x \]
. Vilde man, for at hævde den
opstillede Grundsætnings exacte Gyldighed, opfatte den variable
Størrelse, $x$, som en foreløbig Betegnelse, et collectivt
Udtryk, der siger, at $x$ ved at antage hvilkesomhelst Værdier
opløser sig i en Fleerhed af constante Størrelser: $a=a$, $b=b$, $c=c$ o. s. v.; da havde man herved ikke blot, hvad
der i andre Henseender er mathematisk utilladeligt, gjort den
variable Størrelse til Eet med den vilkaarlige, men endog
aldeles ophævet den variable Størrelses Begreb. At Bestemmelsen:
\emph{variabel} ligesaa væsentlig tilkommer Størrelsen,
som Bestemmelsen \emph{constant}, sees deraf, at Størrelsen
netop defineres som det, der kan forøges og formindskes.\footnote[1]{Alt, hvad der kan tænkes forøget eller formindsket, og som kan opløses i eensartede Dele, kaldes Størrelse (quantum), og siges at besidde Storhed (quantitas). Ramus: Elementær Geometrie. 1850. Indledn.}

% --- p. 6 ---
\setcounter{footnote}{0}%
\opage{6}Vilde man i Tillid til den formelle Logiks øverste Grundsætning
forklare $x=x$ saaledes: det Variable maa, forsaavidt
det ikke er constant, nødvendigviis ligne sig selv deri, at
det er variabelt; da vilde man med samme Ret, som naar man
siger: det Variable er variabelt, altsaa: $x=x$, ogsaa kunne
sige: det uendelige er uendeligt, altsaa: $\infty=\infty$; men imod
en saadan Forklaring maa Mathematiken protestere. Hvad
betyder $\infty$? Dersom $\infty$ ikke betegner en Størrelse, saa er
Ligningen $\infty=\infty$ meningsløs; det maatte da i logisk Almindelighed
hedde: det Uendelige er, hvad det er, en Størrelsesnegation.
Skal det Uendelige derimod opfattes som
Størrelse, da er Ligningen: $\infty=\infty$ incomplet, d. e.
ligesaa ugyldig som gyldig:
\[ \infty \overset{>}{\underset{<}{=}} \infty \]
er det fuldstændige
Udtryk. At lægge Eftertryk paa Forskjellen imellem den
bestemte og den endnu ubestemte Størrelse hjælper ikke.
Er den endelige Størrelse efter sit Begreb bestemt ved Modsætning
til den uendelige, da er den uendelige jo med det
Samme bestemt ved Modsætning til den endelige; ere de
endelige Størrelser at bestemme ved deres Forhold til hverandre
indbyrdes, maa et Tilsvarende jo gjælde om de
uendelige.

Det opstillede Axiom: „enhver Størrelse er lig sig
selv“, har altsaa nu opløst sig i to Sætninger, af hvilke
den første er en Tautologie, den anden en Modsigelse. Det
er nemlig ikke Størrelsen som Størrelse, men Størrelsen
som constant Størrelse, der er lig sig selv; Eftertrykket falder
ikke paa det at være Størrelse, men paa det at være
constant Størrelse. Og hvad er saa en constant Størrelse?
En Størrelse, der er sig selv lig. Altsaa siger Axiomet:
en Størrelse, der er sig selv lig, er sig selv lig; en reen
Tautologie! Opfattes dernæst Størrelsen som det, „der
kan tænkes forøget eller formindsket“, saa falder Eftertrykket
ikke paa det at være Størrelse, men paa det at være variabel
% --- p. 7 ---
\setcounter{footnote}{0}%
\opage{7}Størrelse; det Variable maa vise sig i Varieren; den variable
Størrelse, der ifølge Axiomet jo ogsaa skulde være sig selv
lig, kan altsaa kun ligne sig selv deri, at den bestandig
bliver sig selv ulig, hvilket er en Modsigelse.

Disse Vanskeligheder hæves, naar vi see paa Functionen.
Functionsbegrebet er, dialektisk seet, det sande Størrelsesbegreb,
efterdi Størrelsens dobbeltsidige Bestemmelse, at
være constant, d. e. sig selv lig og dog variabel, d. e. sig
selv ulig, her sættes i Eenhed. Kan man end ikke umiddelbart
sige: $x=x$, forsaavidt $x$ er variabel, saa kan man dog
med Sandhed sige: $f(x)=ax=f(x)$, efterdi den samme
Function maa som Function under alle Forandringer være
sig selv lig.\footnote[1]{Altsaa dog ogsaa $x=x$, forsaavidt $x$ er en Function af $x$.} Idet Tautologien er hævet, er Modsigelsen
med det Samme ophævet. Det Ulige staaer ikke længere
som abstract Negation ligeoverfor det Lige, eller det Variable
ligeoverfor det Constante; Functionen lader just det Lige
reflectere sig i det Ulige, det Variable i det Constante, Forskjelsmomentet
i Identiteten.

Til Functionsbegrebets nødvendige Reqvisiter hører endvidere:

\begin{center}β) \emph{To variable Størrelser, en uafhængig og en afhængig.}\footnote[2]{At en Function som: $y=f(x, z, t, u....)$ indeholder flere Uafhængige, lades i denne Sammenhæng, hvor det blot gjælder en Begrebsudvikling, ude af Betragtningen.}\end{center}

I Ligningen: $f(x)=y=ax$ ere Størrelserne $x$ og $y$ foranderlige, men saaledes, at $x$ er uafhængig, $y$ afhængig.
Forandringen i $x$ forholder sig da til Forandringen i $y$ som
Forudsætning til Conseqvens, som Grund til Følge. Da
Functionsbegrebet netop viser os Forudsætningens og Conseqvensens,
det Uafhængiges og det Afhængiges reflecterede
Eenhed, kan Forskjellen imellem den blotte Forudsætning og
den egentlige Grund i denne Sammenhæng oplyses. Forud-

% --- p. 8 ---
\setcounter{footnote}{0}%
\opage{8}sætningen i dens Umiddelbarhed er et Givet, et Antaget,
der opfattes i Forestillingen, og kun forsaavidt falder under
Reflexionen, som det indsees, at den, for at være Forudsætning,
maa have sine Conseqvenser. Man kan saaledes an-
tage en Forudsætning, fastholde den i Forestillingen som
noget Bekjendt, og dog i høi Grad overraskes af dens
Conseqvenser. Overraskelsen hidrører derfra, at Forudsæt-
ningens Væsen ikke fra Begyndelsen af blev gjennemskuet;
Forudsætningen blev vel tydeliggjort i en Forestilling, og
forsaavidt opfattet fra en vis Side, men ikke alsidig opklaret,
ikke opfattet efter sit Begreb. Først naar vi gjennemskue
Forudsætningen, og opdage de deri skjulte Conseqvenser,
erkjende vi Forudsætningens Væsen, sætte den som
Grund, og Conseqvensen som den begrundede Følge. Hvilken
Forskjel det gjør, om man kun har Grunden som en umiddelbar
Forudsætning, eller man i Forudsætningens Væsen
seer Grunden, kan med Hensyn paa $f(x)=ax=y$ oplyses
ved et saare simpelt Exempel. A spørger B: Hvad vil Du
have for Din Jagttaske? „3 Rdlr.!“ er Svaret. „3 Rdlr.“,
vedbliver A, „vil jeg give Dig strax, og dersom Du vil bie
i otte Dage, indtil en Pengeforstærkning, jeg venter fra
Hjemmet, arriverer, vil jeg endog give Dig $3 x Rdlr.$

„Hvad mener Du med dette $x$? „$x$ er,“ svarer A, „en
foranderlig Størrelse, en varierende Eenhed, der kan antage
forskjellige Værdier, t. Ex. $1, 2, 3, 4 ...$; de $3 x$ Rdlr. ville altsaa, som Du seer, paa denne Maade kunne
blive til $... 3. 1 Rdlr., 3. 2 Rdlr., 3. 3 Rdlr.$ o. s. v.

B tænker efter: „$x$ er altid en Eenhed; de $3. 1 Rdlr.$ er
jeg da vis nok paa; maaskee kan det Belovede, naar Pengeforstærkningen
arriverer, og A er i Lune, stige til $3. 3$, ja
maaskee endog til $3. 4 Rdlr.$"; altsaa: B stoler paa den i
den varierende Eenhed liggende Forudsætning, modtager
Tilbudet, og venter. De otte Dage forløbe; Pengeforstærkningen
udebliver, og A forklarer nu B Betydningen af den
variable Størrelse: „$x$“, siger han, „kan, som sagt, være:

% --- p. 9 ---
\setcounter{footnote}{0}%
\opage{9}$1, 2, 3, 4 \ldots$, men, da det er en varierende Eenhed, kan
det jo ogsaa — det ligger i vor Forudsætning — være
$\frac{1}{2}, \frac{1}{3}, \frac{1}{4} \ldots$ Vi kunne nu denne Gang vælge at sætte $x=\frac{1}{288}$; Du faaer altsaa, ifølge vor udtrykkelige Overeenskomst, d. e.
ifølge vor Forudsætning, $3x$ Rdlr. $= \frac{1}{288}$ Rdlr. siger og % PRINTED AS IS: og og (doubled across the line break)
og skriver en heel Skilling for Din Jagttaske.“ Nu gaaer
der et Lys op for B, et Lys, der vel er sine 2 Rdlr. 95 ß
værd; thi B gjør Bekjendtskab med den varierende Eenhed,
han gjennemskuer Forudsætningen, og opdager Grunden.\footnote[1]{Man vil forstaae, at hvad der i dette Exempel, seet fra et reent mathematisk Standpunkt, vilde være for subjectivt, faaer fra Philosophiens Synspunkt en Betydning derved, at den i Forudsætningen skjulte Illusion psychologisk afsløres.}

At indlade sig med Forudsætninger uden at kjende
Grundforholdene er ofte ligesaa misligt, som at lege med
ladte Pistoler; inden man vel seer sig for, kan en Forudsætning
saa at sige gaae af, og udskyde en Conseqvens, der,
om den end ikke dræber, saa dog saarer. Functionsbegrebet
fordrer da ogsaa, at

\begin{center}γ)

\emph{Forholdet imellem Functionens Afhængige og dens Uafhængige bestemmes som et Grundforhold.}\end{center}

Uagtet Conseqvensens Forhold til Forudsætningen altid
er et Nødvendighedsforhold, bliver det dog først erkjendt som
Grundforhold, naar Forudsætningen viser sig som Grund.
Nødvendighedsforholdet udtrykker vel, at enhver Forudsætning
maa have sin tilsvarende Conseqvens, og enhver Conseqvens
sin tilsvarende Forudsætning; men, saalænge Forholdet
forbliver i ubestemt Almindelighed, er det endnu
uafgjort, hvad Udtrykket: „tilsvarende“ egentlig vil sige. For
at en tilsvarende Conseqvens skal fremgaae af den tilsvarende
Forudsætning, maa Betingelsen træde til; først ved Betingelsen
blive Forudsætningen og Conseqvensen nærmere bestemte.

% --- p. 10 ---
\setcounter{footnote}{0}%
\opage{10}Men Betingelsen, der, fra den ene Side seet, er forskjellig
fra Forudsætningen, er, fra den anden Side seet, identisk
med Forudsætningen. I sin Forskjel fra Betingelserne er
Forudsætningen det Væsentlige, og efter sit Væsen Grund;
men kun i sin Eenhed med Betingelserne er Grunden tilstrækkelig,
og som tilstrækkelig Grund (\textit{ratio sufficiens}) er
den omfattende Forudsætning: ved Forudsætningens Identitet
med Betingelsen bliver det formelle Nødvendighedsforhold et
Grundforhold.

Saalænge Ligningen: $f(x)=y$ kun udsiger, at enhver
Forandring i $y$ er afhængig af en tilsvarende Forandring i $x$, er Afhængighedsforholdet endnu ubestemt og formelt,
umiddelbart lydende paa Forudsætning og Conseqvens. Det
er imidlertid saa langt fra, at det formelle Afhængighedsforhold
skulde i og for sig have nogen sand Almeengyldighed,
at man end ikke er istand til at paavise dets Betydning
uden at indføre Betingelser;\footnote[1]{Allerede ved fra $\Phi(x, y)=0$ at gaae til $y=f(x)$ have vi sat en Betingelse.} men idet Betingelser indføres,
bliver det almene Forhold forvandlet til et særskilt. „Man
kan“, hedder det i Mathematiken, „ikke sige, at $2$ er en
Function af $3$, fordi $2=3^{2}—7$; thi $2$ er ogsaa $= 3—1$, eller $= 3^{3}—25$ o. s. v., saa at $2$ kan afhænge af $3$ paa
mangfoldige Maader. Men derimod kan man vel sige, at
$y$ er en Function af $x$ formedelst Ligningen $y=x^{2}—7$, og
at deri svarer $y=2$ til $x=3$“.\footnote[2]{Chr. Jürgensen. Begyndelsesgrund. af den mathem. Anal. S. 19.} Altsaa: Betingelsen for,
at $y$ skal vise sig afhængig af $x$, er den, at der tillægges $x$ forskjellige Værdier; Betingelsen for, at den anførte Afhængighed
skal betegne en Function, er den, at Afhængighedsmaaden
nærmere fastsættes. Men Betingelsen for den nærmere
Bestemmelse af Afhængighedsmaaden er igjen, naar vi
see nærmere til, denne, at den almene Function med indre

% --- p. 11 ---
\setcounter{footnote}{0}%
\opage{11}Nødvendighed maa sondre sig i særskilte Functioner: $f(x)=y$ faaer først Betydning ved at sættes som: $f_{1}(x)=ax$, $f_{2}(x)=x^{a}$,
$f_{3}(x)=a^{x}$, o. s. v. Kun idet Reflexionen formaaer
at udvikle den almene Function til særskilte Functioner,
og de særskilte Functioner igjen til Identitet med den
almene, kan Afhængighedsforholdet erkjendes som Grundforhold.

Stilles de særskilte Functioner: $f_{1}(x)=ax$, $f_{2}(x)=x^{a}$,
$f_{3}(x)=a^{x}$ o. s. v., imod den almene Function: $f(x)=y$, da viser det sig strax, at den almene Function hverken umiddelbart
kan holde sig udenfor, eller falde sammen med de
særskilte. Fik den almene Function en egen Gyldighed ligeoverfor
de særskilte Functioner, vilde den just derved blive
særskilt, og altsaa ikke almindelig. Faldt den almene Function
umiddelbart sammen med en af de særskilte, fik den jo
derved de øvrige udenfor sig, og var heller ikke den almindelige;
ja, selv om den antoges at falde saaledes sammen
med hver enkelt særskilt, at ingen faldt udenfor dens Omraade,
vilde den just udstykkes i Særskilthed, kunde ikke engang
være en Function for sig, endsige en almindelig Function.

At den almene Function ikke lader sig fastholde i umiddelbar Væren, betyder ingenlunde, at det Almene skulde
være uden sand Gyldighed, et tomt Schema, en subjectiv
Abstraction; det beviser tvertimod, at det Almene har i sig
selv en væsentlig og just derfor reflecteret Gyldighed: i sin
reflecterede Gyldighed er den almene Function bestemt som
Grunden til alle særskilte Functioner, som Grundfunctionen.

Den Maade, hvorpaa de særskilte Functioner, idet de
hver efter sin Eiendommelighed gjengive Grundfunctionen,
reflectere sig i hverandre indbyrdes, viser overalt, at Forudsætningerne
med deres Betingelser paa een Gang ere Reflexionsmomenter
i Grundens Væsen og dog Umiddelbarhedsbestemmelser
i den satte Mangfoldighed, den begrundede

% --- p. 12 ---
\setcounter{footnote}{0}%
\opage{12}Tilværen. Til Forstaaelse af de mange særskilte Functioner
vilde det altsaa, ifølge Grundforholdet, være tilstrækkeligt
at forstaae en eneste, naar Forstaaelsen af denne ene var
udtømmende; men til en udtømmende Forstaaelse af een
Function vilde det, ligeledes ifølge Grundforholdet, være
nødvendigt at forstaae dem alle.

\begin{center}b) Functionen og \emph{Loven}.\end{center}

Til Begrundelsen hører altsaa et Element, der gaaer
sammen med Grunden, og dog falder udenfor Grunden;
dette Element er Betingelsen. Betingelsen, som i Reflexionen
er identisk med Grunden, forsaavidt det Særskilte begrundes,
falder i sin Umiddelbarhed udenfor Grunden, og giver i
denne Henseende en Afskygning af det \emph{Tilfældige}. Er det
Tilfældige paa denne Maade et i Grunden Nødvendigt,
maa det Nødvendige jo med det Samme faae et Anstrøg af
det Tilfældige, og vise sig som betinget Nødvendighed. Det
i \emph{Tilfældene Tilfældige} er det vexlende Enkelte, det
Mangfoldige; det i Tilfældene Nødvendige er det væsentlig
Ene, det Almene, \emph{Loven}; mellem disse Yderled falder
Mellembestemmelsen af det overfladisk Almene, det Fælleds,
\emph{Reglen}. Veien fra det mangfoldigt Givne gjennem Betingelserne
til Grundeenheden er en Vei fra \emph{Tilfældet} gjennem Reglen til Loven.

\begin{center}α) \emph{Tilfældet}.\end{center}

I Erfaringens Verden forefindes Tilfældet som Kjendsgjerning,
og maa tages, som det gives; i den aprioriske
Videnskab sætter den Vidende selv Tilfældet; men i begge
Tilfælde ligner det Tilfældige sig selv deri, at, naar det
engang er sat, saa er det sat, enten det saa har viist sig
som Kjendsgjerning, eller som et \textit{a priori} Antaget. Vel
kunde det ved første Øiekast synes, som var det i Kjendsgjerning
givne Tilfælde anderledes tilfældigt, end det Bevidstheden
selv sætter. Som objectivt Tilfælde er Kjendsgjerningen

% --- p. 13 ---
\setcounter{footnote}{0}%
\opage{13}jo, fordi den er, og saaledes, som den er, uafhængig af
Subjectiviteten; det subjectivt Satte derimod er, fordi Bevidstheden
vil, at det skal være, og er just saaledes, fordi
det er sat saaledes: det Vilkaarlige synes altsaa at fortære
det Tilfældige. Dette er dog kun en Indbildning; at Vilkaarligheden
ikke kan tilintetgjøre Tilfældigheden, er let at
indsee. Vil Subjectiviteten, der sætter et Tilfælde, udelukkende
sætte det i Kraft af Vilkaarlighed, hvad er den
abstracte Vilkaarlighed da Andet, end en bevidst Tilfældighed?
Jeg sætter vilkaarlig et $a$, og er mig bevidst, at jeg
ligesaa godt kunde have sat et $b$; jeg indseer, at det Vilkaarlige
er i sig et Tilfældigt. Jo mere Subjectiviteten
giver sig hen i vilkaarlige Phantasier og Forestillinger, desto
mere giver den sig da ogsaa Tilfældigheden i Vold. Skal
det vilkaarlig Satte faae intellectuel Betydning, maa det
sættes som en Forudsætning, hvis Conseqvens anerkjendes:
saamange særegne Forudsætninger, saamange Conseqvenser;
Tilfældenes Eiendommelighed maa for Conseqvensens Skyld
respecteres. At respectere Tilfældenes Eiendommelighed er
at tillægge dem objectiv Betydning; men, idet Tilfældet
faaer objectiv Betydning, viser sig jo netop det Tilfældige.

Ikke alene i den empiriske Yderverden, ogsaa i den
aprioriske Forestillingskreds kan altsaa Tilfældenes Mængde
vise sig overvældende, og det Uforudsete overraske os. I
Mathematiken er Tallenes Antal utalligt, og Talforbindelsernes
Mulighed uudtømmelig. Idet de enkelte Talstørrelser
fastholdes, saaledes som de umiddelbart haves i Forestillingen,
er hver Talstørrelse et Phænomen, og man fatter uvilkaarlig,
at Talphænomenerne i vexlende Mangfoldighed ikke
staae tilbage for Erfaringsphænomenerne. Tallet $9$ være
t. Ex. det givne Phænomen; man begynde nu at undersøge,
paa hvilken Maade denne Størrelse fremkommer,
og lader sig opløse ved andre Størrelser; og man vil
snart overbevise sig om, at det ikke er dette eller hiint sær%
% --- p. 14 ---
\setcounter{footnote}{0}%
\opage{14}egne Tilfælde, men Millioner Gange Millioner mulige Tilfælde,
man har med at gjøre.

Hvorledes skulle vi da behandle de utallige Tilfælde?
Dette Spørgsmaal er i det Aprioriske væsentlig at besvare
paa samme Maade som i det Empiriske. Hvad de empiriske
Tilfælde angaaer, gjælder det, ved sammenlignende
Abstraction at fremhæve det Fælleds, adskille det Uligeartede,
og forbinde det Ligeartede; paa de aprioriske Tilfælde bliver,
forsaavidt det Tilfældige fremhæves, en lignende Fremgangsmaade
at anvende.

\begin{center}β) \emph{Reglen}.\end{center}

Naar det i mange Tilfælde Lige abstraheres og fastholdes,
fremkommer Reglen. I de fire Regningsarter, saaledes
som disse fremsættes og læres populært til praktisk
Brug, sees en Mængde Regler. Man begynder med Forskrifter
for Behandlingen af de særegne Opgaver; det vises
i enkelte Exempler, hvad man har at gjøre, hvorledes man
bærer sig ad med at opsætte Stykket og komme til Facit.
Ved at anvende den foreskrevne Fremgangsmaade paa en
Række Tilfælde, og see, hvorledes det hvergang slaaer til,
at det rigtige Facit udkommer, overbeviser man sig om Reglens
Gyldighed. Reglen, der er abstraheret af de enkelte
Tilfælde, har i denne Sammenhæng kun Interesse, forsaavidt den lader sig anvende paa de enkelte Tilfælde: en
Undersøgelse af Reglen i dens Almindelighed som Regel og af
Reglernes indbyrdes Forhold forudsætter et høiere Synspunkt.

Saalænge Bevidstheden er fortabt i Tilfælde og Regler,
er Tanken hildet i det Endelige; den videnskabelige Indsigt
mangler. Kjendskab til Reglerne og Færdighed i at
anvende dem, er endnu ikke Videnskab: Regnemesteren er
endnu ikke Mathematiker. Den Fordom, Mange synes at
nære mod Mathematiken, at den, uden væsentligt Forhold
til Ideen og det Ideelle, kun skulde være en tør Gjenstand
for „ „den tørre Forstand“, hidrører vel derfra, at de, der
dømme saaledes, kun kjende det Mathematiske i dets ende%
% --- p. 15 ---
\setcounter{footnote}{0}%
\opage{15}lige, meest overfladiske Form som et Indbegreb af Regnestykker
og Regler; thi deri have Poeter, Fruentimmer og
Geniale unegtelig Ret, at Regnestykker og Regler ere til for
„den tørre Forstand.“

Ved at fremhæve Reglen søger Forstanden at naae det
Almeengyldige, men Reglen som Regel er ikke almeengyldig;
enhver Regel har, eller maa idetmindste kunne have sin
Undtagelse: \textit{nulla regula sine exceptione.}\footnote[1]{Naar der altsaa efter friere Sprogbrug tales om almeengyldige Regler, maa man, forsaavidt det Almeengyldige fremhæves, tænke paa regulerende Love. I denne Forstand tages da Udtryk som: „Almindelige Differentiationsregler.“ Jürgensen. Begyndelsesgrunde af den mathem. Analyse. S. 43.} Vilde Nogen
spørge, om man da ikke af Tilfælde, som dette: $2(3+4)=2.3+2.4$ skulde kunne abstrahere den i sig almeengyldige Regel: $a(b+c)=ab+ac$? da maatte vi svare:
Nei! Det blot Abstraherede er endnu ikke erkjendt som det
Almeengyldige; det i sig Almeengyldige er ikke en Regel.
Hvor misligt det er at ville ved Abstraction fra de enkelte
Tilfælde sikkre sig det Almeengyldige, kan, hvad Mathematiken
angaaer, blandt Andet sees i Læren om Primtallene. Man
finder: $2^{3}-1=7$, $2^{5}-1=31$, $2^{7}-1=127$, lutter
Primtal; men dersom man nu af disse tre Tilfælde uddrog
den Slutning, at $2^{2m-1}-1$ nødvendigviis maatte være et
Primtal, vilde man forregne sig, da allerede det fjerde Tilfælde: $2^{9}-1=511=7.73$, ikke giver noget Primtal.
Fermat antog, at $2^{(2^{n})}+1$ var et Primtal; men Euler
viste, at $2^{(2^{5})}+1=2^{32}+1$ er deleligt ved $641$, og følgelig intet Primtal. At i slige Exempler Undtagelsen fra
Reglen indtræffer snart med det fjerde, snart med det femte
Tilfælde o. s. v., er i den hele Betragtning selv som et Tilfælde.
Var Undtagelsen først indtraadt med det 99de Tilfælde,
vi havde jo dog i den udvidede Regel endnu ikke
fundet det i sig Almeengyldige, men derimod været udsatte
for at bedrages af Tilfældenes Mængde.

% --- p. 16 ---
\setcounter{footnote}{0}%
\opage{16}
\begin{center}γ) \emph{Loven}.\end{center}

Det Almeengyldige, hvorfra der ikke er og ikke kan
gives Undtagelse, er ikke Regel, men Lov. Skal Ligningen $a(b+c)=ab+ac$ opfattes som Lov, maa dens Gyldighed
ikke beroe paa en Abstraction af de enkelte Tilfælde,
men paa Indsigt i Sagens Væsen. Et eneste Tilfælde, der
fører til Indsigt, har i saa Henseende større Værd, end
tusinde Tilfælde, der kun afgive Stof for en sammenlignende
Abstraction. At $\frac{1}{1-x}$ vil ved fortsat Division give en
Række af utallige Led efter stigende Potenser af $x$, nemlig:
\[ 1+x+x^{2}+.... \]
bevises ikke ved at fortsætte Divisionen
i det Uendelige, men ved at tydeliggjøre, hvad Divisionen
paa given Forudsætning nødvendigviis fører med sig.
I sin Umiddelbarhed staaer Tilfældet for Reflexionen som
det, hvorover der reflecteres; Resultatet af den endelige
Reflexion, det Generelle, det relativt Almene, er den af
Tilfældene abstraherede, paa Tilfældene anvendelige Regel;
i Reflexionens Eenhed med Sagens Væsen erkjendes Loven.
I Mathematiken, hvor intelligible Størrelser træde istedetfor
sandselige Phænomener, kan Loven fremstilles i den rene
exacte Form. Hvad er da Loven? Efter sit Begreb udtrykker
Loven den Nødvendighed, hvormed Følgen indtræder,
naar hele Forudsætningen d. e. Grunden med dens Betingelser
er given. Er $a=b$ og $b=c$, saa er $a=c$; den
Nødvendighed, Loven udtrykker, er altsaa hypothetisk: som
den hypothetiske Nødvendigheds Udtryk er Loven det Almeengyldige.
Enten vi sige: „det er en Lov, at Størrelserne $a$ og $c$ maae være indbyrdes lige, naar enhver af dem for sig
er lig med Størrelsen $b$"; eller vi sige: „det er nødvendigt,
at det maa være saaledes“; eller: „denne Sætning er almeengyldig“, betyder Eet og det Samme.

Som det Almeengyldige finder Loven sit Udtryk i \emph{Axiomet}, i \emph{Definitionen}, i \emph{Theoremet}. I enhver af disse
% --- p. 17 ---
\setcounter{footnote}{0}%
\opage{17}Udtryksmaader gjenkjende vi det hypothetisk Nødvendige, der
betegnes ved Forholdet mellem Forudsætning og Conseqvens:
dersom $a=b$, og $b=c$, saa er $a=c$; dersom alle
Punkter i en plan Curve have den samme Afstand fra et
bestemt Punkt i Planet, saa er Curven en Cirkel; dersom
Figuren $ABC$ er en Trekant, saa er Vinklernes Sum $=$ $2R$. Om Udtryksformen selv forøvrigt er hypothetisk eller
thetisk, gjør i denne Henseende ingen Forskjel; den væsentlige
Forskjel beroer paa Begrundelsen.

\emph{Axiomet} udtrykker Loven paa umiddelbar Maade.
„Kun een Linie kan være den korteste Vei imellem to Punkter“.
Vilde man omskrive denne Grundsætning saaledes:
„Det er en Lov, at der imellem to Punkter kun kan være
een Linie, der er den korteste“, da vilde en saadan Omskrivning
ganske vist være overflødig og tvungen. Men at
Omskrivningen udelades maa dog ikke misforstaaes, som om
Loven og Grundsætningen faldt ud fra hinanden, saa at
Grundsætningen var Et, og Loven et Andet; thi det Sande
er, at Loven ikke særlig nævnes som Lov, just fordi den
falder umiddelbart sammen med Grundsætningen: den Lovsætning,
der i sig selv er saa indlysende, at den ikke trænger
til Begrundelse, er Grundsætning.

\emph{Definitionen}, der angiver de Forudsætninger, af
hvilke Lovens Conseqvenser paa givne Betingelser fremgaae,
er i Reflexionsforhold til Loven. Definitionen falder ikke,
saaledes som Grundsætningen, umiddelbart sammen med
Loven, tvertimod, den synes endog ved første Øiekast at
forholde sig ligegyldig til Loven; men, naar Forholdet imellem
Loven og Grundsætningen nærmere skal bestemmes, viser
det sig, at Loven kun hviler i Grundsætningen, forsaavidt
den hviler paa Definitionen. „Hvorfor kan kun een Linie
være den korteste Vei imellem to Punkter?“ Det ligger i
Sagens Natur, fordi det ligger i Begrebet: „den korteste“,
men at det ligger i Begrebet, vil jo med andre Ord sige,
at det oprindelig ligger i Definitionen. Dersom $a^{x}=y$,

% --- p. 18 ---
\setcounter{footnote}{0}%
\opage{18}saa gjælder Loven: $x=\log y$, med $a$ som Grundtal. Hvorpaa
beroer saa Lovens Gyldighed? paa Logarithmens Definition.

I \emph{Theoremet} bliver det ved Definitionen angivne
Forhold paa særegen Maade formidlet, og Loven udtrykkelig
sat som Lov. „En Læresætning eller et Theorem (\textit{theorema})
fremstiller en Sandhed, som ikke er umiddelbart indlysende.
Den indeholder en Forudsætning (\textit{hypothesis}), bestaaende af
en eller flere Betingelser, og en Følgesætning (\textit{thesis}). Den
Række af Slutninger, hvorved Følgesætningen føres tilbage
til Forudsætningen, saa at den derved med Nødvendighed
begrundes, danne Læresætningens Beviis (\textit{demonstratio}).“\footnote[1]{C. Ramus. Elementær Geometrie 1850. S. 7—8.}

For at fatte Lovens Væsen og dens Forhold til Functionen
er det nødvendigt at indsee, hvorledes Beviset og
Definitionen, Læresætningen og Grundsætningen reflectere
sig i hinanden. Haves t. Ex. $Log$ som Logarithme svarende
til Grundtallet $A$, og $log$ svarende til Grundtallet $a$, da
vil den Lov, hvorefter Transformationen fra det ene af
disse Systemer til det andet foregaaer, vise sig baade at ligge
i Logarithmens Begreb, og dog at trænge til et Beviis. Af
Beviset: $A^{x}=y$, og $x=\mathrm{Log}\,y$; $a^{x}=A$, $z=\log A$, $(a^{z})^{x}=y$,
og $xz=\log y$; altsaa: $\mathrm{Log}\,y \log A=\log y$, og $\mathrm{Log}\,y=\log y \frac{1}{\log A}$, sees, at Beviisførelsen er en Udvikling af Loga-
rithmebegrebet, og at denne Begrebsudvikling er en Lovudvikling.
Idet Beviisførelsen væsentlig falder sammen med
Begrebsudviklingen, sees Grundsætningens og Læresætningens
dialektiske Eenhed. Det er Videnskabens Opgave at forvandle
enhver Grundsætning til Læresætning; det i sig Indlysende
er kun indlysende for en Reflexion: det Umiddelbare
maa formidles, det Ubegrundede maa begrundes. At den
exacte Videnskab selv anerkjender denne Fordring, sees t. Ex.
af dens Forsøg paa at føre Beviis for den Sætning, at

% --- p. 19 ---
\setcounter{footnote}{0}%
\opage{19}„den korteste Linie imellem to Punkter er en ret“\footnote[1]{See t. Ex. L. Oppermann. Elementær Plangeom. Kbh. 1845. S. 43.} At
omvendt Læresætningen maa hæves til Identitet med Grundsætningen,
vil med andre Ord sige, at det i Erkjendelsen
Middelbare maa hæves til høiere Umiddelbarhed. Beviset
har just den Betydning, at det ved at gjøre Sagens enkelte
Punkter umiddelbart indlysende viser os det hele i et Lys
af reflecteret Umiddelbarhed.

At Functionen og Loven falde sammen i en Begrebsudvikling,
er da i Qvantitetslæren uimodsigeligt; for at
bestemme, hvad en Function er, maae vi forudsætte Loven,
for at bestemme Loven, maae vi igjen forudsætte Functionen:
saamange Functioner, saamange Love. Ligesom Loven ved
at have Tilfældene i sin Magt er hævet over den blotte
Regel, saaledes er ogsaa Functionen hævet over den blotte
Regel; ligesom Loven ikke kan undvære Tilfældet, fordi den
i sin hypothetiske Nødvendighed ikke kan undvære den umiddelbare
Forudsætning, saaledes kan Functionen i sin betingede
Almeengyldighed heller ikke undvære Forudsætningens
og Betingelsens Umiddelbarhed;\footnote[2]{„Indem die Function über den Begriff des realen continuirlichen Seins im Raume hinausgeht und die einfache Unmittelbarkeit von dessen Begrenzung negirt, so geht sie doch wieder vollständig mit ihm zusammen, weil sie die allgemeine Regel ist, vermöge deren auf die gesammte Folge von Punkten dieser Curve (her henvises til den Maade, hvorpaa en plan Curve bestemmes, idet ethvert Punkt( $x, y$) den kan bestemmes, formedelst Ligningen= $f(x)$) zurückgekommen werden kann, weil sie in der unveränderlichen Folge ihrer Rechnungsoperationen das allgemeine Gesetz, das Bleibende darstellt, wodurch die aus ihrem Flusse herausgerissenen isolirten vielen Punkte wieder zu einer Einheit zusammengeschlossen werden, und weil sie mithin dieselbe befassende Einheit ist, als welche die ebene Curve in der Anschauung vorliegt.“ Versuch einer Philosophie der Mathematik verbunden mit einer Kritik der Aufstellungen Hegel's über den Zweck und die Natur der höheren Analysis von Hermann Schwarz. Halle 1853, S. 26.} ligesom den ene Grund-

% --- p. 20 ---
\setcounter{footnote}{0}%
\opage{20}function formedelst Tilfældenes Forskjellighed og dog med
indre Nødvendighed bliver adskilt i et System af særskilte
Functioner, saaledes bliver den ene Grundlov paa samme
Vilkaar sondret i et System af særskilte Love.% NOTE: no mark printed in the print; the note stands at the foot of the page and fits best here (its last words: "et System af særskilte Love")
\footnote[1]{Hvorledes Functionen og Loven falde sammen i Grundbegrebet, sees af utallige Exempler. At $F(x, y, z) = 0$ er Ligningen for en Overflade, vil jo med andre Ord sige, at denne Function udtrykker Loven for de til et givet System svarende Coordinatforandringer, hvorved Punkterne i den paagjældende Overflade bestemmes. Punkternes Antal er uendeligt, et uendeligt Antal Punkter kan ikke bestemmes enkeltviis; i Valget af de Punkter, man enkeltviis vil bestemme, er der da saaledes noget Vilkaarligt og Tilfældigt, men ifølge Functionen bliver Bestemmelsen af de utallige Punkter netop reflecteret i den Maade, hvorpaa det enkelte Punkt bestemmes: saaledes indgaaer overalt det Endelige under det Uendelige, det Foranderlige under det Bestandige, Phænomenerne under Loven, det Tilfældige under det Nødvendige, Tingen under Begrebet, det Særskilte under det Almene.\par
Hvormange særlige Væsensbestemmelser en Grundfunction og dermed en Grundlov i Begrebet kan indeholde, uden at der endnu er Tale om den Særskilthed, der individualiserer sig for en umiddelbar Anskuelse, er ligeledes indlysende af det anførte Exempel. Den almene Form $F(x, y, z) = 0$ angiver ikke blot Fladens Væsen som Flade, men af denne Form kan endog en Fleerhed af Egenskaber, nemlig, Fladens Grundegenskaber: dens Krumning i et vist Punkt, dens Beliggenhed mod andre Flader i et vist Punkt, Ligningen for dens Tangentplan til et hvilketsomhelst Punkt o. s. v. udledes. Først idet Fladens almindelige Ligning articuleres til Ligningen for særskilte Flader, fremtræder Grundlovens Indhold i et System af særskilte Love.}

\begin{center}c) Functionsbegrebets eiendommelige Fremstilling.\end{center}

Til at angive Functionsbegrebets Reqvisiter er den
mathematiske Betegnelsesmaade allerede anvendt; men det i
Betegnelsesmaaden Eiendommelige kan først blive Gjenstand
for særegen Reflexion, naar Opmærksomheden fæster
sig paa Betydningen af Mathematikens pasigraphiske Sprog,

% --- p. 21 ---
\setcounter{footnote}{0}%
\opage{21}paa den til Sproget svarende Tænkuing, og den af begges % PRINTED AS IS: Tænkuing (for Tænkning)
Forening resulterende Indsigt, den punktuelle Begriben.

\begin{center}α) \emph{Sproget.}\end{center}

Et System af Tegn, hvorved dels Følelser, dels Forestillinger
og Tanker udtrykkes, er et Sprog. Blandt de
mange Slags Sprogforskjelligheder er Forskjellen imellem
det almeenlogiske Talesprog og det pasigraphiske Tegnsprog uden
Tvivl den meest gjennemgaaende. „Til en philosophisk Fremstilling
af Begrebet bruges Ordet, og Ordet lyder forskjelligt
i de forskjellige Folkesprog; men Mathematiken nøies ikke
med Ordet; uafhængig af Folkesprogenes Forskjellighed har
den exacte Videnskab sit eget Sprog: $f(x + y) = f(x). f(y)$ er hverken dansk eller fransk, men mathematisk.“\footnote[1]{Jvfr. Phil. og Mathem. S. 9—15.}

Det mathematiske Sprog, paa hvis Uddannelse Mænd
som Vieta, Cartesius og Leibnitz have arbeidet med saameget
Held, er for Forstanden, hvad Nodesproget er for Følelsen;
kun at det som et Forstandens Sprog er i et væsentligere
Forhold til Logikeu. % PRINTED AS IS: Logikeu (for Logiken)

Mathematikens pasigraphiske Sprog er i en vis Henseende
overvurderet af Leibnitz, men undervurderet af Hegel.

Leibnitz\footnote[2]{Jvfr. Kuno Fischer. Geschichte der neuern Philosophie, zweiter Band. Mannheim 1855. S. 10—16.} der med en saa uhyre Lærdom, at „han
for sig alene forestillede et Akademie“, greb ind i alle sin
Tidsalders videnskabelige Anliggender, var Stifter af Akademierne
i Berlin, Dresden, Wien og Petersborg, fattede
endog i Rom den dristige Plan at indføre naturvidenskabelige
Studier i de italienske Klostre, og lagde overalt an paa
at oprette videnskabelige Stater, Lærdomsrepubliker efter en
overordentlig Maalestok. For et Videnskabernes Universalmonarkie
var imidlertid Folkesprogenes Forskjellighed en
følelig Hindring; den Tanke laae da nær, om det dog ikke
var muligt ved Indførelse af en direct Betegnelsesmaade,

% --- p. 22 ---
\setcounter{footnote}{0}%
\opage{22}en umiddelbar \textit{Signatura rerum}, at hæve sig over denne
sproglige Particularisme. Opfattet i Almindelighed og betragtet
i Frastand maatte en saadan Idee have noget særdeles
Tiltalende for en Genius som Leibnitz. Det var jo
netop slige Tankehieroglypher, de Gamle, navnlig Pythagoræerne,
søgte i de symbolsk-mystiske Tal, og hvorom Kabbalisterne
havde drømt saameget; saadanne Begrebscharakterer,
et saadant signativt Sprog, havde jo netop de exacteste Videnskaber,
navnlig Mathematiken i dens arithmetiske og alge-
braiske Tegn; hvorfor skulde det nu ikke være muligt at udvide
og udtrykke den Verdensviisdom, der udgjør samtlige
Videnskabers Indhold, i en Verdensskrift, et pasigraphisk
Sprog?“ „Der Mathematiker Leibnitz hat, scheint es, den
Philosophen Leibnitz zu dem Probleme der Pasigraphie überredet,
und der nach Universalität strebende Geist des letzteren
begünstigte gern ein Projekt, welches der Weltaufklärung so
viele Schwierigkeiten ersparen und die Idee auf dem kurzesten
Wege befördern wollte.“ Kuno Fischer.

Til den modsatte Yderlighed er Hegel gaaet i sin Logik.
I Læren om „Begrebet“ dadler han Euler og Lambert, der
have søgt at henføre Begrebsbestemmelser til Betegnelser ved
Linier, Figurer o. desl. „Slige Gjenstande“, siger han,
„have det Eiendommelige ved sig, at de ere udvortes for
hinanden, fixe Bestemmelser. Tages Begreberne saaledes,
høre de op at være Begreber.“ ... „Ihre Bestimmungen
sind nicht so ein Todtliegendes, wie Zahlen und Linien, deren
ihre Beziehung nicht selbst angehört; sie sind lebendige Bewegungen“.

Her er Misforstaaelse paa Misforstaaelse. Medens
Leibnitz vil have den hele Cyclus af videnskabelige Erkjendelser
udtrykt ved en med den mathematiske analog Betegnelsesmaade,
gaaer Hegel i sin Indsigelse herimod saavidt,
at han endog anseer det mathematiske Tegnsprog som saa
lidet svarende til Logikens Fordring, at kun Talesproget
alene kan være Begrebernes Sprog. „Da der Mensch die

% --- p. 23 ---
\setcounter{footnote}{0}%
\opage{23}Sprache hat, als das der Vernunft eigenthümliche Bezeichnungsmittel,
so ist es ein müssiger Einfall, sich nach einer unvollkomnern
Darstellungsweise umsehen und damit qvälen zu
wollen. Der Begriff kann als solcher wesentlich nur mit
dem Geiste aufgefaßt werden dessen Eigenthum nicht nur,
sondern dessen reines Selbst er ist. Es ist vergeblich, ihn
durch Raumfiguren und algebraische Zeichen zum Behufe
des äußerlichen Auges und einer \emph{begrifflosen} mechanischen
Behandlungsweise, eines Calculs, festhalten
zu wollen.“)\footnote[1]{Wissenschaft der Logik, zweiter Theil, Berlin 1834. Side 59.}

Det Eensidige i denne Synsmaade er iøinefaldende.
Vistnok er Talesproget et for Begrebsudviklingen uundværligt
Medium; men da „Begrebet som saadant væsentlig kun
kan opfattes med Aanden“, saa følger jo aabenbart heraf,
at det ikke er Midlet som Middel, men Midlets aandige
Brug, hvorpaa det kommer an. Det er da heller ikke
Rumsfigurer og algebraiske Tegn alene, der kunne opfattes
paa aandløs Maade; det Samme kan skee med Sprogets
Ord og Sætninger. Den, der i Sætningen kun seer en
Sum af Ord, i Ordet en Sum af Stavelser, i Stavelsen
en Sum af Bogstaver, seer jo heller ikke Begrebet. Da det
nu altsaa ikke er Midlet i og for sig, men den rette Brug
af Midlet, hvorom det gjælder, saa opstaaer det i denne
Sag afgjørende Spørgsmaal, om der til en alsidig og fuldstændig
Begrebsudvikling ikke foruden de Midler, Talesproget
tilbyder, skulde behøves endnu andre Midler. For at besvare
dette Spørgsmaal maae vi see paa:

\begin{center}β) \emph{Tanken.}\end{center}

Naar Hegel viser det Ufrugtbare i at ville, som Leibnig,
giøre de syllogistiske Figurer til Gjenstand for en combinatorisk
Calcul, og dadler en Ploucquet, der ved mathematiske
Bestemmelser af det Logiske mener at have bragt det
saavidt: \textit{posse etiam rudes mechanice totam logicam}

% --- p. 24 ---
\setcounter{footnote}{0}%
\opage{24}\textit{doceri, ut pueri mathematicam docentur,} da er han i
sin Ret: at forvandle Logiken til Forstandsmechanik er
aandløst. Ligeledes er, hvad Hegel under „Der Schluß des
Daseyns“ angaaende „den fjerde Figur: $A — A — A$ eller
den mathematiske Slutning“ bemærker, at navnlig Mellemleddet,
\textit{terminus medius,} som her skulde være det Formidlende,
ikke har nogen egen Bestemthed ligeoverfor Yderleddene,
og at det just er dette Tautologiske, som gjør, at
Slutningen falder sammen med et mathematisk Axiom,
--- vistnok aldeles berettiget. Men med alt dette har Hegel i
sin Bedømmelse af den mathematiske Udtryksmaade overseet
en Side, som, hvad Tanken og Begrebet angaaer, er den
egentlige Hovedside, nemlig \emph{det Punktuelle.}

Til en alsidig Udvikling af Qvantitetens Kategorier
hører, at det Punktuelle maa gaae i Eet med det Almene:
her spørges ikke blot om et Stort og et Større, her spørges
tillige hvor stort? hvormeget større? Talesproget, der
ikke ved Ordet kan paapege det Enkelte i dets umiddelbare
Tilværen, kan heller ikke ved Ordet tydeliggjøre den punktuelle
Bestemthed; Talesproget har altsaa, naar Tanken
fordrer det Punktuelle, en Mangel, der maa erstattes ved
det mathematiske Tegnsprog. Saadanne Fundamentalligninger
som: $f(x + y) = f(x) + f(y)$, og $f(x + y) = f(x).f(y)$
ere tillige Sætninger, der udtale et Begreb, altsaa logiske
Domme. Hvad der skal forstaaes ved $f(x + y)$, bliver i
disse Ligninger væsentlig bestemt paa samme Maade, som
naar det i Dommen bestemmes ved Prædicatet, hvad der
skal forstaaes ved Subjectet. Idet Ligningen udtrykker en
Lov, er den en Dom, der udsiger det Almene, og Tegnsproget kan i saa Henseende siges at congruere med Talesproget;
men, forsaavidt Ligningen er en Formel, har den
i sig, hvad Talesproget ikke har, en Mulighed til exacte
Bestemmelser i det Uendelige. Hvad Hastighed er, kan man
vel paa en Maade udtrykke i Ord, men Hastighedsbegrebet
fordrer, naar det med sine exacte Bestemmelser skal tydelig%

% --- p. 25 ---
\setcounter{footnote}{0}%
\opage{25}gjøres, den exacte Analyse. Det gjør en Sphæreforandring
i Erkjendelsen, naar man, efter at have forklaret, hvad
Hastighed er, i Ord, anvender den mathematiske Betegnelse, kalder Rummet $s$, Tiden $t$, Hastigheden $v$, og sætter $v=\frac{s}{t}$ som Definition paa Hastigheden; ja, for at Definitionen
skal passe ikke blot paa den jævne, men ogsaa paa den
ujævne Bevægelse, gaaer videre, sætter istedetfor Rummet $s$ Rumselementet $ds$, istedetfor Tiden $t$ Tidselementet $dt$, og
saaledes definerer Hastigheden ved $v=\frac{ds}{dt}$. Den, der kun
har Hastighedsbegrebet, forsaavidt det kan fremstilles og
tydeliggjøres i Ord, har det kun i svævende Almindelighed,
og altsaa paa en overfladisk Maade. Kategorien „Hastighed“
hører til en Kreds af Begreber, der have det Særegne ved
sig, at de ikke lade sig tydeliggjøre ved almindelig Reflexion,
men til deres væsentlige og punktuelle Udvikling fordre den
intellectuelle Handling, den mathematiske Operation: alle
Functionsbegreber ere Operationsbegreber.

\begin{center}γ) \emph{Sproget og Tanken.}\end{center}

Udvikler man, hvad Hegel i sin Retning har gjort,
Qvantitetens og Maalets Kategorier i logisk- dialektisk Form
saaledes, at Operationsmomenterne kun flygtig berøres, maa
Begrebsudviklingen blive eensidig og i flere Punkter, endog
hvad det Almindelige angaaer, uholdbar. Begreber, der
have deres videnskabelige Rødder i de exacte Analysers Jordbund,
forqvakles, naar de exacte Momenter maae savnes i
deres Dialektik. At en dialektisk Udvikling af Qvantitetsbegreberne
ikke erstatter Manglen paa Indsigt i Operationens Grundforhold, sees deraf, at, medens der vel gives
en Overgang fra den exacte Analyse til den dialektiske, gives
der ikke omvendt en Overgang fra den almeenlogiske og
dialektiske Analyse til den exacte: >>\textit{Non datur metaphysica
in mathesin via}«.\footnote[1]{C. Ramus: „Differential= og Integralregning“. Fortale.}

% --- p. 26 ---
\setcounter{footnote}{0}%
\opage{26}Opstiller man, for at slippe for videre Uleilighed med
Operationsmomenterne, den Paastand, at de til Qvantiteten
svarende Functionsbegreber skulle opfattes fra to saa forskjellige Hovedsider, at den ene udelukkende vedkommer
Mathematiken, den anden Philosophien: hvad følger saa
heraf? Enten maa man da forudsætte, at den Side, der
hører under Mathematiken, er som Begrebsside uvæsentlig,
og altsaa miskjende den exacte Videnskab; eller ogsaa indrømme,
at de omhandlede Begreber vel i Mathematiken
have en for Begribningen væsentlig, men dog for Philosophien
utilgjængelig Side. Men indrømmer man, at her i
Videnskaben skulde gives en for Philosophien utilgjængelig
Begrebsside, et \textit{cognoscendi genus}, hvis Væsen Philosophien
ikke formaaer at gjennemskue: da har man jo med
det Samme indrømmet, at Philosophien, om den end kan
raisonnere en Deel over Tænkning i Almindelighed og
Granskning i Særdeleshed, dog alligevel, naar det kommer
til Stykket, hverken kan afgive nogen almindelig d. v. s.
gjennemgaaende Erkjendelseslære eller omfattende Videnskabslære.

Den Maade, hvorpaa det Punktuelle i Functionsbegrebet
gaaer sammen til Eenhed med det Almene, viser,
at Fremstillingen ikke kan undvære den exacte Formel.
Naar Platos Timæus vil forklare, hvorledes det gik til, at
Guden, der dannede Verden, „ „udfyldte de dobbelte og tredobbelte
Mellemrum, idet han afskar Dele af det Hele, og
satte dem imellem disse, saa at der i ethvert Mellemrum
vare to Midter, af hvilke den ene overgik det første Yderled,
og blev overgaaet af det andet ved den samme Brøk af
enhver af dem“ o.s. v.; samt, hvorledes han (Guden) lod
„det Sammes Kreds bevæge sig efter Siden af et Paralle-
logram fra Venstre til Høire, men det Andets fra Høire til
Venstre efter Diogonalen“: da føler man det Besværlige,
for ikke at sige Uoverkommelige, i at skulle uden Hjælp af
Tegnsproget, beskrive med Ord, hvad Ord ikke kunne beskrive.
% --- p. 27 ---
\setcounter{footnote}{0}%
\opage{27}Indsees nu paa den anden Side Nødvendigheden af at
benytte Tegnsproget, og pointere ved Formler, da indsees
tillige, at Fremstillingen bør holdes saaledes, at den hverken
hildes i mathematiske Operationer, eller anbringer Formlerne
som løse, uforstaaelige Aphorismer. Hvad hermed
menes, maa et Exempel oplyse. I sin Logik, i det Afsnit
af Ideelæren, der omhandler Erkjendelsen\footnote[1]{Wissenschaft der Logik, zweiter Theil, Berlin 1834. S. 286.} vil Hegel godtgjøre,
at det kun er i den høiere Analysis, at Læresætningerne
blive synthetiske, og siger saaledes: „So ist die Aufløsung
der Gleichung $X^{m}-1=0$ mit Hülfe der Sinus ....

\dots{} eine synthetische Auflösung, weil die zu Hülfe genomenen
Bestimmungen, die Sinus\dots{}. nicht eine Bestimmung
der Aufgabe selbst ist.“

Men hvad vil det nu sige, at Sinus maa tages med,
naar man skal opløse Ligningen: $X^{m}-1=0$? En udførlig
Anvendelse af Eulers og Moivres Formler, en fuld-
stændig Deduction af Ligningen:
\[ X^{m}-1=(X-1)(X-\cos\frac{2\pi}{m}-\sin\frac{2\pi}{m}\sqrt{-1})(X-\cos\frac{4\pi}{m}-\sin\frac{4\pi}{m}\sqrt{-1})\ldots, \]
vilde her forvandle Tankegangen fra en philosophisk til en
mathematisk, medens fuldstændig Mangel paa Oplysning
vilde, hvad man vistnok i dette Tilfælde kan indvende mod
den hegelske Fremstilling, hindre den Læser, der ikke særlig
beskæftiger sig med det Mathematiske, fra at fatte Pointen.

I hvilket Omfang den mathematiske Oplysning maa meddeles,
maa da væsentlig afhænge af det Eftertryk, der i en
given Sammenhæng\footnote[2]{Vilde Hegel ved det anførte Exempel virkelig oplyse Læseren om, hvad det vil sige, at en mathematisk Læresætning bliver synthetisk; da havde det været ham let at simplificere Fremstillingen. Sættes i Ligningen: $X^{m}-1=0$, $m=2$, faaes $X^{2}-1=0$, d.e. $X=\pm\sqrt{1}=\pm 1$; men $\cos 0+\sin 0\sqrt{-1}=+1$, og $\cos\frac{2\pi}{2}+$ $\sin\frac{2\pi}{2}\sqrt{-1}=-1$; sættes endvidere $m=3$, faaes $X=\sqrt[3]{1}=1^{\frac{1}{3}}=+1$, $\frac{-1+\sqrt{-3}}{2}$, $\frac{-1-\sqrt{-3}}{2}$; men $\cos 0+\sin 0\sqrt{-1}=+1$, $\cos\frac{2\pi}{3}+\sin\frac{2\pi}{3}\sqrt{-1}=\frac{-1+\sqrt{3}\ \sqrt{-1}}{2}$, og $\cos\frac{4\pi}{3}+\sin\frac{4\pi}{3}\sqrt{-1}=\frac{-1-\sqrt{-3}}{2}$. Læseren behøver altsaa kun almindelige Forkundskaber, for at indsee, hvorledes det gaaer til, at Sinus og Cosinus indføres i Løsningen af anførte Ligninger. Havde den philosophiske Opgave berimod været, ved Udtrykket: $X^{m}-1=0$ exempelviis at oplyse, at Rødderne i en algebraisk Ligning kunne henføres til Formen $a\pm b\sqrt{-1}$, som ogsaa for $b=0$ indbefatter de reelle Rødder, og derved opklare, hvad det vil sige, at de imaginære Rødder forekomme parviis, saaledes at hvert Par kan betragtes som Rødder i en algebraisk Ligning af 2den Grad ($(x-a)^{2}+b^{2}=0$), og Ligningen altsaa vise sig som sammensat af lutter reelle Factorer af i det høieste 2den Grad: da havde det valgte Udtryk været at belyse fra en anden Side.} lægges paa Udviklingen: som Tanken
er, maa ogsaa Sproget være.

% --- p. 28 ---
\setcounter{footnote}{0}%
\opage{28}
\begin{center}II.\end{center}

\begin{center}Uendelighedsproblemet.\end{center}

I Functionen er Eenheden af den constante
variable Størrelse sat saaledes, at Muligheden til en Udvikling
af Forholdet mellem det Endelige og det uendelige
derved er given. Idet Størrelsen sættes constant, sættes den
med værende Grændse, og opfattes som endelig; idet Størrelsen
sættes variabel, sættes den med vordende Grændse,
og antyder en Overgang til det Uendelige. De Forandringer,
hvis Mulighed Functionen angiver, ere utallige, og
vise os Uendelighedens phænomenale Side; den Lov, Forandringerne
ifølge Functionen ere underkastede, er det i sig
Uforanderlige, og fremstiller det Uendelige fra den væsentlige
Side. „Uendelighedsproblemet er Mathematikens og Philo-
% the note 2) of p. 27 runs over onto the foot of this page: set in full at its mark

% --- p. 29 ---
\setcounter{footnote}{0}%
\opage{29}sophiens store Fælledsproblem. Men før man ret kan tænke
paa Løsningen af et Problem, maa man allerførst see at
lade Problemet blive til. Nærværende Problem er nu af
den Beskaffenhed, at det kun da kan staae klart for Bevidstheden,
naar Mathematiken vil hjælpe Philosophien med
at stille det frem“.\footnote[1]{Phil. og Mathem. S. 20.}

For at Uendelighedsproblemet kan blive tydeligt, og
Betingelserne for dets Løsning indsees., vil det fornemmelig
komme an paa at opklare følgende Punkter: det Endeliges
Forskiel fra det Uendelige, Tilnærmelsen af det Endelige til
det Uendelige, og den Maade, hvorpaa denne Tilnærmelse
lader sig articulere.

\begin{center}a) Det Endelige i dets umiddelbare Forskjel fra det Uendelige.\end{center}

For den umiddelbare Betragtning viser det Endelige sig
som det Værende, det Positive, det Uendelige som det Ikke-Værende,
det Negative; her er forsaavidt en qualitativ
Grændse imellem det Endelige og det Uendelige. Ved at
fastholde den blot qvalitative Grændse vilde det imidlertid
være umuligt at udvikle det uendeliges Begreb; berøvet
ethvert Endelighedens Gjenskin, falder det Uendelige sammen
med den tomme Negation: Problemet kan ikke blive til.
Opfattes det Endelige derimod ikke blot som et Værende,
men som en værende Størrelse, som den Størrelse, hvis
Grændser kunne overskrides, da bliver det uendelige ikke
længere en absolut Negation af Størrelsen, men en Negation
af det i Størrelsen Endelige, en Negation af den
endelige Størrelse, en Størrelsesnegation, der væsentlig
hører med til Begrebet: Størrelse. Idet den qvalitative
Grændse da saaledes viser sig som qvantiterende, fremkommer
Problemet.

Opfattes nemlig det Uendelige, ikke som absolut Størrelsesnegation,
men som Negation af den endelige Størrelse,

% --- p. 30 ---
\setcounter{footnote}{0}%
\opage{30}maa dette Uendelige for den umiddelbare Anskuen vise sig
fra aldeles modsatte Sider. Som det, der i Storhed og
Lidenhed overgaaer det Endelige, falder det udenfor Endeligheden,
som det, hvori enhver endelig Størrelse lader sig
opløse, falder det indenfor, ja sammen med Endeligheden.
Ved at gjennemskue de modsatte Synsmaader opdager Reflexionen
det Endeliges dialektiske Eenhed med det Uendelige.

\begin{center}α) \emph{Det Uendelige udenfor det Endelige.}\end{center}

Hvorledes det Uendelige, derved at det i Storhed og
Lidenhed overgaaer enhver endelig Størrelse, kommer til at
falde udenfor Endeligheden, og holde sig som et Hiinsidigt,
anskueliggjøres ved Rækker, hvis Led ere Functioner af den
Plads, de hver i sin bestemte Række indtage. Ligningerne for
den arithmetiske, den harmoniske og geometriske Række:

\[ u_{n}=f_{1}(n)=a+nx, u_{n}=f_{2}(n)=\frac{1}{a+nx}, u_{n}=f_{3}(n)=ax^{n} \]
vise, at der imellem det størst mulige endelige Led og
det uendelig Store, ligesom ogsaa imellem det mindst mulige
endelige Led og det uendelig Lille, er en uendelig Afstand.
Idet Grændsen mellem det Endelige og det uendelige er
paa een Gang baade qvalitativ og qvantitativ, danner her
sig en Overgang, som dog, saalænge de to Bestemmelser
holdes umiddelbart imod hinanden, ingen Overgang er.

I Ligningen $u_{n}=f(n)=ax^{n}$ gjør det vel i een Henseende
en Forskjel, om man imellem $ax^{0}$ og $ax^{\infty}$ udelukker
ethvert endeligt Led, eller anbringer Millioner Gange
Millioner Mellemled; men denne Forskjel er ikke destomindre
aldeles forsvindende. Forsaavidt Grændsen imellem det
Endelige og det Uendelige er qvantitativ, synes det vistnok
ved første Øiekast, som om de mange Mellemled udfyldte
Kløften, og bragte det Endelige nærmere sammen med det
Uendelige; men, naar saa det i Grændsen Qvalitative:
„enten et Endeligt eller et Uendeligt“, betones, tabe de
Millioner Led, der skulde udfylde Kløften, deres Betydning
som Overgangsled. Istedetfor at danne en Overgang fra

% --- p. 31 ---
\setcounter{footnote}{0}%
\opage{31}det Endelige til det Uendelige, betegne de mange Led kun
en Fortsættelse af det Endelige; naar $ax^{\infty}$, $ax^{-\infty}$ sættes, sættes det udtrykkelig ved en Afbrydelse,
en Negation, et Spring. Den qvalitative Grændse, der
adskiller det ene Endelige fra det andet, Hvidt fra Sort,
Dag fra Nat, Land fra Hav o. s. v. kan, saalænge her i
qvantitativ Henseende kun reflecteres paa et Endeligt, være
snart nærmere snart fjernere: Den, der kun er en Miil fra
Kysten, er naturligviis Grændsen nærmere, end Den, der er
mange Mile ude paa Havet; men, naar Grændsen er imellem
Endeligt og Uendeligt, altsaa uendelig qvantiterende, er den
som qvalitativ Grændse overalt lige nær og lige fjern:
den Negation, hvormed det Endelige afbrydes, og det
Uendelige sættes, er for hvert endeligt Led lige ubetinget,
lige negativ. At her er en Overgang, som ingen Overgang
er, viser, at den angivne Synsmaade er urolig, utilfredsstillende.
Det Utilfredsstillende sees blandt Andet deraf, at
man, istedetfor at lade det uendelige falde udenfor det Endelige,
med ligesaa megen Ret kan vælge den modsatte Synsmaade,
og søge:

\begin{center}β) \emph{Det Uendelige indenfor det Endelige.}\end{center}

Hvorledes Overgangen fra Endeligt til Uendeligt d. e.
fra Endeligt til $0$, og fra Endeligt til $\infty$ falder sammen
med en tilsvarende Overgang fra Endeligt til Endeligt,
kunne de trigonometriske Functioner anskueliggjøre. Er t. Ex. $x=90^{\circ}$, saa er $\cos x=0$, $\sin x=1$, $\mathrm{tg}\,x=\infty$ o. s. v.
Gjøres nu, hvad $\angle x$ angaaer, en Overgang fra $89^{\circ}$ til $90^{\circ}$, saa gjør $\sin x$ en Overgang fra Endeligt til Endeligt, $\cos x$ fra Endeligt til $0$, og $\mathrm{tg}\,x$ fra Endeligt til $\infty$. Ja
Overgangen fra Endeligt til Endeligt er selv en Overgang i
et Uendeligt. Imellem $\angle x=89^{\circ}$ og $\angle x=90^{\circ}$ ligge
jo uendelig mange uendelig smaae Vinkler, der alle maae
gjennemløbes, naar Overgangen skal skee: den endelige Vinkel
er, som i det Hele enhver endelig Størrelse, deelbar i det
Uendelige. Ifølge den uendelige Deelbarhed bliver det

% --- p. 32 ---
\setcounter{footnote}{0}%
\opage{32}Uendelige omsluttet af endelige Grændser; idet Delene blive
uendelig smaae, bliver deres Antal uendelig stort: de endelige
Grændser omslutte det uendelig Store, just fordi de
omslutte det uendelig Lille. Hvad er altsaa det Endelige?
Et Produkt af Uendeligheden, en Synthese af det uendelig
Lille med det uendelig Store. Og det Uendelige? Et
Endelighedens Product, en gjennemført Udfoldning og derved
en gjennemført Opløsning af det Endelige. Hvorledes
det Endelige, idet det viser sig at være uden Ende, kan
frembringe det Uendelige, sees t. Ex. af en Qvotientrække, som

\[ \frac{1}{1-1}=1+1+1+.... \]
thi selv naar Rækken
tænkes fortsat i det Uendelige, have vi dog
\[ \frac{1}{1-1} \]
til Rest
og maae begynde forfra. Det uendelige fremkommer ved
endelige Leds uophørlige Gjentagelse: idet det Endelige formedelst
Deelbarheden bliver opløst i det Uendelige, bliver
det uendelige med det Samme sat formedelst det Endelige.
Heraf følger:
\begin{center}γ) \emph{Det Endeliges og det Uendeliges dialektiske
Eenhed.}\end{center}

„Den formelle Logik holder fast paa den Grundsætning,
at det Endelige ikke er uendeligt, og det uendelige
ikke endeligt. Men saa uimodsigelige disse af Contradictionsprincipet
garanterede negative Domme end synes at
være, saa fattige, indholdsløse ere de tilsvarende positive:
„det Endelige er det Endelige, det Uendelige det Uendelige“,
thi deres saakaldte Identitet er en reen Tautologie. Som
dialektisk Videnskab har Philosophien ikke kunnet blive
staaende ved slige Tautologier. At det uendelige ikke er
endeligt, vil jo med andre Ord sige, at det er en Grændse
for det Endelige, altsaa det Ikke- Endelige. Det Endelige,
der er det Ikke-Uendelige, er da med det Samme en Grændse
for det Uendelige. Men, naar det Endelige og det Uendelige
paa denne Maade ere hinandens Grændser, maa

% --- p. 33 ---
\setcounter{footnote}{0}%
\opage{33}Grændsebestemmelsen blive dialektisk. Har nemlig det Uendelige
sin Grændse i det Endelige, saa har det jo en Endelighedsgrændse,
en endelig Grændse, og er selv endeligt.
Naaer det Endelige omvendt først sin Grændse i det Uendelige,
da har det jo ingen endelig Grændse, og er følgelig uendeligt.“\footnote[1]{Philos. og Mathem. S. 21—22.}

Forsaavidt det Uendelige omgiver det Endelige, er
Grændsen bestemt ved det Uendelige; forsaavidt det Endelige
omgiver det Uendelige, er Grændsen bestemt ved det Endelige;
at denne dobbeltsidige Bestemmelse nødvendig hører med
i den gjennemreflecterede Grændse, er let at paavise. En
Vinkel paa $90^{\circ}$ kan deles i uendelig mange uendelig smaae
Vinkler, en Vinkel paa $1^{\circ}$ ligeledes; en Qvadratmiil kan
deles i uendelig mange uendelig smaae Qvadrater, en Qvadrattomme
ligeledes: hvad følger heraf? Lægges Eftertrykket
paa det Endeliges Opløsning i det Uendelige, da synes
vistnok Følgen at være, at her ikke længere er Forskjel imellem
Vinklen paa $90^{\circ}$ og Vinklen paa $1^{\circ}$, at Delen bliver
lig det Hele, at Grændsen, med andre Ord, er udslettet.
Men lægges saa Eftertrykket igjen paa det Uendeliges Tilværen
i og ved det Endelige, da bliver der ikke blot Forskjel
imellem Endeligt og Endeligt, imellem Delen og det
Hele, men her viser sig med det Samme en bestemt Forskjel
imellem Uendeligt og Uendeligt saavel med Hensyn til det
uendelig Lille som til det uendelig Store. Det ene uendelige
Antal uendelig smaae Vinkler kan da være ligesaa forskjelligt
fra det andet, som $1^{\circ}$ fra $90^{\circ}$; det ene uendelige
Antal uendelig smaae Qvadrater ligesaa forskjelligt fra det
andet, som en Qvadrattomme fra en Qvadratmiil. At
Grændsen i det Endelige ophæves, betyder, at den sættes i
det Uendelige, at det Uendelige i sig optager Grændsen,
betyder, at det Endelige sættes. Den dialektiske Eenhed viser
sig deri, at Grændsen ligesaa meget tilhører det Uendelige
som det Endelige.

% --- p. 34 ---
\setcounter{footnote}{0}%
\opage{34}
\begin{center}b) Det Endeliges \emph{Overgang} i det Uendelige: \emph{Tilnærmelsen.}\end{center}

Dersom Grændsen imellem det Endelige og det Uendelige
blot var qvalitativ, kunde her ikke gives nogen Overgang,
men kun en Afbrydelse; var Grændsen blot qvantitativ,
vilde den ved Negationen satte Afbrydelse forflygtiges,
og Overgangen tabe sin Betydning. Kategorien: \emph{Overgang}
forudsætter en dialektisk Eenhed af det Qvalitative
og det Qvantitative, der finder sin Udvikling i \emph{Tilnærmelsen.}
$A$ være t. Ex. adskilt fra $B$ ved en qvalitativ
Grændse; $A$ er altsaa ikke $B$, $B$ ikke $A$; forsaavidt Grændsen
udelukkende er qvalitativ, er her intet Mellemværende:
hvor $A$ ophører, begynder $B$, her er ingen Overgang, ingen
Tilnærmelse. Bliver Grændsen mellem $A$ og $B$ derimod
tillige qvantitativ, da viser sig et Mellemværende; det vil
sige: man kan, dersom det Qvantiterende er en Gradforskjel,
forstærke $B$ ved at svække $A$, og omvendt; eller, dersom det
Qvantiterende er et Afstandsforhold, nærme sig til $B$ ved
at fjerne sig fra $A$, og omvendt. Idet her er en Qvalitetsbestemmelse,
er her altsaa en bestemt Modsætning; idet her
i det Qvalitative er en Qvantiteren, er her imellem de modsatte
Led et Mellemværende; idet her er et Mellemværende,
er her en Overgang, og i Overgangen en Tilnærmelse.

Da her nu ikke alene gives en qvalitativ Forskjel imellem
Endeligt og Uendeligt, men indenfor Endelighedens
Enemærker tillige en qvalitativ Forskjel imellem de endelige
Størrelser indbyrdes: kan Tilnærmelsen, som Følge heraf,
være deels en Tilnærmelse af Endeligt ved Endeligt til
Endeligt, og derved til et Opnaaeligt; deels en Tilnærmelse
enten af Endeligt ved Uendeligt til Endeligt, eller af Endeligt
ved Endeligt til Uendeligt, og derved til et Uopnaaeligt.
Ved den uendelige Tilnærmelse bliver det Uopnaaelige dog
opnaaet.

\begin{center}α) Et Opnaaeligt: Tilnærmelse af Endeligt til Endeligt.\end{center}

Naar Diogenes, idet han disputerede med Eleaterne

% --- p. 35 ---
\setcounter{footnote}{0}%
\opage{35}om Bevægelsen, gik, efter Sigende, frem og tilbage, for
saaledes i Gjerning at bevise Bevægelsens Realitet; da
var dette en overfladisk Beviismaade. Ved at gaae frem og
tilbage viste han kun i det Sandselige en Overgang fra
Endeligt til Endeligt, en Tilnærmelse til et Opnaaeligt, der
egentlig ikke var omtvistet. Her er paa en saadan Vei t. Ex.
12 Skridt fra Udgangspunktet til Maalet; naar 1 Skridt
er gjort, er der kun 11 Skridt tilbage; naar 11 Skridt ere
gjorte, er der kun 1 Skridt tilbage; 1 Skridt er en endelig
Længde, 12 Skridt er en endelig Afstand; Den, der allerede
har gjort 11 Skridt, er naturligviis Maalet 10 Skridt
nærmere, end Den, der kun har gjort 1 Skridt: Bevægelsen
i det Endelige, den endelige Tilnærmelse, er Noget, der
i den Grad forstaaes af sig selv, at her endnu ikke kan
være Tale om noget Problem. Problemet viser sig først,
naar Uendeligheden kommer med. I det Endelige kan en
Tilnærmelse unegtelig begynde, naar En blot vil gjøre det
første Skridt; men just i Forhold til Uendeligheden gjælder
det: \textit{c'est le premier pas qui coute,} thi i det første Skridt
ligger hele Problemet. For at kunne gjøre et heelt Skridt,
maa man først gjøre et halvt Skridt; for at gjøre et halvt
Skridt, maa man først gjøre det Halve af et halvt Skridt;
forud for det Halve af et halvt Skridt gaaer det Halve af
det Halve af det Halve ― i det uendelige --- af et halvt
Skridt. Det er altsaa ikke Spidsfindighed, men logisk
Conseqvens, naar Bevægelsesproblemet føres tilbage til en
uendelig lille Deel af det første Skridt, Grændsen imellem
Begyndelsens Nulpunkt og Veiens mindste endelige Stykke.

\begin{center}β) Et Uopnaaeligt: Tilnærmelse af Endeligt til Uendeligt.\end{center}

Det uendelig Store er Endeligheden for stort, thi
Endeligheden rummer kun det endelig Store; det uendelig
Lille er Endeligheden for lille: det Uendelige er et for
Endeligheden Uopnaaeligt. Kan der da gives en Tilnærmelse
til et Uopnaaeligt? At der gives en Tilnærmelse 3* til
% --- p. 36 ---
\setcounter{footnote}{0}%
\opage{36}et Opnaaeligt, forstaaer sig. „Tilnærmelse til et Opnaaeligt
indeholder ingen Modsigelse; et Uopnaaeligt uden Tilnærmelse
er heller ikke nogen Modsigelse. Der er saaledes intet
Modsigende i den Sætning, at parallele Linier, hvormeget
de end forlænges, aldrig ville skjære hinanden. Naar Mathematikeren
undertiden bruger det Udtryk om parallele Linier,
at de, forlængede i det Uendelige, ville skjære hinanden:
da kan man herved ikke egentlig tænke paa nogen Tilnærmelse;
det er en Umulighed, her fremstilles under Mulighedens
Form. Tilnærmelse til et Uopnaaeligt synes derimod
at frembyde den haardeste af alle Modsigelser. Ved Tilnærmelse
maa man jo dog i en eller anden Forstand komme
Maalet nærmere; men at komme Maalet nærmere maa dog
vel være det Samme som at komme Opnaaelsen nærmere;
men fra Opnaaelsen af det Uopnaaelige er man jo altid
lige langt borte.“\footnote[1]{Phil. og Mathem. S. 25—26}

Skulle Begreberne: „Uopnaaeligt“ og „Tilnærmelse“
bringes i logisk Forbindelse, da kunde man ifølge »\textit{diversus
respectus}« maaskee begynde med at tænke sig Sagen saaledes:
Kategorien „Tilnærmelse“ maa altid sees fra to
Synspunkter: 1) det Punkt, hvorfra man gaaer ud, Begyndelsespunktet,
og 2) det Punkt, der skal naaes, Maalet,
Endepunktet. Gives der overhovedet nogen Bevægelse fra
et bestemt udgangspunkt henimod et Maal, saa kan man,
dersom Maalet er uendelig langt borte, jo i en Forstand
gjerne indrømme, at enhver Tilnærmelse er som et Intet
imod Maalets uendelige Fjernhed, og dog, naar man seer
tilbage, betragte den fortsatte Fjernelse fra Udgangspunktet
som en Tilnærmelse til Maalet. Her er altsaa i en Henseende,
forsaavidt man nemlig seer fremad, ingen Tilnærmelse,
men dog i en anden Henseende, naar man, vel at
mærke, seer tilbage, en virkelig Tilnærmelse“. Ifølge denne
Forklaring skulde man altsaa see Tilnærmelsen som en Bort-

% --- p. 37 ---
\setcounter{footnote}{0}%
\opage{37}fjernelse, og see fremad ved at see tilbage; hvilket er en
Modsigelse.

\begin{center}γ) Den uendelige Tilnærmelse.\end{center}

Tilnærmelse til et Uopnaaeligt er en Modsigelse, naar
man nemlig forudsætter, at den skulde kunne standse ved dette
eller hiint endelige Punkt. Naar Tilnærmelsens Maal er
opnaaeligt, fuldendes den ved Maalets Opnaaelse, og faaer
altsaa en Ende; men den Tilnærmelse, der er rettet imod
det Uopnaaelige, kan aldrig ende, dens Fuldendelse er dens
uafbrudte Fortsættelse: den er og maa være \emph{uendelig}. At
Overgangen fra Endeligt til Uendeligt kan bestemmes ved
uendelig Tilnærmelse, har sin Grund i Endelighedens Væsen.
Det Endelige kan nemlig paa dobbelt Maade fortsætte
sig saaledes, at det netop ved denne sin Fortsættelse
ophæver sig selv, og gaaer over i det Uendelige.

I den til Qvotienten $\frac{1}{1-1}$ svarende Række $1+1+1+\ldots$ sees det Endelige at fortsætte sig i endelige Led
uden Ende; standses her ved et vilkaarligt Led, da kan Leddenes
Antal være saa stort man vil, det er dog alligevel
kun et endeligt Antal: Rækken har kun endelig Værdi, her
er ingen Tilnærmelse til det Uendelige. Fortsættes Rækken
derimod uden Afbrydelse, da er Fortsættelsen en uophørlig
Overskriden af den endelige Grændse, men en fortsat Overskriden
af den endelige Grændse, er en Tilnærmelse til det
Uendelige, en \emph{uendelig} Tilnærmelse.\footnote[1]{Ved denne uendelige Tilnærmelse tages endnu ikke Hensyn til Forskjellen imellem \textit{convergente} og \textit{divergente} Rækker, men ene og alene til det Endeliges egen Uendeliggjørelse.} I Rækken $\frac{1}{1-\frac{1}{2}}=1+\frac{1}{2}+\left(\frac{1}{2}\right)^{2}+\ldots$ er hele Størrelsen indesluttet mellem endelige
Grændser, da $\frac{1}{1-\frac{1}{2}}=2$; men for Leddenes For-
mindskelser er her ingen værende Grændse, saalidt som her

% --- p. 38 ---
\setcounter{footnote}{0}%
\opage{38}i det Endelige er nogen Grændse for Leddenes Antal. At
Uendelighedsgrændsen er uopnaaelig, betegnes derved, at
enhver endelig Grændse maa overskrides; men, idet Grændsen
overskrides, er her jo dog en Tilnærmelse til det Uendelige,
forsaavidt Tilnærmelsen selv er uendelig.

I geometrisk Form kan den uendelige Tilnærmelse anskueliggjøres
ved Forholdet imellem en Hyperbelgreen og dens
Asymptote. Til Hyperblens Ligning: $y=\frac{b}{a}\sqrt{x^{2}-a^{2}}$ svarer Assymptotens Ligning: $y_{1}=\frac{b}{a}x$, og Ligningen: % PRINTED AS IS: Assymptotens (Rettelser: Asymptote)
\[ y_{1}-y=\frac{ba}{x+\sqrt{x^{2}-a^{2}}} \]
til den mellem begge efter
haanden forsvindende Afstand.. I et saadant Exempel sees
det eiendommelige Forhold mellem Anskuen og Tænken, et
Forhold, som for Uendelighedsproblemet og dets Løsning
er afgjørende. Forsaavidt Curven med dens Assymptote % PRINTED AS IS: Assymptote (Rettelser: Asymptote)
fremstilles for Anskuelsen, maa den Paastand, at disse Linier
umulig under endelig Fortsættelse nogensinde skulde kunne
berøre hinanden, forekomme som et Paradox; men den anførte
Ligning for $y_{1}-y$ beviser, at dette Paradox er en
uimodsigelig Sandhed. Her er en fortsat Tilnærmelse, thi
jo større $x$ bliver, desto mindre bliver $y_{1}-y$; men da $y_{1}-y$ først kan blive $0$, naar $x$ bliver $\infty$, er Maalet
uopnaaeligt, og Tilnærmelsen uendelig.

\begin{center}c) Tilnærmelsen i punktuel Udvikling.\end{center}

Det Endelige er i sig et Ufuldendt, dets Fuldendelse
en Overgang i det Uendelige; som det Endeliges Negation,
som Abstraction, er det Uendelige ligeledes et Ufuldendt,
dets Fuldendelse en Fremgang af det Endelige; det Endeliges
og det uendeliges reflecterede Eenhed er altsaa ikke en
Væren, men en \emph{Vorden}: i \emph{Vorden} har den uendelige
Tilnærmelse sit kategoriske Udtryk. Men en almindelig
Vorden uden en Tilværen, hvori Vorden foregaaer, er
% --- p. 39 ---
\setcounter{footnote}{0}%
\opage{39}intetsigende, og en almindelig Tilværen uden tilværende Led
umulig; skal altsaa den Vorden, hvori det Endelige og
det Uendelige enes som i uendelig Tilnærmelse, tydeliggjøres,
da maa denne Vordens Begreb vise sig som et Punktlighedsbegreb.

Punktlighedsbegrebet kan nu vistnok som Begreb ikke
fordre en udtrykkelig Opregnen af alle enkelte Led og Forhold
i alle enkelte Tilfælde, thi, da en saadan Opregnen er
umulig, vilde det Almene gaae tabt, og Begrebet forsvinde;
derimod kan og maa Punktualiteten fordre, at Muligheden
til en Bestemmelse af de enkelte Led i deres eiendommelige
Forhold reflecterer sig i det almene Udtryk som i en exact
Formel. De Overgange, hvori det Endeliges og det Uendeliges
Vorden punktualiseres, blive da at opfatte saavel fra
den mathematiske som fra den logiske Side, og deres lovbestemte
Mangfoldighed at paavise, ikke blot i Functioner,
men i Systemer af Functioner.

\begin{center}α) Mathematisk-logiske Overgange.\end{center}

Begrebsudviklingens Overgange lade sig i Almindelighed
charakterisere derved, at der af en given Forudsætning uddrages
en Conseqvens, og at der i og med Conseqvensen
sættes en ny Bestemmelse. Falder Eftertrykket paa Conseqvensens
Udvikling af Forudsætningen, er Overgangen analytisk;
falder Eftertrykket paa den nye, særlig tilføiede Bestemmelse,
er Overgangen synthetisk: det Analytiske forholder
sig til den i Forskjellen fremherskende Identitet, som det
Synthetiske til den i Identiteten fremherskende Forskjel. Med
Identitet og Forskjel er al Tænkning beskæftiget; men i Logiken
anderledes end i Mathematiken. I Logiken er Tænkningen
\emph{Reflexion}, i Mathematiken \emph{Operation}; har nu det Analytiske
Overvægt over det Synthetiske, da bliver Tænkningen
i logisk Forstand \emph{objectiv} Reflexion, forsaavidt den
mathematiske Handling bliver en Transformation; har omvendt
det Synthetiske Overvægt over det Analytiske, da
% --- p. 40 ---
\setcounter{footnote}{0}%
\opage{40}bliver Tænkningen, logisk seet, \emph{teleologisk} Reflexion, forsaavidt
den mathematiske Handling bliver en Combination.

\begin{center}αα) Analysens og Synthesens exacte Anvendelse: Operation.\end{center}

Ifølge den kantiske Fornuftkritik beroer den rene Mathematik
paa „ synthetiske Domme a priori“. Sætninger som: „$5+7=12$“; „den lige Linie er den korteste Vei imellem
to Punkter ,“ o. fl., grunde sig paa en „Synthesis a priori“.
Ifølge den hegelske Logik ere de arithmetiske Læresætninger
derimod, paa enkelte Undtagelser nær, formelt analytiske.
Sandheden er, at den mathematiske Erkjendelse maa, hvad
der nødvendig følger af Operationens Væsen, være analythisk-synthetisk % PRINTED AS IS: analythisk (Rettelser: analytisk)
fra Først til Sidst.\footnote[1]{Spørgsmaalet om Forholdet imellem den analytiske og synthetiske Dom falder vistnok ikke umiddelbart sammen med Spørgsmaalet om den analytiske Methodes Forhold til den synthetiske; ikke desto mindre er det Dialektiske dog i lige Grad gjældende for begge Forhold. Til nærmere Oplysning lade vi \emph{Kant} og Hegel udtale sig udførligere. Kant (Kritik der rein. V. Gesammtausgab. Leipzig 1838. Einleitung S. 46—48): „Man sollte anfänglich zwar denken: daß der Satz $7+5=12$ ein bloß analytischer Satz sei, der aus dem Begriffe einer Summe von Sieben und Fünf nach dem Satze des Widerspruchs erfolge. Allein, wenn man es näher betrachtet, so findet man, daß der Begriff der Summe von 7 und 5 nichts weiter enthalte, als die Vereinigung beider Zahlen in eine einzige, wodurch ganz und gar nicht gedacht wird, welches diese einzige Zahl sei, die beide zusammengefaßt. Der Begriff von Zwölf ist keineswegs dadurch schon gedacht, daß ich mir jene Vereinigung von Sieben und Fünf denke, und ich mag meinen Begriff von einer solchen möglichen Summe noch so lange zergliedern, so werde ich doch darin die Zwölf nicht antreffen. Man muß über diese Begriffe hinausgehen, indem man die Anschauung zu Hülfe nimmt, die einem von beiden correspondirt, etwa seine fünf Finger, oder (wie Segner in seiner Arithmetik) fünf Puncte, und so nach und nach die Einheiten, der in der Anschauung gegebenen Fünf zu dem Begriffe der Sieben hinzuthun. (Jvfr. Phil. Propæd. S. 180—181). Hegel derimod (Wissensch. der Log. 2te Theil. Berlin 1834. S. 282—83 o. flg.): „Bekanntlich wird die Arithmetik und die allgemeineren Wissenschaften der discreten Größe vorzugsweise analytische Wissenschaft und Analysis genannt .... der Fortschritt ist die Reduktion des Ungleichen auf immer größere Gleichheit. Um ein Beispiel an den ersten Elementen zu geben, so ist die Addition das Zusammenfassen ganz zufällig ungleicher Zahlen, die Multiplikation dagegen von gleichen, worauf noch das Verhältniß der Gleichheit von der Anzahl und der Einheit folgt, und das Potenzen-Verhältniß eintritt.“ „Weil nun die Bestimmtheit des Gegenstandes und der Verhältnisse eine gesetzte ist, so ist die weitere Operation mit ihnen auch ganz analytisch, und die analytische Wissenschaft hat daher nicht sowohl Lehrsätze, als Aufgaben. Der analytische Lehrsatz enthält die Aufgabe schon für sich selbst als gelöst, und der ganz äußerliche Unterschied, der den beiden Seiten, die er gleich setzt, zukommt, ist so unwesentlich, daß ein solcher Lehrsatz als eine triviale Identität erscheinen würde. Kant hat zwar den Satz $5+7=12$ für einen synthetischen Satz erklärt .... Allein wenn das Analytische nicht das ganz abstract Identische und Tautologische $12=12$ bedeuten, und ein Fortgang in demselben überhaupt seyn soll, so muß irgend ein Unterschied vorhanden seyn, jedoch ein solcher, der sich auf keine Qualität, keine Bestimmtheit der Reflexion und noch weniger des Begriffs gründet. $5+7$ und 12 sind durchaus ganz derselbe Inhalt .... Hier ist im Geringsten kein Uebergang zu einem Andern; es ist ein bloßes Fortsetzen, d. h. Wiederholen derselben Operation, durch welche 5 und 7 entstanden ist.“ „Der Beweis eines solchen Lehrsatzes — einen solchen erforderte er, wenn er ein synthetischer Satz wäre — würde nur in der Operation des durch 7 bestimmten Fortzählens von 5 an, und in dem Erkennen der Uebereinstimmung dieses Fortgezählten mit dem bestehen, was man sonst 12 nennt, und was wieder weiter nichts, als eben jenes bestimmte Fortzählen selbst ist“ .... „Es ist daher ein höchst überflüssiges Gerüste, hier die Form der geometrischen Methode, welche sich auf synthetische Sätze bezieht, anzuwenden und der Aufgabe außer der Auflösung auch noch einen Beweis folgen zu lassen. Er kann nichts als die Tautologie ausdrücken, daß die Auflösung richtig ist, weil man operirt hat, wie aufgegeben war. Wenn die Aufgabe ist, man soll mehrere Zahlen addiren; so ist die Auflösung: man addire sie; der Beweis zeigt, daß die Auflösung richtig ist, darum weil aufgegeben, war zu addiren, und man addirt hat.“ — Hvad den synthetiske Erkjendelse angaaer, lægger Hegel (Anfrt. Skrift S. 288—313) særdeles Eftertryk paa Forholdet imellem Definitionen, Inddelingen og Læresætningen. „Die Definition enthält nur Eine Bestimmtheit, die Eintheilung die Bestimmtheit gegen andere; in der Vereinzelung ist der Gegenstand in sich selbst aus einander gegangen. Insofern die Defintion beim allgemeinen Begriffe stehen bleibt, so ist dagegen in den Lehrsätzen der Gegenstand in seiner Realität, in den Bedingungen und Formen seines reellen Daseyns erkannt“ ... „Der Lehrsatz nun nach der angegebenen Bestimmung ist das eigentlich Synthetische eines Gegenstandes, insofern die Verhältnisse seines Bestimmtheiten nothwendig, das ist, in der innern Indentität des Begriffes gegründet sind. Das Synthetische in der Definition und Eintheilung ist eine äußerlich aufgenommene Verknüpfung; das Vorgefundene wird in die Form des Begriffes gebracht, aber als vorgefunden wird der ganze Inhalt nur monstrirt; der Lehrsatz aber soll demonstrirt werden.“ .... „Das glänzende Beispiel der synthetischen Methode ist die geometrische Wissenschaft (S. 313). .... „Wenn man die Lehrsätze“ — hedder det S. 307 — „einer synthetischen Wissenschaft, und namentlich der Geometrie, näher vergleicht, so wird sich dieser Unterschied zeigen, daß einige ihrer Lehrsätze nur einzelne Verhältnisse des Gegenstandes enthalten, andere aber solche Verhältnisse, in welchen die vollständige Bestimmtheit des Gegenstandes ausgedrückt ist. Es ist eine sehr oberflächliche Ansicht, wenn die sämmtlichen Sätze an Werth einander gleichgeachtet werden, weil überhaupt jeder eine Wahrheit enthalte, und im formellen Gange, im Zusammenhange des Beweisens gleich wesentlich sey. At denne paa Abstillelse af Definition, Inddeling og Læresætning hvilende Opfattelse af den synthetiske Methode imidlertid ligesaavel passer paa Arithmetiken som paa Geometrien, er da blandt Andet indlysende deraf, at Arithmetiken ligesaavel behøver Definitioner, Inddelinger og Læresætninger, som Geometrien. I Elementær Algebra ved C. Ramus, Kbhvn 1855, navnlig i „Andet Afsnit. Fundamentalsætninger om hele Tal“, samt i „Tredie Afsnit. Irrationalitet og Approximation; Decimalbrøker“, forekomme arithmetiske Theoremer, og Methoden er altsaa forsaavidt synthetisk. Er nu den synthetiske Methode anvendelig i Arithmetiken, da er, hvad den nyere Videnskab fornemmelig skylder Cartesius, den analytiske Methode omvendt jo ogsaa anvendelig paa Geometrien. Hvad her paa den ene Side opstilles som Læresætning, og behandles synthetisk, kan da paa den anden Side fremstilles som Opgave, og behandles analytisk. „Den Linie, som halverer Linien AB og staaer lodret paa samme, er det geometriske Sted for de Punkter, som have ligestor Afstand fra de to faste Punkter A og B.“ At man, idet man af de retvinklene Trianglers Congruens beviser denne Læresætning, skulde komme til en indholdsrigere Erkjendelse, end ved at forvandle Læresætningen til Opgave, og løse Opgaven analytisk, er ikke til at indsee. Tvertimod! idet Opgaven behandles analytisk, idet man vælger et retvinklet Coordinatsystem t. Ex. det, der har AB til X-Axe og A til Begyndelsespunkt, og lader Punktet (x, y) være et af de søgte Punkter, har man, for AB=a, $x^{2}+y^{2}=(a-x)^{2}+y^{2}$, og i Ligningen $x=\pm(a-x)$ en dobbelt Løsning, nemlig $x=\frac{a}{2}$, der betegner en ret Linie gjennem Midtpunktet af AB og lodret derpaa, og $x=\infty$, d. e. en ret Linie ligeledes lodret paa AB, men i en uendelig Afstand fra Punkterne A og B. Medens nu den synthetiske Læresætning kun angav een Linie som det geometriske Sted, har Analysen derimod viist, at her kunne være to Linier; den analytiske Methode har altsaa i dette Tilfælde givet en rigere Erkjendelse end den synthetiske. Det vil da uden Tvivl, jo omhyggeligere man undersøger Sagen, blive stedse tydeligere, at, forsaavidt man ikke lader Abstillelsen af analytisk og synthetisk Methode gjælde efter Tradition og Vedtægt, men seer paa Erkjendelsens Væsen, kunne de to Methoder ikke holde sig afsluttede ligeoverfor hinanden, men det Analytiske maa, selv hvor det fremhersker, overalt forudsætte det Synthetiske, og omvendt, det Synthetiske det Analytiske.} Ligningen: $a+a=2a$ % PRINTED AS IS: retvinklene (Rettelser: retvinklede)

% --- p. 41 ---
\setcounter{footnote}{0}%
\opage{41}er i Grunden en analytisk Dom, der siger, at Multiplicationen
kun er en Addition; men da Multiplicationen jo
% the note 1) of p. 40 runs over onto the foot of this page: set in full at its mark

% --- p. 42 ---
\setcounter{footnote}{0}%
\opage{42}dog alligevel er en særegen Handling, en eiendommelig Function,
er Dommen tillige synthetisk. Det Samme gjælder om
% the note 1) of p. 40 runs over onto the foot of this page: set in full at its mark
% --- p. 43 ---
\setcounter{footnote}{0}%
\opage{43}$a\cdot a=a^{2}$ o. s. fr. Hvorledes det Synthetiske under en
fremadskridende Operation kan vore med det Analytiske, kan
udtrykkelig paavises. I en Ligning, som: $a^{x}=A_{0}+A_{1}x+A_{2}x^{2}+\ldots$ kan Relationen imellem Coefficienterne først
% the note 1) of p. 40 runs over onto the foot of this page: set in full at its mark

% --- p. 44 ---
\setcounter{footnote}{0}%
\opage{44}udfindes, naar man hermed sammenstiller en \emph{anden} Ligning
af samme Form. $a^{z}=A_{0}+A_{1}z+A_{2}z^{2}+\ldots$ Multipliceres nu de to Ligninger med hinanden saaledes,
at Productet erholder en dobbelt Form, nemlig: $a^{x}a^{z}=A_{0}+A_{1}^{2}xz+A_{2}^{2}x^{2}z^{2}+\ldots.$ og $a^{x+z}=A_{0}+A_{1}(x+z)+A_{2}(x+z)^{2}+\ldots.$, da sees allerede heraf, hvorledes det Synthetiske under Operationen er blevet forstærket
ved den fortsatte Analyse. Synthesens Medvirken bliver
især kjendelig, naar man, efter at have udviklet det ene Led $A_{2}(x+z)^{2}$, sammenligner $2A_{2}xz$ med $A_{1}^{2}xz$ den første
Række, thi $2A_{2}$ og $A_{1}^{2}$ stille sig da som Coefficienter, der
tilhøre hver sin Række, med en saa umiddelbar Selvstændighed
ligeoverfor hinanden, at man, for at bevise deres
Identitet, maa anvende de ubestemte Coefficienters Methode.

Hvad det Endeliges Overgang i det Uendelige selv angaaer,
beroer det i Tilnærmelsen Punktuelle væsentlig paa
Forholdet imellem Analysen og Synthesen. I Qvotienten $\frac{1}{1-x}$ ere de endelige Led: $1$ og $x$ saaledes sammenstillede
(Synthese), at en uendelig Række heraf lader sig udvikle
(Analyse). Udviklingens intellectuelle Handling, Operationen,
er „en Synthesis \textit{a priori}“; men en Synthesis \textit{a priori} er
en analyserende Synthesis: det Aprioriske er analytisk.

\begin{center}ββ) Analyse paa synthetisk Grundlag: Transformation.\end{center}

Er det nu det samme Indhold, der fremstilles i forskjellig
Form, har Analysen Overvægt over Synthesen. I
Ligningen: $\frac{a+b}{c+d}=\frac{1}{c+d}(a+b)$ er det ene Udtryk kun
en Omskrivning af det andet; selv i Ligningen $\frac{a+b}{c+d}=\frac{1}{c}(a+b)\left(1-\frac{d}{c}\ldots\right)$ hersker det Analytiske. Forsaavidt Analysen
anvendes til Løsning af en bestemt Opgave, er en Ombytten
af de forskjellige Udtryksmaader dog alligevel mere end en
% --- p. 45 ---
\setcounter{footnote}{0}%
\opage{45}vilkaarlig Omskrivning; den er en Omformen, der betinger
Indholdets Udvikling, en væsentlig Omdannelse, en Transformation.
Qvotienten $\frac{a+b}{c+d}$ udtrykker en bestemt Relation
imellem Størrelserne $a, b, c$ og $d$; Produktet $\frac{1}{c}(a+b)\left(1-\frac{d}{c}\ldots\right)$, hvori en uendelig Tilnærmelse er antydet, udtrykker
Relationen imellem de samme Størrelser paa en anden
Maade, og bidrager derved til at give Relationen en
fuldstændig Gjennemsigtighed; den til fuldstændig Gjennemsigtighed
klarede Relation er \emph{objectiv} Reflexion.

Hvorledes Transformationen tjener den objective Reflexion,
og betinger Indsigten, kan en nærmere Behandling
af Rækken: $a^{x}=A_{0}+A_{1}x+A_{2}x^{2}+\ldots$ vise. At Rækken
gaaer frem efter stigende Potenser af $x$, er saa indlysende,
at man i denne Henseende strax kan angive
det $n$te Led ved $A_{n-1}x^{n-1}$, og i dette ene Udtryk finde
Rækkens utallige Potenser af $x$ reflecterede. En saadan
Reflexion finder derimod endnu ikke Sted, hvad Coefficienterne $A_{0}, A_{1}, A_{2}, \ldots$ angaaer. At her maa være en Relation
imellem disse Coefficienter og $a$ i Udtrykket: $a^{x}$, tør
forudsættes; men for at denne Relation skal kunne paavises
og klares i Reflexion, maa Transformationen anvendes.
Ved at sætte $x=0$, og finde $a^{0}=A_{0}=1$, kan man da
for det Første transformere Rækken til $a^{x}=1+A_{1}x+A_{2}x^{2}+\ldots$; ved dernæst, i Conseqvens af vort Foregaaende
$A_{2}=\frac{A_{1}^{2}}{1.2}$ at sætte $A_{3}=\frac{A_{1}^{3}}{1.2.3}$, $A_{m}=\frac{A_{1}^{m}}{1.2.3\ldots m}$, kan man
transformere Rækken til $a^{x}=1+A_{1}\frac{x}{1}+A_{1}^{2}\frac{x^{2}}{1.2}+A_{1}^{3}\frac{x^{3}}{1.2.3}+\ldots$, og paa denne Maade see Rækkens utallige Coefficienter
reflecterede i $A$; sætter man nu endvidere $x=\frac{1}{a}$, og hvorved % PRINTED AS IS: x=\frac{1}{a} (for x=\frac{1}{A_{1}})
Rækken forandres til: $a^{\frac{1}{A_{1}}}=1+\frac{A_{1}}{1}\left(\frac{1}{A_{1}}\right)+\frac{A_{1}^{2}}{1.2}\left(\frac{1}{A_{1}}\right)^{2}+\ldots$

% --- p. 46 ---
\setcounter{footnote}{0}%
\opage{46}$=1+\frac{1}{1}+\frac{1}{1.2}+\frac{1}{1.2.3}+\ldots=e$, samt $A_{1}=\frac{\log a}{\log e}$, eller, i det naturlige System, $A_{1}=l.a$, saa transformeres
Rækken endelig til $a^{x}=1+l.a\frac{x}{1}+(l.a)^{2}\frac{x^{2}}{1.2}+(l.a)^{3}\frac{x^{3}}{1.2.3}+\ldots$; og man har da en Form, der hæver Relationen
til alsidig Reflexion.\footnote[1]{At Reflexionen nu er gjennemført, sees da ogsaa deraf, at alle særskilte Led med deres eiendommelige Bestemthed afspeile sig i det $n$ eller almindelige Led: $l(a)^{n-1}\frac{x^{n-1}}{1.2\ldots(n-1)}$.}

Idet Forholdet mellem det Endelige og det uendelige
punktuelt bestemmes ved Tilnærmelse, svæver det dialektisk
imellem Forestilling og Reflexion. De særskilte Led i den
uendelige Række ere til for Forestillingen, og i Forestillingen
hersker det Endelige; men naar den alsidige Relation er
paaviist, og Rækkens utallige Led sees at reflectere sig saavel i
$a^{x}$, som i hverandre indbyrdes, kommer Uendeligheden
tilsyne: det Uendelige er ikke til for Forestillingen, det er
kun i Reflexionen for den rene Indseen.

\begin{center}γγ) Synthese paa analytisk Grundlag: Combination.\end{center}

Aabner Analysen nu for Indsigten en saadan Mængde
Muligheder, at Synthesen gaaer i Eet med Valget af de
Tilfælde, der netop behøves til Opgavens Løsning, bliver
Reflexionen \emph{teleologisk}. Har man saaledes ved Analysen
forvisset sig om, at Coefficienterne i Rækken: $a^{x}=A_{0}+A_{1}x+A_{2}x^{2}+\ldots.$ alle ere indeholdte i Formlen: $A_{m}=\frac{A_{1}^{m}}{1.2.3\ldots m}$, da er en Bestemmen af den Relation, hvori $A_{1}$ staaer til $a^{x}$, nærmeste Opgave. Denne Opgaves Løsning
er altsaa et Formaal; men til Formaalets Opnaaelse
gjælder det her, iblandt Millioner mulige Tilfælde at vælge
just det Tilfælde, der afgiver Midlet; den omfattende Opgave

% --- p. 47 ---
\setcounter{footnote}{0}%
\opage{47}medfører en Mangfoldighed af mindre Opgaver: den Udvei,
at sætte $x=\frac{1}{A_{1}}$ er valgt blandt utallige Muligheder, og
forudsætter en oversvævende Bevidsthed.

Idet den teleologiske Reflexion saaledes bestemmer den exacte
Handling, og opgaaer i den, er Operationen combinerende.
Ere de ueensartede Rækker: $e^{x}=1+\frac{x}{1}+\frac{x^{2}}{1.2}+\ldots.$ $\sin x=x-\frac{x^{3}}{1.2.3}+\frac{x^{5}}{1.2.3.4.5}-\ldots$,
$\cos x=1-\frac{x^{2}}{1.2}+\frac{x^{4}}{1.2.3.4}-\ldots$, hver især udviklede og bestemte ved en
Synthesen overveiende Analyse; da kunne disse Rækker,
hvad Euler har viist, igjen combineres ved en Analysen overveiende
Synthese. I Combinationen hersker det Synthetiske;
men de Ligheder, Muligheder og Mellemled, hvorpaa Overgangene
beroe, maae prøves og retfærdiggjøres ved Analysen.
Det Mellemled, Combinationen her fordrer,\footnote[1]{Medens de for Udviklingen særegne Udgangspunkter angives og fastsættes ved combinerende Synthese, beroer Udviklingen selv derimod paa Analysen. Under Combinationen kunne de levende Synspunkter fastsættes saaledes, at Synthesen ikke blot faaer eet, men flere endog meget forskjellige Mellemled. Udvikles til Ex. $e^{\pm x\sqrt{-1}}$ i Række, faaes ved en Abstillelse af den reelle og den imaginære Deel: $\left\{1-\frac{x^{2}}{1.2}+\frac{x^{4}}{1.2.3.4}-\ldots\right\}\pm\left\{x-\frac{x^{3}}{1.2.3}+\frac{x^{5}}{1.2.3.4.5}-\ldots\right\}\sqrt{-1}$, som kan betegnes ved $C(x)\pm S(x)\sqrt{-1}$. Viser det sig nu ved hensigtsvarende Behandling af dette Udtryk, at nemlig: 1) $C^{2}(x)+S^{2}(x)=1$, 2) $S(x\pm y)=S(x)C(y)\pm C(x)S(y)$ og $C(x\pm y)=C(x)C(y)\mp S(x)S(y)$, 3) $C(4pq+x)=C(x)$ og $S(4pq+x)=S(x)$: da slutte vi, at dersom $q$„ findes at have Værdien $\frac{\pi}{2}$, for hvilken $\cos=0$, $\sin=1$, saa kunne $C(x)$ og cos x, $S(x)$ og cos x, S(x) og sin (x) gjennemløbe de samme Værdier for de samme $x$, og altsaa være identiske.“ $\pi$ er altsaa her, idet vi slutte saaledes, et Synthesens Mellemled. Jvfr. A. Steen. Reen Mathematik. Kbhvn 1856. S. 318—26.} nemlig $\sqrt{-1}$

% --- p. 48 ---
\setcounter{footnote}{0}%
\opage{48}betegner en Periodicitet, thi ved fortsat Potenseren, giver $\sqrt{-1}:\ \sqrt{-1},-1,-\sqrt{-1},+1,\sqrt{-1}$ o. s. fr.

Sættes altsaa $e^{x\sqrt{-1}}$ istedetfor $e^{x}$, erholdes følgende Resultater:
$e^{x\sqrt{-1}}=\cos x+\sin x\sqrt{-1}$, hvorved den exponentielle
Function er combineret med de trigonometriske; dernæst $x\sqrt{-1}=l(\cos x+\sin x\sqrt{-1})$, hvorved Logarithmens
Begreb lader sig udvikle derhen, at $l(+a)=l^{1}a+2p\pi\sqrt{-1}$ viser sig reel for $p=0$, hvorimod $l(-a)=l^{1}a+(2p+1)\pi\sqrt{-1}$ aldrig kan blive reel for nogensomhelst
heel Værdi af $p$. Ifølge Combinationen blive altsaa
Overgangene fra det Endelige i det uendelige ikke blot
forskjelligartede, men ogsaa planmæssige, og de planmæssige
Overgange utallige.

\begin{center}β) Functionssystemer: den indre Uendelighed.\end{center}

Den mathematiske Function, der, ifølge sit Begreb,
udtrykker det Foranderliges dialektiske Eenhed med det Uforanderlige,
udtrykker med det Samme det Endeliges Overgang
i det Uendelige: saamange Functioner, saamange
Overgange, saa mange Former for den punktuelle Vorden,
saamange Tilnærmelsens Veie. Disse Veie løbe ikke blot
ved Siden af hinanden; de krydse hinanden, de slynge sig
i hinanden, de danne et Net, hvori det Endelige, saa at
sige, fanger det Uendelige. For Operationen er nu dette
Net et System af Functioner, ja, af Functionssystemer;
for Reflexionen er det Udfoldningen af det Uendeliges Indhold,
den indre Uendelighed. De logiske Grundbevægelser i
Eenhed med de exacte Handlinger ere da: en Henføren af
de mange Functioner under en Grundformel, en Udvikling
af de særskilte Functioner, forsaavidt disse hver for sig lade
sig aflede af den fælleds Grundformel, og endelig en Paa%
% the note 1) of p. 47 runs over onto the foot of this page: set in full at its mark

% --- p. 49 ---
\setcounter{footnote}{0}%
\opage{49}viisning af den Betydning, de punktuelle Tilvæxter og Rester
i hver særskilt Række erholde.

\begin{center}αα) Functionssystemernes Grundformel: Reduction.\end{center}

Ligesom de mange endelige Tilfælde kunne henføres
under en særskilt Function, saaledes kunne de mange særskilte
Functioner henføres under en Grundformel. Ved denne
Reduction bliver det Enkelte og Særskilte jo vistnok opløst
i det Almene, og gaaer forsaavidt, logisk talt, til Grunde.
Men at gaae til Grunde betyder, hvad Schelling og Hegel
i metaphysisk Retning have paaviist, et Dobbelt: at tilintetgjøres,
det vil sige: \emph{opløses}, og at ophæves d. e.
\emph{opbevares} i sin Grund. For at reducere, maa man
omforme: enhver Reduction er en Transformation, men
ikke omvendt. Hvad den exacte Udvikling angaaer, kunne
alle mulige Functioner af $x+h$ henføres til den saakaldte
taylorske Række: $f(x+h)=f(x)+f'(x)\frac{h}{1}+f''(x)\frac{h^{2}}{1.2}+\ldots$ Denne Række, der har det Eiendommelige ved sig,
at Coefficienterne til de stigende hele og positive Potenser af $h$ ere Functioner af $x$, viser hvorledes en Function af $x$ kan udfolde sig i et System af Functioner, naar $x$ faaer Tilvæxten $h$. Forsaavidt $f(x)$ endnu ikke er nærmere bestemt,
er det hele af $f(x+h)$ fremgaaende Functionssystem
ligeledes ubestemt; men denne Ubestemthed er, hvad der ligger
i Grundens Væsen, ikke desto mindre en negativ Bestemthed,
en Reflexion af opløste og i Opløsningen bevarede Functioner.
Den opstillede Formel er ikke et umiddelbart Første;
den er ifølge Reductionen et til væsentlig Forudsætning forvandlet
Resultat, og som saadan Grundformel.

Hvilken Forskjel det gjør, om man i Erkjendelsen kommer
til det Almene ved Abstraction alene, eller ved Reduction,
kan sees i nærværende Sammenhæng. Ifølge Abstractionen
er det Almene tomt, den opstillede Række et dødt
Schema, der falder udenfor det Indhold, hvormed det vil%
% --- p. 50 ---
\setcounter{footnote}{0}%
\opage{50}kaarlig kan udfyldes; ifølge Reductionen er det Almene det
i Grunden Indholdsrige, der stadfæster sin Almeengyldighed
ved Begrundelsen af det Særskilte. At de særskilte Functionsrækker
reduceres til een Grundformel, vil med andre
Ord sige, at de særskilte Love have deres Bestaaen i een
Grundlov.\footnote[1]{Fra Bestemtheden af denne for Anskuelsen givne Kugleflade: $x^{2}+y^{2}+z^{2}=3^{2}$ gaaer man ved Reduction til en Bestemmelse af Fladens Art som Kugle: $x^{2}+y^{2}+z^{2}=r^{2}$, fra denne ved en ny Reduction til Kuglen som Omdreiningsflade: $z=f(x^{2}+y^{2})$ og fra denne særegne Art til Fladen som mathematisk Genus: $F(x,y,z)=0$. Hvad Deductionen, den omvendte Fremgangsmaade, angaaer, maa bemærkes, at de mathematiske Individer lade sig deducere, de impiriske derimod ikke.} % PRINTED AS IS: impiriske (Rettelser: empiriske)

\begin{center}ββ) Særskilte Functionssystemer: Deduction.\end{center}

Kun hvad der indeholdes i Grunden kan under givne
Betingelser udledes af den: Deductionen forudsætter Reductionen.
Betingelsen for et særeget Systems Udtræden af
Grundformlen er en udtrykkelig Fastsætten af den oprindelige
Function. Kjender man nemlig den Lov, hvorefter $f'(x)$, udviklet i Række ifølge Grundformlen, fremkommer af Rækkeudviklingen
for $f(x)$, saa kjender man med det Samme den
Lov, hvorefter $f''(x)$ afledes af $f'(x)$, $f'''(x)$ af $f''(x)$ o. s. fr.

Sættes t. Ex. $f(x)=ax$, bliver Rækken for $a(x+h)$ til $ax+ah$; sættes $f(x)=x^{a}$ og $f(x+h)=(x+h)^{a}$, faaes
Systemet: $x^{a}+ax^{a-1}\frac{h}{1}+a(a-1)x^{a-2}\frac{h^{2}}{1.2}+\ldots.$; til $f(x)=a^{x}$ og $f(x+h)=a^{x+h}$ svarer Systemet: $a^{x}\left(1+la\frac{h}{1}+(l.a)^{2}\frac{h^{2}}{1.2}+\ldots\right)$; til $f(x)=l.x$, og $f(x+h)=l.(x+h)$ Systemet: $lx+\frac{1}{1}\frac{h}{x}\div\frac{1}{2}\frac{h^{2}}{x^{2}}+\ldots$ o. s. fr. I Forhold til Reductionen er Muligheden Grund, i Forhold
til Deductionen er Grunden Mulighed. Identiteten af Grund
og Mulighed sees af den disjunctive Doms forskjellige For-

% --- p. 51 ---
\setcounter{footnote}{0}%
\opage{51}mer: $A$ er hverken $B$ eller $C$ eller $D\ldots,$ $f(x)$ er hverken $ax$ eller $x^{a}$ eller $a^{x}\ldots;$ de særskilte Bestemmelser ere opløste
i det Almene, og derved gaaede til Grunde; $A$ er, i Grunden,
baade $B$ og $C$ og $D\ldots,$ $f(x)$ er, i Mulighed, baade $ax, x^{a}, a^{x}\ldots$ Forsaavidt det i Grundens Indhold Særstilte
lader sig tilsyne, reflecteres det i Mulighed; hvor eet
Tilfælde er muligt, maae idetmindste to Tilfælde være mulige,
thi dersom det mulige Tilfældes Modsatte ikke ogsaa
var muligt, var Tilfældet selv nødvendigt: Muligheden har
alternerende Led. Skal nu et af Leddene sættes som tilværende,
da viser den disjunctive Dom i dens sædvanlige Form: $A$ er enten $B$ eller $C$ eller $D\ldots,$ $f(x)$ er enten $ax$ eller $x^{a}$ eller $a^{x}\ldots$ Spændingen imellem de alternerende Led. Hvorledes
det Tilværende ved en udtrykkelig Sætten afledes af
Grunden, det Virkelige af Muligheden, det Særskilte af
den concrete Almindelighed, kan sees af den disjunctive
Slutnings dobbelte \textit{modus}: $f(x)$ er enten $ax$ eller $x^{a}$ eller $a^{x}\ldots,$ nu er $f(x(=ax,$ altsaa hverken $x^{a}$ eller
$a^{x}\ldots,$ (\textit{modus ponendo tollens}), og: $f(x)$ er enten $ax$ eller
$x^{a}$ eller $a^{x};$ nu er $f(x)$ hverken $x^{a}$ eller $a^{x},$ altsaa maa
den, forudsat, at her kun ere tre Muligheder, eller rettere
at den paagjældende Mulighed kun har tre alternerende
Led, nødvendigviis være $ax$ (\textit{modus tollendo
ponens}). At Functionerne: $ax, x^{a}, a^{x}\ldots\ldots$ udelukke
hinanden (\textit{principium exclusi medii}) har, ifølge Deductionen
sin Betydning deri, at de i Grunden (\textit{principium
inclusi medii}) forudsætte hinanden. Det Samme, der gjør
at $f(x+h)=(x+h)^{a}$ giver Rækken: $x^{a}+ax^{a-1}\frac{h}{1}+\ldots,$ gjør ogsaa, at $a(x+h)$ giver $ax+ah,$ at $a^{x+h}$ giver $a^{x}+a^{x}\,l.\,a\,\frac{h}{1}+\ldots$ o. s. v., thi det er just Grundloven
selv, Lovens Grundtanke.

De særskilte Systemer ere Udtryk for særskilte Love.
Viser nu Reductionen, at de særskilte Love maae have deres

% --- p. 52 ---
\setcounter{footnote}{0}%
\opage{52}væsentlige Forudsætning i den ene Grundlov, saa viser
Deductionen omvendt, at den ene Grundlov har sin Bestaaen,
sin endelige Tilværen i de særskilte Love.

\begin{center}γγ)

Tilværter og Rester: \emph{Articulation.}\end{center}

Som det, hvori Endeligheden opløses, er det Uendelige,
væsentlig seet, det Endeliges Grund og Mulighed; en
gjennemført Udvikling af Endelighedens Former er da med
det Samme en Udvikling af Uendelighedens Indhold. Fordringen
er, at Udviklingen bliver punktuel. Kaster man
Punktualiteten bort, og holder sig eensidig ved det Almindelige,
bliver man let færdig med Forskjellen imellem det i
Forestillingen og i Begrebet Uendelige, eller, for at bruge
Spinozas Udtryk, imellem \textit{infinitum imaginationis} og \textit{infinitum
actu;} men anerkjendes Punktualitetens Fordringer,
indlyser det, at Forestillingens Led høre med til Begrebet selv.

Hvad Qvotienten $\frac{1}{1-1}$ angaaer, er Rækken $1+1+1+\ldots$ nødvendig til en Bestemmelse af det Uendelige; hvad
Rækken angaaer, er Qvotienten nødvendig. Holde vi os
til Rækken $1+1+1+\ldots$ alene, da see vi kun den ene
Side af det Uendelige, nemlig det Uopnaaelige, det, hvis
endelige Fremstilling aldrig kan fuldendes, \textit{infinitum imaginationis}.
Holde vi os udelukkende til Betegnelsen $\frac{1}{1-1}$=

$=\frac{1}{0}=\infty,$ uden at betænke, hvorledes 00 først fremkommer
ifølge en uendelig Formerelse af en endelig Qvotient,
idet $0$ fremkommer ifølge en uendelig Formindskelse af en
endelig Divisor, faae vi vel formelt et \textit{infinitum actu,} et
Uendelighedsbegreb, men et dødt Begreb. Til Indsigt i
Uendelighedens Væsen, den indre Uendelighed, hører altsaa,
at infinitum imaginationis og \textit{infinitum actu} maae sættes i
dialektisk Eenhed, at en Række som: $1+1+1+\ldots$ og
en Qvotient som: $\frac{1}{1-1}$ maae reflectere sig i hinanden.
% --- p. 53 ---
\setcounter{footnote}{0}%
\opage{53}Er nu Forestillingen incommensurabel med Begrebet, maa
jo Articulationen af Uendelighedens Indhold blive en Articulation
af den uendelige Tilnærmelse; en saadan Articulation
er betinget af, at Tilvæxternes forskjellige Betydning i
de forskjellige Functionssystemer paavises, samt deraf, at de
uendelige Rækkers Convergens eller Divergens exact angives
ved Bestemmelse af Rester og Grændser.

At en endelig Tilvært $(h),$ hvormed $x$ i $f(x)$ forøges,
maa medføre en til Functionens Beskaffenhed svarende Tilvært
til Functionen selv, og denne Tilvæxt $(k)$ altsaa være
forskjellig for de forskjellige Functioner, er en Selvfølge.
Sættes $f(x)=y,$ og $f(x+h)=y+k,$ da indsees, at
medens $k,$ for $f(x+h)=a(x+h),$ bliver til $ah,$ og kun
udgjør eet Led, og for $f(x+h)=(x+h)^{a},$ for $a$ positiv
heel, til $ax^{a-1}\frac{h}{1}+a(a-1)x^{a-2}\frac{h^{2}}{1.2}+\ldots h^{a},$ og giver en
Fleerhed af Led, nemlig $a$ Led, maa $k$ t. Ex. for $f(x+h)=a^{x+h}$ blive til $a^{x}+a^{x}(la)\frac{h}{1}+a^{x}(la)^{2}\frac{h^{2}}{1.2}+\ldots$ og udvikle sig i utallige Led. Men sees nu paa Rækkerne med
de utallige Led, altsaa paa de uendelige Rækker, da synes
en exact Udvikling af deres Værdier ved første Øiekast umulig:
skal en uendelig Række udfoldes, maa den jo føres til
Ende, men det er umuligt; kan den uendelige Række aldrig
føres til Ende, maa der jo blive en Rest tilbage, der ikke
udfoldes; men bliver der en Rest tilbage, der ikke udfoldes,
maa Tilnærmelsen jo, saa synes det, tabe sig i det Omtrentlige.

I den til $\frac{1}{1-x}$ svarende Række $1+x+x^{2}+\ldots.\ x^{n-1}$ kan man jo gjerne tilføie Resten $\frac{x^{n}}{1-x};$ dog, hvorfor
blive staaende ved $x^{n-1}?$ man kunde jo ligesaa godt gaae til $x^{2n-1}$ og saa tilføie Resten $\frac{x^{2n}}{1-x};$ men, naar Afbrydelsen

% --- p. 54 ---
\setcounter{footnote}{0}%
\opage{54}er saa vilkaarlig, hvorledes kan Resultatet da hæve sig over
det Omtrentlige? Den Rest, der saaledes angives, være
sig nu som $\frac{x^{n}}{1-x}$ eller som $\frac{x^{n+1}}{1-x}$ eller som $\frac{x^{2n}}{1-x}$ indeholder i sig et Noget, som endnu ikke er udviklet.
Behøver det i Resten Indeholdte ingen Udvikling, saa behøver Qvotienten
$\frac{1}{1-x}$ vel heller ingen Udvikling, og hele
Rækken bliver overflødig; behøver Resten, hvorvidt man end
gaaer, bestandig en Udvikling, da bliver her et Noget tilbage,
som aldrig kan udvikles. Vil man bestemme Restens
Værdi ved at mærke sig Forskjellen imellem saadanne uendelige
Rækker, hvori hvert følgende Led er større end sit foregaaende,
og saadanne, hvori det Omvendte finder Sted:
da synes heller ikke denne Fremgangsmaade at kunne i hvert
Tilfælde føre til et exact Resultat. Sættes t. Ex. $x=2$ i $\frac{1}{1-x},$ saa at Rækken bliver $1+2+2^{2}+2^{3}+\ldots.;$ sættes dernæst $x=\frac{1}{2},$ saa at Rækken bliver $1+\frac{1}{2}+\left(\frac{1}{2}\right)^{2}+\left(\frac{1}{2}\right)^{3}+\ldots,$ da indseer man let, at Resten i den
første af disse Rækker vil staae i et ganske andet Forhold til
Leddenes Sum, end Resten i den anden. I første Række
vil Resten, hvorlænge man end fortsætter Divisionen, beholde
en uendelig stor Værdi, medens den i sidste Række vil aftage.
For Rækken $1+\frac{1}{2}+\left(\frac{1}{2}\right)^{2}+\left(\frac{1}{2}\right)^{3}+\ldots.$ bliver Resten
imidlertid ikke bestemt derved, at vi fortsætte Divisionen,
men derved, at vi ad `anden Vei sættes i Kundskab om den
til hele Rækken svarende Værdi. Hvor lidet man nu kan
stole paa en stadig Formindskelse af Resten, fordi man har
en Række med stadig aftagende Led, kan sees af et andet
Tilfælde, hvori vi ligeledes ad særegen Vei ere satte i

% --- p. 55 ---
\setcounter{footnote}{0}%
\opage{55}Kundskab om Rækkens sande Værdi. Da nemlig $a^{-\infty}=\left(\frac{1}{a}\right)^{\infty}=0,$ saa er $\mp\infty=\mathrm{Log}\,0;$ men $\mathrm{Log}\,0=\mathrm{Log}\,(1-1)=-\mathrm{Log}\,e\left(1+\frac{1}{2}+\frac{1}{3}+\ldots.\frac{1}{n}+\ldots\right),$ og, i det naturlige % PRINTED AS IS: „Da nemlig“ (Rettelser: Er nemlig)
System, $l(0)=-\left(1+\frac{1}{2}+\frac{1}{3}+\ldots.\frac{1}{n}\ldots.\right)=-\infty;$ heraf sees, at ikke enhver Række med aftagende Led giver en
stedse aftagende Rest. Føies nu hertil, at der jo ere mangfoldige
Rækker, hvis Rester, som f. Ex. Resten for anførte
logarithmiske, eller som Resten for den til $e$ svarende
Række: $1+\frac{1}{1}+\frac{1}{1.2}+\ldots$ ikke engang lade sig fremstille
under endelig Form: da maa det vistnok forekomme
den over Punktualiteten svævende Reflexion, som fik Omtrentligheden
overalt et betænkeligt Spillerum: ved Spændingen
mellem det i Tilnærmelsen Omtrentlige og det Exacte,
d. e. ved den i „Articulationen“ skjulte Modsigelse, er
Uendelighedsproblemet charakteriseret.

I hvilken Grad Uendelighedsproblemet har kunnet forbause
endog den med Analysen fortrolige Mathematiker, naar
han for Alvor begyndte at ville gjennemtænke det, er paa en
ret slaaende Maade anskueliggjort i følgende fra \textit{les Elémens
de Mathématique du P. Lami} hentede Indklædning: \textit{»Le
mauvais Riche brûlé de soif, prie, dit on, Abraham de
lui laisser distiller une goute d'eau. Abraham obtient
ce qu'il demande. Il laisse tomber une goute d'eau,
qui fait 100 lieuës à la premiere minute, 99 à la seconde,
ainsi de suite, toujours en même proportion à
chaque minute de 100 à 99. Or la distance des lieux
où sont Abraham \& le mauvais Riche étant supposée
infinie, on demande en combien de tems cette goute
arrivera au mauvais Riche? On répond à cela, jamais.
En effet, l'espace étant infini, la goute d'eau ne peut
pas parvenir à un terme. On conçoit bien que la chose
ne peut être autrement. Mais ce qui a droit de sur}%
% italic French quotation continues on next page (set as \textit there)

% --- p. 56 ---
\setcounter{footnote}{0}%
\opage{56}\textit{prendre, c'est qu'on détermine le nombre des lieues que
feroit la goute en tombant pendant une éternité; \& ce
nombre est 100,000 lieues. Voila donc une quantité
finie déterminée, quoiqu'il n'y ait point de terme auquel
elle soit limitée, le tems de la chûte d'eau étant infini.
Cela est assurement inconcevable. Afin de rendre la
chose sensible, on fait observer, que la chûte de la
goute d'eau est rétardée, suivant l'hypothese, à chaque
minute. Ce rétardement augmentant toujours, il devient
à la fin si peu considérable, qu'on peut le regarder
comme nul. Alors la goute d'eau ne se meut pas;
puisqu'on néglige son mouvement, qui est infiniment petit.
La question est de savoir, si on peut négliger cet Infiniment-Petit, \& si dans un tems infini, ce rétardement,
quoique toûjours plus petit, n'a pas une valeur réelle.
Notre esprit se perd dans ce raisonnement. Nous concevons
que la progression decroissant toûjours, la valeur
des termes doit diminuer tellement qu'elle ne soit pas
sensible, \& cependant nous voulons que le nombre des
termes, qui composent cette progression, soit infini.
Ceci est tout-à-fait contradictoire.«}
\textit{Application de la
Géométrie ordinaire et des Calculs différentiel et intégral
à la résolution de plusieurs Problêmes. Par feu
M. Robillard. Paris MDCCLIII. (Histoire critique du Calcul
des Infiniment-Petits contenant la Métaphysique et la
Théorie de ce Calcul. P. XVII--XVIII.)}

\begin{center}III.\end{center}

\begin{center}Det Uendeliges Analyse.\end{center}

At en uendelig Tilnærmelse dog kan ende, det Uopnaaelige
opnaaes, og det Omtrentlige opløse sig i det Exacte,
viser det Uendeliges Analyse. For at gjøre Uendelighedsproblemet
kjendeligt, maa Reflexionen støtte sig til

% --- p. 57 ---
\setcounter{footnote}{0}%
\opage{57}Operationen; for at løse Problemet, maa den tilegne sig
den i Operationen udtalte Intelligens. At der imellem
uendelige Størrelser kan være et endeligt Forhold, at en
Sum af uendelig mange Led kan i nogle Tilfælde udgjøre
en endelig, i andre en uendelig Størrelse, at en Række,
hvis Led forholde sig saaledes, at det efterfølgende divideret
med det foregaaende, giver Qvotienten $1,$ kan være enten
convergerende eller divergerende o. s. v., Sligt vil overalt
undgaae den Tænkning, der ikke ved Operationens Vink
faaer Øie paa de exacte Bestemmelser. At nu disse Bestemmelser
i philosophisk Form maae blive dialektiske, medens
de i mathematisk Form hverken ere, eller kunne være det,
viser en charakteristisk Forskjel imellem Mathematik og Philosophie,
men ingenlunde nogen Strid. Intet Begreb kan,
hvor exact det end defineres, holde sig i en mod de øvrige
Begreber ligegyldig Væren: Begrebernes Gyldighed sees i
deres Overgange; men de Overgange, den exacte Analyse
netop gjør ved Operationen, gjør den Operationsbegreberne
gjennemtrængende Tænkning ved Dialektiken.

Den exacte Bestemthed, hvori det Endelige berører det
Uendelige: \emph{Functionsgrændsen}; det exacte Punktualitetsudtryk,
hvorved det Endelige opløses i det Uendelige som i
sit oprindelige Moment: \emph{Differentialet}; den exacte
Form, hvorunder den uendelige Discretions Elementer forenes
til et inden mulige Grændser fuldendt Continuum,
saa at det Endelige derved kan restitueres af det Uendelige:
\emph{Integralet,} ere da de Grundbegreber, Dialektiken i denne
Sammenhæng har at belyse.

\begin{center}a) Functionsgrændsen: det Endeliges dialektiske Berøring med det Uendelige.\end{center}

Den Grændse, hvortil en Function nærmer sig, idet
dens Uafhængige enten voxer eller aftager i det Uendelige,
bliver ved Reflexionen sat paa een Gang baade som vor%

% --- p. 58 ---
\setcounter{footnote}{0}%
\opage{58}dende og som værende, og kan forøvrigt være snart $0,$ snart $\infty,$ snart endelig. Betegner man saaledes Grændsen for $y$ i $f(x)=y,$ naar $x=\infty,$ ved $\mathrm{Lim}\,f(x)=f(\infty),$ og, idet $x=0,$ ved $\mathrm{lim}\,f(x)=f(0);$ da indsees, at her til $f(x)=ax=y$ maa svare $\mathrm{Lim}\,f(x)=\infty$ og $\mathrm{lim}\,f(x)=0,$ til $f(x)=\frac{a}{x}=y,$
$\mathrm{Lim}\,f(x)=0$ og $\mathrm{lim}\,f(x)=\infty,$ til $f(x)=\frac{x}{x}=y$
derimod $\mathrm{Lim}\,f(x)=\mathrm{lim}\,f(x)=1.$ Hvorledes det
gaaer til, at den exacte Handling ikke kan ændse det Dialektiske,
uagtet den gjennemgaaende Tænkning med Nødvendighed
maa opdage det Dialektiske, kan i nærværende
Sammenhæng paavises. Har man t. Ex. $\mathrm{Lim}\,f(x)=\mathrm{Lim}\frac{a_{0}+a_{1}x+a_{2}x^{2}\ldots+a_{n}x^{n}}{b_{0}+b_{1}x+b_{2}x^{2}\ldots+b_{n}x^{n}},$ og sætter uden videre $x=\infty,$ faaer man i $\frac{a_{0}+\infty+\infty\ldots}{b_{0}+\infty+\infty\ldots}=\frac{\infty}{\infty}$ et aldeles ubestemt Resultat; ligeledes vil, dersom man i $\mathrm{lim}\,f(x)=\mathrm{lim}\frac{a_{1}x+a_{2}x^{2}\ldots+a_{n}x^{n}}{b_{1}x+b_{2}x^{2}\ldots+b_{n}x^{n}}$ uden al Forberedelse sætter $x=0,$ Resultatet: $\frac{0+0+\ldots}{0+0+\ldots}=\frac{0}{0}$ forblive ubestemt.

Hvor aldeles Bestemtheden er opløst, hvor svævende Grændsen
er, sees i anførte Tilfælde deraf, at de vundne Resultater,
skiøndt indbyrdes forskjellige, kunne omdannes til skuffende
indbyrdes Lighed, nemlig $\mathrm{Lim}\,f(x)=\frac{\infty}{\infty}=\frac{1}{\infty}\infty.=0\infty,$ og $\mathrm{lim}\,f(x)=\frac{0}{0}=0.\frac{1}{0}=0.\infty.$ Ved en
forberedende Handling, en bestemt Reduction,
kommer man derimod til en aldeles bestemt Grændse.
\[
\mathrm{Lim}\,f(x)=\mathrm{Lim}\frac{\frac{a_{0}}{x^{n}}+\frac{a_{1}}{x^{n-1}}\ldots+a_{n}}{\frac{b_{0}}{x^{n}}+\frac{b_{1}}{x^{n-1}}\ldots+b_{n}}=\frac{a_{n}}{b_{n}}=c,
\]
og

% --- p. 59 ---
\setcounter{footnote}{0}%
\opage{59}\[
\mathrm{lim}\,f(x)=\frac{a_{1}+a_{2}x+\ldots a_{n}x^{n-1}}{b_{1}+b_{2}x+\ldots b_{n}x^{n-1}}=\frac{a_{1}}{b_{1}}=d,
\]
som ikke kan forvandles til $c.$ Sætter man, for at anføre
endnu et Exempel, i $\mathrm{Lim}\left(\frac{1}{x^{x}}\right)^{\frac{1}{x}}$ uden videre $x=\infty,$ faaer
man kun $\left(\frac{1}{\infty}\right)^{\frac{1}{\infty}}=0^{0},$ som i dette Tilfælde er ubestemt;
foretager man derimod først den fornødne Reduction erholder
man $\mathrm{Lim}\left(\frac{1}{x^{x}}\right)^{\frac{1}{x}}=\mathrm{Lim}\frac{1}{x}=0,$ som er en bestemt Grændse.

De exacte Handlinger følge altsaa paa hianden i en % PRINTED AS IS: hianden (Rettelser: hinanden; also „Reduction“ for „Reduction,“ in the line above)
vis Orden: de ere, saa at sige, intellectuelle Foretagender,
der udføres i Tiden: først skeer det Ene, derefter det Andet;
først anvendes den ene Betragtningsmaade, derefter den
anden; først behandler man $x$ som endelig Størrelse, siden
sættes den som $\infty$ eller som $0;$ først antager man Grændserne $\infty$ og $0$ for uopnaaelige, sidenefter, naar man ved
behørig Reduction af de endnu endelige Led har gjort Oprindelsens
Lov kjendelig, betragter man Sagen, som om
den for $x$ i $f(x)=y$ uopnaaelige Grændse dog alligevel var
opnaaet, og indseer saa, hvilken Indflydelse dette maa have
paa $y.$ Det er denne de exacte Handlingers Adskillelse i
Tiden, der lader de modsatte Betragtningsmaader optræde
ved Siden af hinanden, vexle med hinanden, falde ud fra
hinanden, og saaledes forhindrer det Dialektiske fra at komme
tilsyne.

I den gjennemgaaende Reflexion falder Tidsadskillelsen
bort; den endelige Udstykning forsvinder i Begrebets Eenhed,
de modsatte Bestemmelser gjennemskinne hinanden, og det
Dialektiske viser sig. Ligningerne: $\mathrm{Lim}\,f(x)=f(\infty)=c,$ $\mathrm{lim}\,f(x)=f(0)=d$ udsige, naar de i Begrebet opfattes
som logiske Domme, at det Uopnaaelige er opnaaeligt.
Denne Modsigelse, der, løsreven fra sin Forudsætning, er

% --- p. 60 ---
\setcounter{footnote}{0}%
\opage{60}ligesaa utaalelig, som den Modsigelse at Firkanten er rund,
har i sin væsentlige Sammenhæng den Betydning, at den
sammentrænger to modsatte, hinanden udelukkende Bestemmelser
i een Grundbestemmelse, ikke for at gjøre dem ligefrem
til Eet, men for at vise deres fælleds Udspring. At
Grændsen for $x,$ og følgelig ogsaa Grændsen for $f(x)$ i $\mathrm{Lim}\,f(x)=f(\infty)=c$ er uopnaaelig, er jo netop begrundet i $f(x),$ som fordrer, at de utallige Mellemled, hvorved $x$ som endelig Størrelse adskiller sig fra $\infty$ skulle gjennemløbes;
men Ligningen selv, hvori der end ikke er Skygge af
det Omtrentlige, angiver ved den exacte Værdi $c$ tillige den
opnaaede, altsaa dog opnaaelige Grændse. Den Forklaring,
at Grændsen er naaet ved et Spring, idet man har sprunget
utallige Mellemled over, er eensidig og vildledende; i Ligninger
som: $\mathrm{Lim}\,f(x)=c$ og $\mathrm{lim}\,f(x)=d,$ er, væsentlig
seet, intet Mellemled oversprunget; den samme Lov, der i $\mathrm{Lim}\,f(x)=c$ gjælder for $x=\infty,$ gjælder ogsaa for $x$ som
endelig, og den Lov, der i $\mathrm{lim}\,f(x)=d$ gjælder for $x$ som
endelig, gjælder i Conseqvensen ogsaa for $x=0.$ Sætningen: det Uopnaaelige er opnaaeligt, har sin Grund i Functionens
Væsen; naar denne Begrundelse er indseet, forsvinder
Modsigelsen, thi de contradictorisk modsatte Bestemmelser
blive kun forenede for igjen at adskilles, og i endelig
Adskillelse at udelukke hinanden.

I hvilken Betydning det Uopnaaelige udelukker Opnaaelsen, kan sees af de uendelige Rækkers Afbrydelse;
i hvilken Betydning det Uopnaaelige dog bliver opnaaet,
kan sees af disse Rækkers væsentlige Afslutning;
hvor inderlig Begrændsningens modsatte Bestemmelser gjennemtrænge hinanden, kan endelig sees i de
uendelige Størrelser.

\begin{center}α)

Uendelige Rækkers Afbrydelse: Grændsen som uopnaaelig.\end{center}

Da en uendelig Række ikke kan føres til Ende, maa
den afbrydes; den Rest, som herved fremkommer, er en

% --- p. 61 ---
\setcounter{footnote}{0}%
\opage{61}ubekjendt Størrelse. Her møder altsaa den ovenfor berørte
Vanskelighed, at man, for at kjende den fuldstændige Sum;
forud maatte kiende Resten, og omvendt, for at kjende Resten,
først maatte kjende den fuldstændige Sum. Denne
Vanskelighed hæves, naar Leddenes Forhold under en endelig
Fortsættelse viser, til hvilken Grændse den uendelige
Fortsættelse vilde føre. Det er Forskjellen imellem convergente
og divergente Rækker, som her afgiver det bestemte
Udgangspunkt: „nærmer $s_{n},$ i Summen $s_{n}=u_{0}+u_{1}\ldots+u_{n},$ sig for $n$ bestandig voxende, en bestemt endelig
Grændse, siges Rækken selv at være convergent; er en Række
ikke convergent, siges den at være divergent.“ Den exacte
Opfattelse er nu i Korthed følgende:

„I enhver convergent Række er $\mathrm{Lim}\,u_{n}=0,$ men dette
er alene en nødvendig, ikke en tilstrækkelig Betingelse for
Convergens, thi et uendeligt Antal af positive Led convergerende
til $0$ kan give en Sum $=\infty.$ Derimod naar $r_{n}$ betyder Summen af alle de efter $u_{n}$ følgende Led $u_{n+1},u_{n+2},u_{n+3},\ldots.$ i det Uendelige, hvilket kaldes Resten
af Rækken, saa vil $\mathrm{Lim}\,r_{n}=0$ udtrykke baade en nødvendig
og tilstrækkelig Betingelse for Rækkens Convergens, thi
man kan da stedse gjøre $n$ saa stor, at $r_{n}$ er mindre end
en hvilkensomhelst opgiven Størrelse, hvorved $s_{n}$ bliver
bragt ligesaa nær, som man vil, til en vis Grændse.“\footnote[1]{Ramus. Algebra og Functionslære. S. 80.}
Har man nu ved Anvendelse af dertil fundne Udtryk $\mathrm{Lim}\,\frac{u_{n+1}}{u_{n}}$ og $\mathrm{Lim}\,n\left(\frac{u_{n}}{u_{n+1}}\div1\right)$ forvisset sig om en Rækkes
Convergens, da vil en Paaviisning af Convergensen igjen
kunne føre til en nærmere Bestemmelse af Resten. Hvad
Qvotienten $\frac{1}{1-x}$ angaaer, tilkjendegiver Udtrykket $\frac{1}{1-x}=1+x+x^{2}+x^{3}+\ldots.\left\{\begin{array}{l}x=-1\\x=+1\end{array}\right.$ „at denne Form

% --- p. 62 ---
\setcounter{footnote}{0}%
\opage{62}tilfredsstiller Betingelsen $\mathrm{Lim}\,r_{n}=0,$ naar $x$ har en Værdi
indesluttet imellem Grændserne: $\pm1$ exclusive. Giver nu $u_{n}=x^{n}$ som Rest $r_{n}=\frac{x^{n+1}}{1-x}=u_{n}\frac{x}{1-x},$ saa indsees,
„at saasnart i en Række Leddene aftage saa stærkt, at $\frac{u_{n+1}}{u_{n}}<x$ for en vis Index $n$ og for alle høiere, saa er
ogsaa $r_{n}<u_{n}\frac{x}{1-x}.$“ „Dette“, siger Ramus, „kan tjene
til at bestemme den høiere Grændse for Feilen, som begaaes
derved, at en uendelig Række afbrydes ved $u_{n}$.“

I den exacte Analyse begaaer man altsaa en Feil ved
at afbryde en uendelig Række; Index $n$ kan i $u_{n}$ være saa
stor, man vil: saalænge $n$ er endelig, medfører Afbrydelsen
en Feil. Saaledes giver ovenstaaende Række for $x=\frac{1}{2},$ $r_{n}=u_{n},$ for $x=\frac{1}{3},$ $r_{n}=\frac{1}{2}u_{n}\ldots.$ for $x=\frac{1}{p},$ $r_{n}=\frac{1}{p-1}u_{n}.$ Er $\frac{u_{n+1}}{u_{n}}$ ikke constant, men „for en vis Index $n$ og for alle høiere“ stadig aftagende, vil Resten naturligviis
blive endnu mindre. Man kan i slige Tilfælde bestemme
Feilen, forsaavidt man kan bestemme Resten; men
at her altid bliver en Rest, beviser, at Rækkens Afbrydelse
altid giver en Feil. For Rækken: $e=1+\frac{1}{1}+\frac{1}{1.2}+\ldots,$ t. Ex. har man nu vistnok beviist, baade hvad Summen
angaaer, at $2<e<3,$ og at Resten $r_{n}<\frac{1}{1.2.3\ldots.n.n'},$ naar Rækken er afbrudt ved $u_{n};$ men man har da ogsaa
med det Samme beviist Umuligheden af at undgaae en
Feil, idet man nemlig har beviist, at hele Størrelsen og
følgelig ogsaa Resten er incommensurabel med Eenheden. I
dette Punkt er Modsætningen mellem Operations- og Reflexionsbegrebet
ret paafaldende. Operationsbegrebet, der

% --- p. 63 ---
\setcounter{footnote}{0}%
\opage{63}er anlagt paa Beregning, tilfredsstilles exact derved, at
Feilens Grændser bestemmes; Reflexionsbegrebet, der er
anlagt paa logisk Indseen, fordrer en dialektisk Tilfredsstillelse.
Forsaavidt den exacte Analyse maa indskrænke sig
til at sætte Grændser for Feilen uden at kunne hæve Feilen,
indrømmer den selv det Ufuldkomne, opgiver sin exacte
Charakteer, og nøies med det Omtrentlige; men idet Analysen
saa igjen angiver Forholdet imellem større og mindre
Feil, bliver sig Feilen tydelig bevidst, og veed at sætte den
saa snævre Grændser, det skal være, har den jo netop
Feilen i sin Magt, og kan altsaa rette Feilen\footnote[1]{Rækken for den Newtonske Binomialformel, $(a+x)^{n}=a^{n}+\frac{n}{1}a^{n-1}x+\ldots.,$ som, naar n er negativ eller bruden, bliver uendelig, er convergent, naar $\mathrm{Lim}\frac{A_{r+1}}{A_{r}}x=\mathrm{Lim}\left(\frac{n-r}{r+1}\right)\frac{x}{a}=\frac{x}{a}<1,$ altsaa, naar $a>x>-a.$ Benyttes Rækken til „approximeret Beregning, ibet man f. Ex. tager Summen af p Led, bliver Resten af Rækken eller Feilen, der begaaes, udtrykt ved: $r_{p}=A_{p}a^{n-p}x^{p}\left\{1+\frac{n-p}{p+1}\frac{x}{a}+\frac{(n-p)(n-p-1)}{(p+1)(p+2)}\left(\frac{x}{a}\right)^{2}+\ldots\right\}<A_{p}a^{n-p}x^{p}\left(1+\frac{x}{a}+\left(\frac{x}{a}\right)^{2}+\ldots\right)=A_{p}a^{n-p+1}\frac{x^{p}}{a-x}\left(=A_{p}a^{n-p}x^{p}\frac{a}{a-x}\right)$ „ibet man blot iagttager, at $n-p<p+1$ eller $n-1<2p,$ hvilket altid kan opnaaes ved at fortsætte Rækken langt nok.“ A. Steen. Reen Mathm. S. 307—308.

„En nødvendig Forberedelse for de reelle Rødders approximerede Beregning i numeriske Ligninger er deres Udskillelse d. e. Bestemmelsen af de forskjellige Intervaller, hvori de ere enkelt beliggende. Er $f(x)=0$ en Ligning med blot ulige Rødder, vil Udskillelsen kunne opnaaes ved i $f(x)$ at substituere for x de forskjellige mellem de yderste Grændser beliggende reelle Størrelser, som ere Led i en arithmetisk Række, hvis Differens er mindre end hvilkensomhelst to Rødders Differens i $f(x)=0$.. (Ramus. Algebra og Functionslære. S. 129.) Haves altsaa i Ligningen $f(x)=0$ (I) en Rod fuldstændig udskilt t. Ex. imellem $a_{0}$ og $a_{0}+1$, da maa Ligningen tilfredsstilles ved $x=a_{0}+\frac{1}{x_{1}}$ $(\alpha)$. Indsættes denne Værdi i (I), faaes idetmindste $f(x_{1})=0$ (II); heri maa være en reel Rod $>1$, men ogsaa kun een, da vi jo gik ud fra, at der imellem $a_{0}$ og $a_{0}+1$ var en Rod fuldstændig udskilt: $x_{1}$ kan med andre Ord, kun have een Værdi $>1$. Udskilles nu denne Rod fuldstændig, finde vi den beliggende mellem $a_{1}$ og $a_{1}+1$, d. e. $a_{1}<x_{1}<a_{1}+1$, saa at $x_{1}=a_{1}+\frac{1}{x_{2}}$ $(\beta)$. Indsættes $(\beta)$ i $(\alpha)$, faaer man: $x_{0}=a_{0}+\cfrac{1}{a_{1}+\cfrac{1}{x_{2}}}$ $(\gamma)$. Indsættes Udtrykket $(\beta)$ i (II), faaes $f(x_{2})=0$ (III), som ligeledes kun har een reel Rod, der antages fuldstændig udskilt mellem $a_{2}$ og $a_{2}+1$, saa at $a_{2}<x_{2}<a_{2}+1$, og $x_{2}=a_{2}+\frac{1}{x_{3}}$. Ved at indsætte det sidste Udtryk i $(\gamma)$ erholder man $x=a_{0}+\cfrac{1}{a_{1}+\cfrac{1}{a_{2}+\cfrac{1}{x_{3}}}}$ osv. Her danner sig altsaa, som man seer, en Kjædebrøk, der kan fortsættes i det Uendelige; saaledes haves $\sqrt{2}=1+\cfrac{1}{2+\cfrac{1}{2+\ldots}}$ Afbryder man nu en saadan Kjædebrøk, og antager to paa hinanden følgende Convergenter til den ved Tilnærmelsen søgte Størrelse x at være $\frac{y_{r-1}}{z_{r-1}}$ og $\frac{y_{r}}{z_{r}}$, da haves $\frac{y_{r-1}>\ \ >y_{r}}{z_{r-1}<x<z_{r}}$. Man begaaer nu unegtelig en Feil, hvor man saa end afbryder den uendelige Kjædebrøk, men den Feil man begaaer, ved t. Ex. at tage $\frac{y_{r-1}}{z_{r-1}}$ for x, er i numerisk Værdi mindre end $\frac{y_{r}}{z_{r}}-\frac{y_{r-1}}{z_{r-1}}=\frac{1}{z_{r-1}z_{r}}$, og forsaavidt nøiagtig bestemt. (Jvfr. Ramus. Elementær Algebra, S. 103—112. og Algebr. Functionsl. S. 25—35).}; den exacte

% --- p. 64 ---
\setcounter{footnote}{0}%
\opage{64}Analyse befinder sig saaledes, naar den sees fra Reflexionens
Synspunkt, i det dialektiske Tilfælde, at den ved at
afbryde de uendelige Rækker begaaer Feil uden at tage feil,
at den feiler uden at feile.
% the note 1) of p. 63 runs over onto the foot of this page: set in full at its mark

% --- p. 65 ---
\setcounter{footnote}{0}%
\opage{65}
\begin{center}β)

Rækkernes væsentlige Afslutning: Grændsen \emph{som opnaaelig.}\end{center}

At det ved en uendelig Rækkes væsentlige Afslutning
kommer an paa Indseen og Begriben er ligesaa uimodsigeligt,
som at det Punktlighedsbegreb, hvorom det her gjælder,
er forskjelligt fra et hvilketsomhelst almeenlogisk Begreb.
Betegnelser som: $u_{n},$ det nte Led, $\mathrm{Lim}\frac{u_{n+1}}{u_{n}},$
$\mathrm{Lim}\,r_{n}$ o. desl. ere, væsentlig forstaaet, mere end vedtagne Beregningsudtryk;
de ere Tankehandlingens prægede Former,
Operationskategorier, hvori det Endelige paa en saa eiendommelig
Maade gjennemskinnes af det Uendelige, at det
Endeløses Vidtløftighed hæves, og Rækken kommer til Afslutning.
Hvorledes det nte Led i Begrebets Almindelighed
afspeiler de forestillede Enkeltleds uendelige Vrimmel,
er ved Bestemmelsen af Forholdet mellem Operationen og
den objective Reflexion allerede fremhævet. For at fatte
det Væsentlige i en uendelig Rækkes Afslutning vil det
imidlertid her være nødvendigt at betragte Forholdet imellem
Operations- og Reflexionsbegrebet lidt nærmere.

Til Bedømmelse af en Rækkes Convergens opstilles,
som bekjendt, den Læresætning, at Rækken er convergent,
naar $\mathrm{Lim}\frac{u_{n+1}}{u_{n}}<1,$ og derimod divergent, naar $\mathrm{Lim}\frac{u_{n+1}}{u_{n}}>1.$ Hvorvel man nu vistnok kan besidde en vis Færdighed
i at bruge de opstillede Formler, og forsaavidt have de
deri udtrykte Operationsbegreber, uden endnu at gjennemskue
Sagens Væsen, saa ligger dog Reflexionsbegrebet i nærværende
Tilfælde saa nær, at det under fortsat Udøvelse af
Operationen omsider maa komme som af sig selv. At nemlig
en Række, hvis Led saaledes vore, at $\mathrm{Lim}\frac{u_{n+1}}{u_{n}}>1,$ maa
divergere, er iøinefaldende; heller ikke kan den Forstaaelse
ligge fjernt, at en Række, hvis Led aftage saaledes, at

% --- p. 66 ---
\setcounter{footnote}{0}%
\opage{66}$\mathrm{Lim}\frac{u_{n+1}}{u_{n}}<1,$ maa convergere. I hvilken Grad Opera-
tionskategorierne bidrage til at klare den objective Reflexion,
kan her sees ved en Sammenstilling af saadanne Rækker
som: $1+\frac{1}{2}+\left(\frac{1}{2}\right)^{2}+\ldots$ og $1+\frac{1}{2}+\frac{1}{3}+\ldots$ Da Leddene
i begge disse Rækker ere aftagende, skulde man ved
første Blik formode, at begge maatte være convergente, den
sidste ligesaavel som den første. Ved Operationens Formler
oplyses vi imidlertid ikke blot om det Modsatte, men vi
ledes tillige til at indsee Grunden. Betragte vi de almindelige
Led $\left(\frac{1}{2}\right)^{n}$ og $\frac{1}{n}$ lidt nøiere, overbevises vi om,
at $\left(\frac{1}{2}\right)^{n+1}$ maa, naar $n$ er overordentlig stor, være i
Forhold til $\left(\frac{1}{2}\right)^{n}$ ganske anderledes forsvindende, end $\frac{1}{n+1}$ i Forhold til $\frac{1}{n}.$ Rækken $1+\frac{1}{2}+\left(\frac{1}{2}\right)^{2}\ldots+\left(\frac{1}{2}\right)^{n}+\ldots$ staaer altsaa Formindskelsen af hvert følgende Led i
et saadant Forhold til Formerelsen af Leddenes Antal, at
det Første i det Uendelige ophæver det Sidste, og den hele
Sum falder indenfor endelige Grændser. I Rækken: $1+\frac{1}{2}+\frac{1}{3}\ldots+\frac{1}{n}\ldots$ bliver derimod Formindskelsen af hvert følgende
Led desto ubetydeligere, jo mere Leddenes Antal
voxer; hvor stor $n$ end bliver, vil $r_{n}$ paa disse Vilkaar ved
sine utallige Led bestandig kunne give en Sum $=\infty,$ og
Rækken altsaa være divergent. Det bliver da paa denne
Maade indlysende, at en Række, hvori $\mathrm{Lim}\frac{u_{n+1}}{u_{n}}=1,$ kan,
saaledes som i det givne Exempel, divergere, medens en
Række, hvori $\mathrm{Lim}\frac{u_{n+1}}{u_{n}}<1$ maa, ifølge Sagens Natur

% --- p. 67 ---
\setcounter{footnote}{0}%
\opage{67}være convergent. Her kan i strengeste Forstand være Tale
om Indseen og Begriben.

En saa fuldstændig Congruens af Reflexion med
Operation finder derimod ikke Sted, hvad Udtrykket: $\mathrm{Lim}\ n\left(\frac{u_{n}}{u_{n+1}}-1\right)>1$ angaaer. En Række, hvori Leddene
forholde sig saaledes, at $\mathrm{Lim}\ \frac{u_{n+1}}{u_{n}}=1$, kan, som vi have
seet, være divergent, men den kan paa given Betingelse
ogsaa være convergent. Naar da en saadan Rækkes Convergens
skal undersøges, vælger man til Norm en anden
Række, hvis Convergens særlig lader sig paavise, uafhængig
af almindelige Theorier, og hvori t. Ex. $\mathrm{Lim}\ \frac{t_{n+1}}{t_{n}}=1$, idet $t_{n}$ er det almindelige Led. En saadan Rækkes
nte Led være nu: % PRINTED AS IS: nte Led (Rettelser: mte Led)
\[ \frac{1.2.3\ldots m}{a(a+1)..(a+m-1)}\cdot\frac{a-1}{a+m} \]
\[ =\left(\frac{1.2.3\ldots m}{a(a+1)..(a+m-1)}-\frac{1.2.3..(m+1)}{a(a+1)\ldots(a+m)}\right) ; \]
sættes heri $m$ efterhaanden $=1,2,3\ldots n$, faaer man
Rækkens Sum udtrykt, paa den ene Side ved de enkelte
Led, paa den anden Side ved en Reduction, nemlig:
\[ \frac{a-1}{a}\left(\frac{1}{a+1}+\frac{1.2}{(a+\cdot 1)(a+2)}\ldots+\frac{1.2.3..n}{(a+1)\ldots(a+n)}+\ldots\right) \]
\[ =\frac{1}{a}-\frac{1.2.3..(n+1)}{a(a+1)\ldots(a+n)} ; \]
ved at sammenligne Summer
med flere og flere Led overbeviser man sig om, at Rækken er
convergent, for $a>1$, og hele Summen en ægte Brøk.
I denne Række haves da $\frac{t_{n+1}}{t_{n}}=\frac{n+1}{a+n+1}$, $\frac{t_{n}}{t_{n+1}}=\frac{a}{n+1}+1$, og altsaa, i Overeensstemmelse med den op-
stillede Fordring, $\mathrm{Lim}\ \frac{t_{n+1}}{t_{n}}=1$. Naar nu $\frac{u_{n}}{u_{n+1}}$ i en

% --- p. 68 ---
\setcounter{footnote}{0}%
\opage{68}Række, hvis almindelige Led er $u_{n}$, er større end, eller lig
med $\frac{t_{n}}{t_{n+1}}=\frac{a}{n+1}+1$ (hvori $a>1$), da er $u$-Rækken convergent, idet $t$-Rækken er det. Denne Convergens-
betingelse udtrykkes ved $\frac{u_{n}}{u_{n+1}}\geq\frac{a}{n+1}+1$; altsaa:
\[ \mathrm{Lim}\ n\left(\frac{u_{n}}{u_{n+1}}-1\right)+\mathrm{Lim}\left(\frac{u_{n}}{u_{n+1}}-1\right)\geq a, \]
eller
\[ \mathrm{Lim}\ n\left(\frac{u_{n}}{u_{n+1}}-1\right)\geq a. \]
Da her angaaende $a$ kun fordres
$a>1$, kan udtrykket $\mathrm{Lim}\ n\left(\frac{u_{n}}{u_{n+1}}-1\right)>1$ i fuld-
stændig Almindelighed angive Betingelsen for Convergens i
det omspurgte Tilfælde. Hermed er Operationsbegrebet
væsentlig afsluttet; forsøger Reflexionen at gaae videre,
viser sig en Skranke. Af
$\frac{t_{n}}{t_{n+1}}=\frac{a}{n+1}+1$ kan $\mathrm{Lim}\ n\left(\frac{t_{n}}{t_{n+1}}-1\right)$ unegtelig meget let afledes, og naar dette
Udtryk er fundet, falder dets Anvendelse paa $u$-Rækken, ifølge ovenstaaende Sammenligning, som af sig selv; men
opkaster man saa det Spørgsmaal, hvorledes det egentlig
er begrundet i $\frac{t_{n}}{t_{n+1}}$, at $\mathrm{Lim}\ n\left(\frac{t_{n}}{t_{n+1}}-1\right)$ maa give
$a>1$, for at Rækken kan være convergent, vil den Opererende, der kun ad særlig praktisk Vei seer sig istand til at
godtgjøre $t$-Rækkens Convergens, uden Tvivl indskrænke sig
til at pege paa den omhandlede Formel, og svare: „jeg
kan bevise det, men jeg kan ikke begribe det.“\footnote[1]{Hvorvidt den her antydede Vanskelighed er for Mathematiken selv uovervindelig eller ikke, lades i denne Sammenhæng uafgjort; kun dette maa udtrykkelig fremhæves, at den synthetiske Methode just hæver sin afsluttende Selvstændighed ved Udelukkelse af den analytiske i saadanne Tilfælde, hvor Beviisførelsen ikke kan hæves til Begrebsudvikling, og at den forsaavidt, istedetfor at betegne det Fuldkomnere, tvertimod beroer paa en Ufuldkommenhed.}

% --- p. 69 ---
\setcounter{footnote}{0}%
\opage{69}Ifølge Reflexionsbegrebet bliver den uendelige Række
væsentlig afsluttet derved, at Uendelighedsgrændsen paa alle
Punkter viser sig at berøre det Endelige. Ligninger som:
\[ \frac{1}{1-x}=\Sigma x^{n}+\frac{x^{n+1}}{1-x} \]
\[ e=1+\Sigma\frac{1}{1.2.3....n}+r_{n}, \]
idet $r_{n}<\frac{1}{1.2.3...nn}$, o.s.v., er det med Hensyn paa
Sagens Væsen ligegyldigt, om her i Summerne anføres
to Led, eller Myriader af Led. Uendelighedsgrændsen
er med alt dette lige fjern og lige nær, lige uopnaaelig,
lige opnaaelig. Dersom Sagen ikke forholdt sig saaledes,
vilde en uendelig Række vel kunne afbrydes, men aldrig
afsluttes. Sætter man t. Ex. i Rækken: $1+x+x^{2}+..+x^{n}+\frac{x^{n+1}}{1-x}$, $n=\infty$, maa jo ikke alene $x^{n+1}$ og $x^{n}$, men tillige $x^{n-1}, x^{n-2}, x^{n-3}\ldots$ forvandles til $x^{\infty}$; det Uendelige
vilde da paa denne Maade udøve en saa uhyre Tilbagevirkning,
at Rækken med Myriader af forskjellige Led tilsidst kun fik Led
af en eneste Form, nemlig $x^{\infty-\infty}$; og hvad er saa $x^{\infty-\infty}$?
Af $\frac{1}{1-x}$ faaes, for $x=\frac{3}{2}$, Rækken: $1+\frac{3}{2}+\left(\frac{3}{2}\right)^{2}....$
\[ +\left(\frac{3}{2}\right)^{n}+\frac{\left(\frac{3}{2}\right)^{n+1}}{1-\frac{3}{2}} ; \]
sættes nu $n=\infty$, bliver jo $\left(\frac{3}{2}\right)^{n}$
\[ =\infty, \left(\frac{3}{2}\right)^{n-1}=\infty, \left(\frac{3}{2}\right)^{n-2}=\infty \]
o.s.v., hele Summen
indtil $\left(\frac{3}{2}\right)^{n}$ inclusive $=\infty$, Resten $\frac{\left(\frac{3}{2}\right)^{n+1}}{1-\frac{3}{2}}=-\infty$; hvad

% the note 1) of p. 68 runs over onto the foot of this page: set in full at its mark

% --- p. 70 ---
\setcounter{footnote}{0}%
\opage{70}betyder dette? Af $\frac{1}{1-x}$ faaes, for $x=\frac{1}{2}$, Rækken $1+\frac{1}{2}$
\[ +\left(\frac{1}{2}\right)^{2}\ldots+\left(\frac{1}{2}\right)^{n}+\frac{\left(\frac{1}{2}\right)^{n+1}}{1-\frac{1}{2}} ; \]
sættes heri $n=\infty$, bliver
ikke alene $\frac{\left(\frac{1}{2}\right)^{n+1}}{1-\frac{1}{2}}=0$ og $\left(\frac{1}{2}\right)^{n}=0$, men $\left(\frac{1}{2}\right)^{n-1}$
\[ \left(\frac{1}{2}\right)^{n-2}, \left(\frac{1}{2}\right)^{n-3}.... \]
i det Uendelige $=0$; hvad er da $0$, hvad er $\infty$?

\begin{center}γ) \emph{Uendelige Størrelser: Begrændsningens Dialektik.}\end{center}

Ved skarp Adskillelse af Størrelsen selv og dens Grændse
kunne Betegnelserne $0$ og $\infty$ fastholdes, ikke som Størrelser,
men som Størrelsesgrændser saaledes, at $0$ udelukker og $\infty$ opsluger enhver mulig Størrelse. I et Udtryk som
\[ \frac{a}{0}=\infty \]
betyder da $0$ ikke, at her er, men at her har været
en Divisor, som nu er forsvunden; heller ikke kan $\infty$ betragtes
som en egentlig Qvotient, men som Tegn paa en i
det uendelige opløst Qvotient, en Qvotient, der ophørte at
være en Størrelse just i det Øieblik, da den tilsvarende
Divisor indsvandt til Intet. Men idet $0$ og $\infty$ paa denne
Maade opfattes som skarpe Negationer, som rene Grændser,
bliver Grændsebegrebet selv eensidig opfattet; Grændsen er
nemlig ifølge sit Begreb ikke blot negativ, men ogsaa positiv:
det Negative er Grændse derved, at det Positive er
det. Det er nemlig ikke blot $0$ og $\infty$, der ere Grændser;
Størrelser som $+1$ og $-1$ ere jo ogsaa Grændser, hine
negative, disse positive og negative. En Størrelse, der
falder imellem Grændserne $0$ og $\infty$ vil altsaa, ifølge den
angivne.Synsmaade, være tilintetgjort, naar den har naaet

% --- p. 71 ---
\setcounter{footnote}{0}%
\opage{71}en af disse Grændser; men en Størrelse, der falder imellem
Grændserne $+1$ og $-1$, vil netop ved at opnaae en af
disse Grændser vise sig som en bestemt tilværende Størrelse,
thi $-1$ er dog vel ligesaafuldt en Størrelse, som $+1$. Imellem Grændserne $+1$ og $-1$ er $0$ den Grændse,
hvori $+$ og $-$ berøre hinanden; at fastholde Adskillelsen
af to Slags Grændser, negative og positive, er altsaa umuligt,
efterdi det just er det Positive, som i Grændsen selv
grændser op til det Negative. Men er det umuligt at adskille
det i Grændsen Negative fra sit Positive, da er det
ogsaa umuligt at holde Størrelsens Grændse endelig adskilt
fra Størrelsen selv.

I Erkjendelsen af, at en Størrelses Grændse nødvendig
hører med til at bestemme den som Størrelse, kunde
man dog, uden endnu at opfatte $0$ og $\infty$ som Størrelser, prøve
paa at gaae et lille Skridt videre, og opfatte dem som
Størrelsers Overgange.\footnote[1]{„Da intet Tal opløftet til 0te Potens frembringer a, kan $\sqrt[0]{a}$ ikke angives ved nogen tidligere bekjendt Størrelse. Roden kan ikke være rational, fordi ethvert Potenstal af $0$te Grad er $1$, og den kan heller ikke være irrational, fordi intet Tal kan være indbefattet mellem to paa hinanden følgende Potenstal af $n$ Grad. Men $\sqrt[0]{a}$ kan bestemmes ved Anvendelse af Ligningen $\sqrt[q]{a^{p}}=a^{\frac{p}{q}}$, idet man sætter q=0, p=1, som giver $\sqrt[0]{a}=a^{\frac{1}{0}}=a^{\infty}$, som enten er $0$ eller $\infty$, efterfom $a<1$ eller $a>1$. Man seer heraf, at $0$ og $\infty$ henhøre hverken til de rationale eller til de irrationale Størrelser; tilmed ere de jo egentlig slet ikke Størrelser, men angive blot Størrelsens Ophør, dens Overgangsformer.“ A. Steen. Reen Mathem. Kbhvn 1856. S. 115—116.} Idet her da handles om forskjellige
Functioner af $0$, om $\sin 0$, $\cos 0$, $\mathrm{tg}\,0$ o.s.v., er
Meningen ikke, at $0$ i sig selv er Noget, heller ikke, at $0$ er slet Intet, thi $0$ er en Overgang, men er $0$ i $\mathrm{tg}\,0$ en Overgang, saa er $\infty$ i $\cot 0=\infty$ jo ogsaa en Overgang.

% --- p. 72 ---
\setcounter{footnote}{0}%
\opage{72}Da nu en Overgang, hvad Modsætningen af Positivt og
Negativt angaaer, kun faaer Betydning derved, at det
Positive gaaer over i det Negative, og omvendt det Negative
i det Positive, altsaa derved, at det Positive og det
Negative begge ere i Overgangen, og desaarsag høre med i
Overgangens Begreb: saa burde saadanne Ligninger, som: $\sin 0=0$, $\cot 0=\infty$, strengt taget, skrives: $\sin(\pm 0)=\pm 0$, $\cot(\pm 0)=\pm\infty$, o.s.v. Ved denne Adskillelse af $+$ og $-$ i $0$ og $\infty$ opkommer der i disse Grændsebetegnelser
selv en indre Forskjellighed, der er i høieste Grad
dialektisk. Betegnelsen $\pm 0$ udtrykker jo, naar vi see nærmere
til, at det ene Nul har skilt sig ad i tre forskjellige
Nuller, et $+0$, et $-0$, og det rene, neutrale $0$, den skarpe Negation af $+$ og $\div$; den samme Adskillelse gjentager
sig, hvad $\pm\infty$ angaaer; ogsaa her er et $+\infty$, et $-\infty$, og et neutralt $\infty$, hvori $+$ og $-$ ere forsvundne.
Naar det altsaa hedder, at $0$ og $\infty$ ere Overgangsbetegnelser, kan det ikke gjælde om det rene $0$ eller det neutrale $\infty$, men kun om $\pm 0$ og $\pm\infty$

Den Maade, hvorpaa de saakaldte ubestemte Former:
\[ \frac{0}{0}, 0.\infty, \frac{\infty}{\infty}, 0^{0}, \infty^{0} \]
og $1^{\infty}$, behandles i Mathematiken,
viser tydelig nok, hvad Overgangsformer ville sige. Disse
Former, som, idet de ere Functionsgrændser, altsaa ikke
umiddelbare Grændser, men i sig reflecterede Grændsebestemmelser,
kunne tænkes fremkomne derved, at $x$ i $f(x)$ sættes = $a$, og da bestemmes ifølge Grundligningen $f(a)=\lim f(a+h)$, dele sig jo væsentlig i to Grupper, af hvilke den
første: $\frac{0}{0}, 0.\infty, \frac{\infty}{\infty}$ og $\infty-\infty$ gaaer ind under Formlen:
\[ \frac{\Phi(a)}{\Psi(a)}=\frac{0}{0} \]
paa Betingelse af $\Phi(a)=0$, og $\Psi(a)=0$;
den anden: $0^{0}, \infty^{0}$ og $1^{\infty}$ paa samme Betingelse under Formlen $\Phi(a)^{\Psi(a)}$.

% --- p. 73 ---
\setcounter{footnote}{0}%
\opage{73}Hvorledes de under første Gruppe hørende Former
lade sig bestemme, sees af følgende Rækker:
\[ \frac{\Phi(a+h)}{\Psi(a+h)} \]
\[ =\frac{A_{0}h^{\alpha_{0}}+A_{1}h^{\alpha_{1}}+..}{B_{0}h^{\beta_{0}}+B_{1}h^{\beta_{1}}+..}=\frac{A_{0}}{B_{0}}h^{\alpha_{0}-\beta_{0}}\frac{1+\frac{A_{1}}{A_{0}}h^{\alpha_{1}-\alpha_{0}}+\ldots}{1+\frac{B_{1}}{B_{0}}h^{\beta_{1}-\beta_{0}}+\ldots}, \]
hvor Leddene, vel at mærke, ere ordnede efter stigende
Potenser af $h$, saaledes, at $\alpha_{0}$ er den mindste Exponent i øverste, $\beta_{0}$ den mindste Exponent i nederste Række. Eftersom
nu $\alpha_{0}-\beta_{0}\gtreqless 0$, faaer man
\[ \frac{\Phi(a)}{\Psi(a)}=\frac{0}{0}=\begin{cases}0\\ \frac{A_{0}}{B_{0}}\\ \infty\end{cases}. \]

Hvad den anden Gruppe angaaer, kommer det væsentlig
an paa at bestemme $\lim\Psi(a+h)\,l.\,\Phi(a+h)$, idet
\[ \Phi(a)^{\Psi(a)}=\lim\Phi(a+h)^{\Psi(a+h)}=e^{\lim\Psi(a+h)\,l\,\Phi(a+h)}. \]
udvikles
denne Exponent i Rækken:
\[ (B_{0}h^{\beta_{0}}+B_{1}h^{\beta_{1}}+\ldots)\,l.A_{0}h^{\alpha_{0}}+A_{1}h^{\alpha_{1}}+\ldots) \]
, vil det efter
behørig Transformation sees, at saavel $0^{0}$ som $\infty^{0}$, altid
% NOTE: no mark printed; the note stands at the foot of the page
\footnote[1]{Formerne: $\infty^{0}$ og $0^{0}$ bestemmes nemlig saaledes: $\Phi(x)^{\Psi(x)}=e^{\Psi(x)\,l\Phi(x)}$, thi $e^{l.\Phi(x)}=\Phi(x)$, $\Phi(a)^{\Psi(a)}=\lim\Phi(a+h)^{\Psi(a+h)}=\lim e^{\Psi(a+h)\,l(\Phi(a+h))}=e^{\lim\Psi(a+h)\,l.\Phi(a+h)}$; det kommer altsaa an paa at bestemme $\Psi(a+h)\,l\,\Phi(a+h)$. $\Psi(a+h)\,l\,\Phi(a+h)=(B_{0}h^{\beta_{0}}+B_{1}h^{\beta_{1}}+\ldots..)\,l.(A_{0}h^{\alpha_{0}}+A_{1}h^{\alpha_{1}}+\ldots..)=(B_{0}h^{\beta_{0}}+B_{1}h^{\beta_{1}}+\ldots).(l.A_{0}h^{\alpha_{0}}+l(1+\frac{A_{1}}{A_{0}}h^{\alpha_{1}-\alpha_{0}}+\ldots..)=(B_{0}h^{\beta_{0}}+B_{1}h^{\beta_{1}}+..)l.A_{0}\,(\mathrm{I})+\alpha_{0}(B_{0}h^{\beta_{0}}+B_{1}h^{\beta_{1}}+\ldots)l.h\,(\mathrm{II})+(B_{0}h^{\beta_{0}}+B_{1}h^{\beta_{1}}+\ldots)\left(\frac{\frac{A_{1}}{A_{0}}h^{\alpha_{1}-\alpha_{0}}+..}{1}-\ldots\right)(\mathrm{III})$. Da $\Psi(a)=0$, sees, at alle $\beta$ ere positive, følgelig lim (I) $=0$; da $\lim x^{n}\,l\,x=0$, er lim (II) $=0$; da endelig $\beta_{0}$ er positiv og $\alpha_{1}-\alpha_{0}$ positiv, idet Rækken er ordnet efter stigende Potenser (altsaa: $\alpha_{1}>\alpha_{0}$), saa kommer (III) til at indeholde lutter positive Exponenter af h, altsaa: lim (III) $=0$. Saaledes har man da: $\Phi(a)^{\Psi(a)}=e^{\lim\Psi(a+h)\,l(\Phi(a+h))}=e^{0}=1$. Exempel: $f(x)=(1+x)^{1+x^{2}}$ give, for $x=-1$, $0^{0}$; altsaa: $f(-1+h)=e^{(2h-h^{2})l.h}=e^{2h\,l\,h-h^{2}l.h}$, og $f(-1)=e^{0}=1$. Det kunde altsaa synes, som om $0^{0}$ og $\infty^{0}$ aldeles ikke vare ubestemte Former; dette vilde de heller ikke være, dersom $\Phi(a+h)$ og $\Psi(a+h)$ altid lode sig udvikle i Række efter stigende Potenser af h. Men der gives, som bekjendt, enkelte Tilfælde, hvor en Rækkeudvikling er umulig, og hvori man altsaa ved særegne Bestemmelser kan for $0^{0}$ eller for $\infty^{0}$ erholde en anden Værdi end 1. Saaledes giver $\left(\frac{1}{(\frac{1}{x})^{x}}\right)^{\frac{1}{x}}$, for $x=\infty$, $\infty^{0}$; men Lim $\left(\frac{1}{(\frac{1}{x})^{x}}\right)^{\frac{1}{x}}=\mathrm{Lim}\frac{1}{((\frac{1}{x})^{x})^{\frac{1}{x}}}=\frac{1}{\mathrm{Lim}\frac{1}{x}}=\infty$.}

% --- p. 74 ---
\setcounter{footnote}{0}%
\opage{74}have den bestemte Værdi 1, naar blot, ifølge Forudsætningen,
ovenstaaende Rækkeudvikling af $\Phi(a+h)$ og $\Psi(a+h)$ er mulig, og at $1^{\infty}$ maa, efterfom $\alpha_{1}+\beta_{0}\gtreqless 0$ give
\[ \Psi(a)\,l\,\Phi(a)=\begin{cases}0\\ A_{1}B_{0}\\ \infty\end{cases} ; \]
altsaa:
\[ \varphi(a)^{\Psi(a)}=1^{\infty}=\begin{cases}e^{0}=1\\ e^{A_{1}B_{0}}\\ e^{\infty}=\infty\end{cases}. \]\footnote[1]{Exempler: 1) $\left(1+\left(\frac{x}{a}\right)^{3}\right)^{\frac{a}{x}}=f(x)$ giver, for $x=0$, $1^{\infty}$; $f(0+h)=e^{\frac{a}{h}l(1+(\frac{h}{a})^{3})}=e^{\frac{a}{h}(\frac{h^{3}}{a^{3}}-\frac{h^{6}}{2a^{6}}....)}=e^{\frac{h^{2}}{a^{2}}-\frac{h^{5}}{2a^{5}}..}$, og $\lim f(0+h)=e^{0}=1$. 2) $f(x)=(2-x)^{\frac{1}{1-x}}$, som, for $x=1$, giver $1^{\infty}$. Sættes $1+h$ istedetfor $x$, haves $f(1+h)=(1-h)^{-\frac{1}{h}}=e^{-h^{-1}l.(1-h)}=e^{h^{-1}(\frac{h}{1}+\frac{h^{2}}{2}....)}=e^{1+\frac{h}{2}+\frac{h^{2}}{3}+....}$; $f(1)=\lim f(1+h)=e$. 3) $f(x)=(1+x)^{\frac{1}{x^{2}}}=1^{\infty}$, for $x=0$. $f(0+h)=e^{\frac{1}{h^{2}}l(1+h)}=e^{h^{-2}(\frac{h}{1}-\frac{h^{2}}{2}+\frac{h^{3}}{3}...)}=e^{h^{-1}-\frac{1}{2}+\frac{h}{3},..}$ og $\lim f(0+h)=f(0)=\infty$.}
% printed mark 1) stands after e^{A_1B_0} inside the cases; footnote set after the display

% the note 1) of p. 73 runs over onto the foot of this page: set in full at its mark

% --- p. 75 ---
\setcounter{footnote}{0}%
\opage{75}Da Functionsgrændsen, idet den viser det Endeliges
Overgang i det Uendelige, og i Forsvundetheden endnu
afspeiler det Forsvundne, ligesaa meget er i Vorden som i
Væren, spiller Tilvæxten $h$ den egentlige Hovedrolle: det
kommer i $f(a)=\lim f(a\pm h)=f(a\pm 0)$, saa at sige, an
paa Bevægelsen i $h$'s Gjennemgang gjennem $0$

Exempler kunne oplyse Sagen:
\[ f(x)=\left(\frac{1}{a-x}\right):\frac{1}{\sqrt{x^{2}-a^{2}}} \]

giver, for $a=x$, Formen $\frac{\infty}{\infty}$, og da saaledes $f(a+h)$
\[ =-\sqrt{\frac{2a}{h}+1}, \]
bliver $f(a)=\lim f(a+h)=\infty$, altsaa
i dette Tilfælde $\frac{\infty}{\infty}=\infty$. Haves derimod $f(x)=\frac{\mathrm{tg}\,x}{\sec x}$
saa bliver for $x=\frac{\pi}{2}$, $f(x)=\frac{\sin x}{\cos x}:\frac{1}{\cos x}=\sin x=1$, og følgelig $\frac{\infty}{\infty}=1$. Haves $f(x)=(a-x)\,l^{\frac{1}{a-x}}$, som for
\[ x=a, \]
giver Formen $0.\infty$, saa bliver $f(a+h)=-he^{-\frac{1}{h}}$
og $f(a)=\lim f(a+h)=0$; altsaa $0.\infty=0$, er derimod
$f(x)=\sin x.\cot x$, som, for $x=0$, ligeledes giver os Formen
\[ 0.\infty, \]
saa er $f(0)=\lim f(0+h)=\lim\sin(0+h)\frac{\cos(0+h)}{\sin(0+h)}$
\[ =\lim\cos h=1, \]
altsaa $0.\infty=1$. Er, for ogsaa at
anføre denne tredie Mulighed $f(x)=\frac{1}{(1-x)^{2}}(1-x)$,
som, for $x=1$, ligeledes giver Formen $0.\infty$, erholder
man ved Reduction
\[ \frac{(1-x)^{2}}{(1-x)^{2}(1-x)}=\frac{1}{1-x}, \]
og følgelig $0.\infty=\infty$

Her er altsaa i $f(a)=\lim f(a\pm h)=f(a\pm 0)$, hvil-
ken af de særskilte Former man end vælger, og hvorledes
end Betingelserne ændres, i den ene Grændse tre Grændser,
hvorpaa det væsentlig kommer an. Forsaavidt $h>0$, ligger

% --- p. 76 ---
\setcounter{footnote}{0}%
\opage{76}her imellem det endelige positive $h$ og $+0$, paa den ene,
og imellem det endelige negative $h$ og $-0$ paa den anden
Side en Størrelse, som saavel fra positiv som fra negativ
Side er i lige Grad endelig og uendelig, efterdi det er den
mindst mulige endelige Størrelse. Det Positive $+0$ betegner altsaa hverken $0$ eller en positiv Størrelse, men Grændsen
imellem den positive Størrelse og dens Negation, $0$; ligeledes betegner det Negative $-0$ Grændsen imellem den
negative Størrelse og dens Negation. Er nu t. Ex. i $\lim f(0\pm h)=\lim\cot(0\pm h)=f(0\pm 0)=\pm\infty$ paa
begge Sider af $0$ en dialektisk Overgangsstørrelse, en Størrelse,
der er uendelig-endelig, efterdi den som den mindst
mulige endelige Størrelse er uendelig lille; da er her ogsaa
paa begge Sider af $\infty$ en dialektisk Overgangsstørrelse, en
Størrelse, der er endelig-uendelig, efterdi den, som den størst
mulige endelige Størrelse, er uendelig stor. Paa denne Maade
bliver enhver Størrelsesovergang forvandlet til en dialektisk
Overgangsstørrelse: uendelig i sin mindst mulige Endelighed er
Overgangsstørrelsen et \emph{Differential}; uendelig i sin størst
mulige Endelighed er Overgangsstørrelsen, forsaavidt den
opfattes som en Differentialernes Sum, et \emph{Integral}.

\begin{center}

b) Differentialet: Momentets Dialektik.\end{center}

„Førend vi gaae videre, vil det være hensigtsmæssigt
nærmere at bestemme Begrebet om uendelig smaae Størrelser.
En \emph{uendelig lille eller forsvindende} Størrelse
er den, der er mindre end enhver nok saa liden Størrelse
af samme Slags. En saadan Størrelse er ikke Nul
(Intet), thi da var den ingen Størrelse, men i Sammenligning
med en endelig Størrelse bliver den at \emph{regne for
Nul}; thi ellers var den ikke mindre end enhver angivelig
Størrelse. Er altsaa $a$ en endelig Størrelse og $\alpha$ en uendelig lille, saa er $a\pm\alpha=a$, $a\alpha=0$

Bevægelse giver Exempel paa uendelig smaae Størrelser, idet
den tilbagelagte Vei hvert Dieblik voxer med en saadan
% --- p. 77 ---
\setcounter{footnote}{0}%
\opage{77}Størrelse. Forholdet imellem to uendelig smaae Størrelser $\alpha$ og $\beta$ kan være endeligt, og da siges de at være af samme
\emph{Orden}. Er $\alpha$ en uendelig lille Størrelse, saa vil $\alpha^{2}$ være
uendelig mindre og forsvindende i Sammenligning med $\alpha$. $\alpha^{2}$ og enhver Størrelse, hvis Forhold til denne er endeligt,
siges at være uendelig lille af anden Orden, idet man tænker
sig $\alpha$ at være af første Orden. Ligeledes siges $\alpha^{3}$ og
enhver Størrelse, hvis Forhold til denne er endeligt, at
være uendelig lille af tredie Orden, og saa fremdeles“.

„Er $a$ en endelig Størrelse, saa siges $\frac{a}{\alpha}$ og enhver Størrelse,
hvis Forhold til denne er endeligt, at være en uendelig
stor Størrelse, $\frac{a}{\alpha^{2}}$ at være uendelig stor af anden Orden
o. s. v. Det forstaaer sig, at man ogsaa kan tænke sig
saavel uendelig smaae som uendelig store Størrelser af Ordener,
hvis Exponenter ikke ere hele Tal. En uendelig
stor Størrelse er ikke det Samme som det, der angives ved
Tegnet $\frac{a}{0}$ eller $\infty$, hvilket man kunde kalde det absolut
Uendelige, den Grændse, der ligger udenfor al Storhed.
Man skriver imidlertid ofte $0$ og $\infty$, hvor der egentlig
menes en uendelig lille og en uendelig stor Størrelse.“\footnote[1]{Chr. Jürgensen. Begyndelsesgruudene af den mathematiske Analyse og dens Anvendelser. Kbhvn 1855. S. 41—43.}

For at overbevise sig om, at her ved Antagelse af
uendelige Størrelser ikke indsmugles uholdbare Begreber,
lægge man Mærke til, at enhver Storhedsbestemmelse væsentlig
er en Forholdsbestemmelse, og at en Forholdsbestemmelse
kan være constant, uagtet Leddene tiltage eller aftage
i det Uendelige. Kun i Læren om de uendelige Størrelser
kan Størrelsesbegrebet fuldstændig udvikles, thi kun i denne
Lære bliver det tilstrækkeligt klart, at Størrelsen er Forhold
og Forholdet Størrelse. I Differentialcoifficienten gaaer
Størrelsen op i det betegnende Forhold som Forholds-

% --- p. 78 ---
\setcounter{footnote}{0}%
\opage{78}størrelse; Differentialcoifficienternes forskjellige Orden beroer % PRINTED AS IS: Differentialcoifficienternes (for Differentialcoefficienternes)
paa et System af Forholdsstørrelser, og Differentiationsprinciperne
angive den ved de eiendommelige Forholdsconseqvenser
bestemte Methode.

\begin{center}α)

\emph{Differential- og Differentialcoefficient: Forholdsstørrelse.}\end{center}

En almindelig Storhed, en almindelig Lidenhed er
endnu ingen Størrelse; først, naar Storheden og Lidenheden
nærmere bestemmes, fremkommer Størrelsen. Da nu
umiddelbar Bestemthed falder sammen med Begrændsning,
og det umiddelbart Begrændsede er et Endeligt, saa ligger
den Opfattelse nær, at kun den begrændsede, det vil sige,
endelige Størrelse er i egentlig Forstand Størrelse. Som
endelig Størrelse kan Størrelsen fastholdes i Forestillingen
og bestemmes for Anskuelsen; ligesom, hvad benævnte Tal
angaaer, en Alen kan bestemmes for Anskuelsen deels umiddelbart,
deels ved Sammenligning med et Qvarteer, saaledes kunne Talrækkens endelige Tal opfattes deels umiddelbart,
deels ved Sammenligning med hverandre indbyrdes;
hvert Tal har da i denne sin Endelighed et eget Navn, et % PRINTED AS IS: et (for en) before "egen Betegnelse"
egen Betegnelse, og kan forsaavidt optræde som et for sig
tilværende Hele. Idet nu de saaledes opfattede Størrelser
indgaae i Forbindelser med hverandre, viser det sig, at, jo
mere den endelige Forestilling har Overvægt over Tænkningen,
desto mere Eftertryk falder der paa de givne Tal,
de tilværende Forholdsled; først efterhaanden som den begribende
Tænkning faaer Magt med den umiddelbare Forestilling,
forvandles Leddene til Momenter, saa at Eftertrykket
tilsidst kommer til at hvile, ikke paa de umiddelbare
Led, men paa Forholdet. I Udtryk som: $1$, $\frac{2}{2}$, $\frac{a}{a}$, $\mathrm{Lim}\frac{x}{x}$
og $\lim\frac{x}{x}$ bliver Eenheden at see fra væsentlig forskjellige
Synspunkter.

% --- p. 79 ---
\setcounter{footnote}{0}%
\opage{79}Saalænge Forestillingen uvilkaarlig hæfter ved Tegnet: $1$, tager det sig ud, som var Eenheden et umiddelbart
Noget, et Tilværende, man kunde anskue og ligefrem paapege;
denne Forestillingens Umiddelbarhed forsvinder, naar
man begynder at opfatte Eenheden som et Forhold: $\frac{2}{2}=\frac{3}{3}$
\[ =\frac{4}{4}\ldots \]
Med Forholdserkjendelsen kommer
Indsigten, men
Leddene have dog endnu særlig Betydning ligeoverfor det
Forhold, hvis Led de ere: Forholdserkjendelsen er endnu
overfladisk. Ved Betegnelsen $\frac{a}{a}$ skeer et betydeligt Fremskridt,
idet det almindelige Udtryk afløser de particulære,
og det bringes til Bevidsthed, at Størrelser ikke blot kunne
indgaae i forskjellige Relationer, men at Relationen nødvendig
hører med i Størrelsens Væsen. De algebraiske
Betegnelser ere imidlertid behæftede med den Biforestilling,
at enhver nok saa vilkaarlig Størrelse maa være endelig.
Ligningerne: $\frac{a}{b}=3$, $\frac{b}{c}=4$, $\frac{c}{d}=e$ have jo vistnok deres
Betydning i Forholdene af $a$ og $b$, $b$ og $c$, $c$ og $d$; men
den Indbildning, at disse bestemte Forhold ville opløse sig
i Ubestemthed, naar $a$, $b$, $c$ og $d$ ophørte at være endelige
Størrelser, er endnu ikke hævet. Denne Indbildning forsvinder:
naar det sees, hvad der ligger i Betegnelsen:
\[ \mathrm{Lim}\frac{x}{x}=\lim\frac{x}{x}=\frac{\infty}{\infty}=1, \]
thi hvorledes man end forklarer sig $\frac{\infty}{\infty}$ og $\frac{0}{0}$, staaer dette dog fast, at Eenheden her bliver
opfattet som det rene Forholdsudtryk. At det ved et saadant Forholdsudtryk aldeles ikke kommer paa de endelige
Led, men ene og alene paa det ved Leddene bestemte Forhold,
kan neppe tilkjendegives paa en mere eftertrykkelig
Maade, end derved, at det bestemte Forhold bliver staaende
uforandret, efterat $\infty$ og $0$ have indtaget Leddenes Plads.

% --- p. 80 ---
\setcounter{footnote}{0}%
\opage{80}At Størrelsesbegrebet maa fuldendes i Functionsbegrebet,
indlyser deraf, at Functionsloven, som ovenfor angivet,
paa een Gang omfatter det Endelige og det Uendelige.
Af $f(x)=y$ og $f(x+h)=y+k$ følger $\frac{f(x+h)-f(x)}{h}$
\[ =\frac{k}{h}, \]
og lim $\frac{f(x+h)-f(x)}{h}=\lim\frac{k}{h}=\frac{dy}{dx}$.
I $f(x+h)-f(x)=k$ er $h$ en Størrelse for sig, $k$ ligeledes; $h$ og $k$ have i endelig Sammenligning umiddelbar Gyldighed; en
saadan Gyldighed kan i $f(x+dx)=y+dy$ derimod ikke tillægges $dx$ og $dy$, hvorfor den exacte Analyse da ogsaa
idelig sætter sig imod, at de uden videre opfattes som
Størrelser; i reen Umiddelbarhed ere de da heller ikke
Størrelser: som Forholdsstørrelser ere de dialektiske. Overseer
man det Dialektiske, kan man beskrive Differentialet,
hvordan man vil; man træffer aldrig Sagen. Differentialet $dy$ er, ifølge $f(x+dx)-f(x)$, en uendelig lille Differens;
men den uendelig lille Differens har i sig den Modsigelse
at være som $0$ og dog større end $0$. Opfattes Differentialet
ikke som Størrelse, men som Størrelsesanlæg, maa
et Anlæg til Størrelse jo dog i Principet selv være Størrelse.
Bestemmer man Differentialet ikke som en vilkaarlig
lille, men som en uendelig forsvindende Størrelse, hvorledes
kan man da i $f(x+dx)=y+dy$ sætte $dx=dx$, medens $dy$ varierer? Hvorledes kan en bestandig vordende, en
uophørlig forsvindende d. v. s. ikke værende Størrelse være
constant? Anerkjender man derimod det Dialektiske, vil
Differentialet, som uendelig lille Differens, som Størrelseselement,
som Forholdsmoment, fra alle Sider belyse Størrelsesbegrebet.
Haves $f(x+dx)=4(x+dx)=4x+4dx=y+dy$, saa er unegtelig $dy=4dx$, og forsaavidt
synes vi i $dy$ og $dx$ at have umiddelbare Størrelser; men idet Eftertrykket falder, ikke paa $dy$ og $dx$ for sig, som om
vi ved $dx+dx+dx+dx$ kunde erholde $dy$, men paa $\frac{dy}{dx}$

% --- p. 81 ---
\setcounter{footnote}{0}%
\opage{81}altsaa paa Forholdet; bive Differentialerne herved udtrykkelig % PRINTED AS IS: bive (for blive)
bestemte som Forholdsmomenter. I Ligningen $dy=4dx$ er $4$ den Coefficient, hvormed $dx$ skal multipliceres for at
give $dy$, men som umiddelbar Størrelse er Coefficienten
endnu ikke svarende til sit Begreb: Differentialcoefficient.
For at erkjendes som Differentialcoefficent maa $4$ fra en % PRINTED AS IS: Differentialcoefficent (for Differentialcoefficient)
umiddelbar Størrelse forvandles til et i sig reflecteret For-
holdsudtryk: $\frac{dy}{dx}$. At $dy=\frac{dy}{dx}dx$, maa ikke forstaaes tau-
tologisk, som om her stod $a=\frac{a}{b}.b$; thi $\frac{dy}{dx}=4$ betegner
udtrykkelig en til et Forhold forvandlet Størrelse, en Forholdsstørrelse:
\emph{Differentialcoefficenten er væsentlig Qvotient.} Differentialcoefficienten, som i det exacte % PRINTED AS IS: Differentialcoefficenten (for Differentialcoefficienten)
Sprog er Grændsen for Forholdet imellem de sammenhørende
og samtidig forsvindende Tilvæxter til den afhængige
og uafhængige Variable, giver altsaa Differentialet sin
Betydning, ikke omvendt.

Dialektikens anskuelige Billede er Bevægelsen, navnlig
Rumsbevægelsen. Newton, som for at løse phoronomiske
Opgaver, tager Geometrien til Hjælp, fremstiller den
ligelig voxende Tid som Abscisse, det i Tiden gjennemløbne
Rum som Ordinat. Curvens Coordinater, der
forestille de stedse foranderlige, af- og tiltagende Størrelser,
Tiden og Rummet, kalder han „Fluenterne“, og
betegner dem ved $x$, $y$ o. s. v., Hastighederne derimod, hvormed
de af- eller tiltage, kalder han „Fluxioner“, og betegner
dem ved $\dot{x}$, $\dot{y}$ o. s. v. Det forstaaer sig af sig selv, at Abscissens
Fluxion $\dot{x}$, da den fremstiller Hastigheden for Tiden,
som den jo repræsenterer, og Tiden tænkes at voxe ligelig
(»\textit{uniformiter fluens}»), maa betragtes som constant. Det
følger da heraf, at den anden, den tredie og alle følgende
Fluxioner af $x$ ere lige med Nul, endvidere, at Fluxionen af $x$ er det Maal, hvormed Fluxionen af $y$ maales, og
altsaa kan antages for Eenhed. Idet nu Tid og Rum og

% --- p. 82 ---
\setcounter{footnote}{0}%
\opage{82}følgelig de Linier, hvorved de repræsenteres, $x$ og $y$, have i Øieblikket, d. e. i Tiden $0$, den ene Hastigheden $\dot{x}$, den anden $\dot{y}$, vil disse Liniers Til- og Aftagen, da denne jo
er ligefrem proportional med Hastigheden, kunne udtrykkes
ved $\dot{x}0$ og $\dot{y}0$. Denne med Hastigheden proportionale Til- eller
Aftagen af Fluenterne, der maa være uendelig lille, da
Hastigheden jo kun med Hensyn paa det enkelte Dieblik
kan ansees for constant, betegner Newton ved Udtrykket:
\emph{Moment.}\footnote[1]{Jvfr. Weissenborn. Die Principien der höheren Analysis. S. 21—57.}

Der Gedanke kann nicht richtiger bestimmt werden,
als Newton ihn gegeben hat.\footnote[2]{Ogsaa her have vi et Exempel, som godtgjør, hvorledes et Begreb, der, saalænge det holdes i logisk Alminbelighed, viser sig som aldeles tilfredsstillende, dog alligevel, naar det gjælder en articulerende Discretion, en Udviling af punktuelle Bestemmelser, kan være mangelfuldt og utilfredsstillende. „Man hat öfter die Fluxionsmethode der jetzt vielfach gebräuchlichen Grenzmethode nahe gestellt, und dies grossentheils mit gutem Rechte, indem Newton stets auf Verhältniss der Fluxionen $\dot{y}$ und $\dot{x}$ kommt. Auf diese Weise umgeht er die Schwierigkeit, die den damaligen Mathematikern ein unüberwindliches Hinderniss schien, nämlich das Rechnen mit unendlich kleinen Grössen oder Null. Denn so lange die Incremente endlich sind und eine angebbare Grösse besitzen, hat ihre Anwendung im Calcul kein Bedenken, Newton lässt sie auch bis zu Null abrechnen, und nennt sie denn Momente oder auch („\textit{Princ. phil. nat.}« \textit{Lib. II. Sect II. Lemma II. pag. 251} der ersten Ausgabe und „\textit{Tractatus de quadratura curvarum.}“ \textit{Prop. I. Problem I. Demonstratio.)} »\textit{incrementa momentanea}“ oder „\textit{momenta nascentia sive evanescentia}“, er hütet sich aber wohl, mit diesen zu rechnen, sondern lässt im Zustande des Verschwindens die Fluxionen, endliche und bestimmte Grössen, gewissermassen als Reserve an die Stelle der zum weiteren Dienste untauglich gewordenen Momente einrücken, auf diese Weise das Rechnen mit unendlich kleinen Grössen, welche Null sein und doch noch einen von Null verschiedenen Werth besitzen sollen, auf das Glücklichste vermeiden; und wenn er (»Tractatus de quadratura curvarum.“ Introductio) sagt: „Volui ostendere, quod in methodo fluxionum non opus sit figuras infinite parvas in Geometriam introducere«, und ebendaselbst: Quantitates mathematicas non ut ex partibus quam minimis constantes, sed ut motu continuo descriptas hic desidero«, so muß man unbedingt zugeben, daß er nicht mehr verspricht als er in der That gehalten hat. Unstreitig würde der große Vortheil, den seine Methode durch die Vermeidung des zu Widersprüchen führenden Begriffs des Uendlichkleinen, wodurch zugleich der unnatürlichen Auffassung der Curven als Polygone vorgebeugt wurde, über die Differenzialmethode, wie sie lange Zeit hindurch dargestellt ward, gewonnen hatte, — unstreitig würde dieser große Vortheil ihr schon frühe den Vorrang vor der Leibnizischen Theorie verschafft haben, wenn ihn nicht Newton dadurch erkauft hätte, daß er die Rechnung aus dem Gebiete der Geometrie weg in das der Phoronomie hinüberspielte. Dazu kam aber noch als zweites sehr wichtiges Moment, daß die Fluxionsrechnung Newtons eines hinlänglich passenden und ausgebildeten Algorithmus entbehrte, was sich besonders beim Uebergehen von den Fluxionen zu den Fluenten als ein Nachtheil sehr bemerklich macht. Denn es war ja hauptsächlich die Auffindung eines solchen Algorithmus, um die sich damals Alles drehte. So kann es denn nicht wunderbar sein, wenn die Leibnizische Differenzialrechnung, welche alle Forderungen, die in dieser Beziehung gestellt werden können, auf das Vollkommenste befriedigt, der Fluxionsrechnung trotz ihrer sichereren Begründung sehr bald den Rang ablief.“ Weissenborn. Anfrt. Skr. S. 57—58.} Ich trenne dabei die

% --- p. 83 ---
\setcounter{footnote}{0}%
\opage{83}Bestimmungen ab, die der Vorstellung der Bewegung und
der Geschwindigkeit angehören, (von welcher er vornehmlich
den Namen Fluxionen nahm), weil der Gedanke hierin
nicht in der gehörigen Abstraction, sondern konkret, vermischt
mit ausserwesentlichen Formen erscheint. Diese Fluxionen
erklärt Newton (\textit{Princ. mathem. phil. nat. L. I. Lemma XI.
Schol.}) dahin, daß er nicht \emph{untheilbare} --- eine Form,
deren sich frühere Mathematiker, Cavalleri und andere, bedienten,
und welche den Begriff eines an sich bestimmten
Quantums enthält, verstehe, sondern verschwindende Theilbare.
Ferner nicht Summen und Verhältnisse bestimmter
% the note 2) of p. 82 runs over onto the foot of this page: set in full at its mark
% --- p. 84 ---
\setcounter{footnote}{0}%
\opage{84}Theile, sondern die \emph{Grenzen} (limites) der \emph{Summen} und
Verhältnisse. Es werde die Einwendung gemacht, daß verschwindende
Grössen kein leztes Verhältniß haben, weil
es, ehe sie verschwunden, nicht das Lezte, und wenn sie verschwunden,
keines mehr ist. Aber unter dem Verhältnisse
verschwindender Grössen sey das Verhältniß zu verstehen,
nicht ehe sie verschwinden, und nicht \emph{nachher}, sondern
mit dem sie verschwinden (quacum evanescunt). Ebenso
ist das erste Verhältniß werdender Grössen, das mit dem
sie werden.“\footnote[1]{Hegel. Wissensch. der Logik. 1ster Th. Berlin 1833. S. 302—303.}

\begin{center}β) Differentialcoefficienternes forskjellige Orden: System af Forholdsstørrelser.\end{center}

Uagtet det, væsentlig seet, er Differentialcoefficienten,
der giver den Afhængiges og Uafhængiges Differentialer
deres Betydning, maa den dog som et Differentialernes
Forhold nødvendig forudsætte Differentialerne selv, og kommer
forsaavidt til at beroe paa disses eiendommelige Gyldighed.
I mathematisk Praxis kan det maaskee have sin
Betydning, at man holder sig exact til det i Grændsen Negative,
og sætter Differentialcoefficienten $\frac{dy}{dx}=\frac{0}{0}$; men herved
er den paa Indseen og Begriben arbeidende Tænkning
ingenlunde tilfredsstillet. I Udtrykket $\frac{0}{0}=4$ kan $\frac{0}{0}$ umulig gaae paa et Nærværende, men kun paa et Forbigangent
og et Tilkommende. Skulde man virkelig antage $\frac{0}{0}=4$, maatte man antage et Forhold uden Forholdsled, et Forhold
af Intet med Intet; $\frac{0}{0}$ er altsaa ikke $4$, betegner heller ikke, at $4$ endnu er, efter at det tilsvarende Forholds-udtryk er forsvundet, men tilkjendegiver kun, at $4$ har været,

% --- p. 85 ---
\setcounter{footnote}{0}%
\opage{85}\begin{center}0\end{center}

idet her har været et nu forsvundet Forhold. Skal $\frac{0}{0}$, strengt
taget, betegne Differentialcoefficienten, saa er Differentialcoefficienten
ingen Størrelse; skal Eftertrykket falde paa $4$, saa er $4$ som umiddelbar Størrelse ingen Differentialcoefficient.
Mathematiken sætter da heller ikke uden videre $\frac{0}{0}=4$; ved Betegnelsen for det Forsvundne mindes den endnu det
Forsvundne, erindrer sig, hvad det egentlig var, der forsvandt, og hvorledes det forsvandt: $\lim\frac{4(x+h)-4x}{h}=\frac{0}{0}=4$. Indsigten er her i strengeste Forstand en Erindring;
i Erindringen sees det Forbigangne som et Nærværende:
det Negerede reflecterer sin Væren i Negationen. Idet $\frac{dy}{dx}$ sættes istedetfor $\frac{0}{0}$, kunne de forsvindende Led udgjøre et
endnu værende Forhold, en Forholdsstørrelse af bestemt
Værdi.

Som bestemt Forholdsstørrelse kan Differentialcoefficienten
være enten endelig eller uendelig, enten constant eller
variabel. Saaledes er i $f(x+dx)=a(x+dx)=ax+adx=y+dy$ Differentialcoefficienten $\frac{dy}{dx}=a$ constant; i $f(x+dx)=(x+dx)^{2}=x^{2}+2xdx+dx^{2}=y+dy$ er $\frac{dy}{dx}=2x$ derimod variabel. Som variabel Størrelse kan
Differentialcoefficienten $f'(x)=\frac{dy}{dx}=2x$ igjen differentieres;
hvorimod $f'(x+dx)=2(x+dx)=f'(x)+f''(x)dx+\ldots=y_{1}+dy_{1}$ giver en Differentialcoefficient: $\frac{dy_{1}}{dx}=2$, der, som con-

stant Størrelse, ikke kan differentieres. Da Differentialcoefficienten allerede som saadan er en Forholdsstørrelse, er den
differentierede Differentialcoefficient som en ved en Forholds%

% --- p. 86 ---
\setcounter{footnote}{0}%
\opage{86}størrelse bestemt Forholdsstørrelse $\frac{dy_{1}}{dx}=\frac{d\frac{dy}{dx}}{dx}$, en Forholds-
størrelse af anden Orden $\frac{d^{2}y}{dx^{2}}$; paa denne Maade kunne
Forholdsstørrelserne udvikle sig i et System af lavere og
høiere Ordener. Af $f(x)=y=x^{n}$ t. Ex. faaes:
\[ f(x+h)=f(x)+f'(x)\frac{h}{1}+f''(x)\frac{h^{2}}{1.2}+f'''(x)\frac{h^{3}}{1.2.3}+\ldots \]
\[ (x+h)^{n}=x^{n}+nx^{n-1}\frac{h}{1}+n(n-1)x^{n-2}\frac{h^{2}}{1.2}+n(n-1)(n-2)x^{n-3}\frac{h^{3}}{1.2.3}+\ldots \]
\[ =y+\frac{dy}{dx}\frac{h}{1}+\frac{d^{2}y}{dx^{2}}\frac{h^{2}}{1.2}+\frac{d^{3}y}{dx^{3}}\frac{h^{3}}{1.2.3}+\ldots \]
Hvilken indre Uendelighed et saadant System af Forholdsstørrelser
kan omfatte, sees, naar vi vælge en Function af
flere indbyrdes uafhængige Variable, t. Ex. $f(x_{1},x_{2},\ldots x_{t})=u$. I den til $f(x_{1}+h_{1},x_{2}+h_{2}\ldots x_{t}+h_{t})\sim e^{\frac{du}{dx_{1}}h_{1}+\frac{du}{dx_{2}}h_{2}+\ldots+\frac{du}{dx_{t}}h_{t}}$ svarende Rækkeudvikling, hvor
$\frac{1}{[r_{1}+r_{2}+\ldots r_{t}]}\left(\frac{du}{dx_{1}}h_{1}+\frac{du}{dx_{2}}+\ldots\frac{du}{dx_{t}}h_{t}\right)^{(r_{1}+r_{2}+\ldots r_{t})}$ er det
almindelige Led, er, idet de til Symboliken svarende
Ændringer, hvorved Potensexponenter forvandles til Diffe-
rentiationsindices, ere foretagne, $\frac{d^{r_{1}+r_{2}\ldots r_{t}}u}{dx_{1}^{r_{1}}dx_{2}^{r_{2}}\ldots dx_{t}^{r_{t}}}\frac{h_{1}^{r_{1}}}{[r_{1}]}\cdot\frac{h_{2}^{r_{2}}}{[r_{2}]}\ldots\frac{h_{t}^{r_{t}}}{[r_{t}]}$ igjen dette almindelige Leds almindelige Led, hvori
Differentialcoefficienten altsaa kan fremstille Forhold af Forhold
i det uendelige.\footnote[1]{At den indre Uendelighed kan fra uendelig mange Sider ligeledes opgaae for Anskuelsen, viser den analytiske Geometrie. Skal, for at vælge et Exempel af Fladers Generation, en Flade danne sig ved Bevægelsen af en vis Linie $\{F_{1}(x,y,z,\alpha,\beta,\gamma\ldots..)=0, F_{2}(x,y,z,\alpha,\beta,\gamma\ldots..)=0$ hvori $\alpha,\beta,\gamma\ldots.$ ere Constanter, der dog variere fra den ene Stilling af Linien til den anden: maae, for at Loven for denne Variation kan være bestemt, visse Betingelser forudsættes, f. Ex. at Frembringeren (Generatrix) under sin Bevægelse skal skjære et Antal Ledelinier. Antage vi, at der i ovenstaaende Ligninger indgaae n Constanter, maae Betingelserne gjøre det muligt at faae disse n Constanter eliminerede, thi den Flade, der skal dannes, er Indbegrebet af samtlige Frembringerstillinger, og kan altsaa ikke afhænge af nogen enkelt især, d. v. s. Fladens Ligning maa være uafhængig af $\alpha,\beta,\gamma\ldots.$ Vi maae altsaa i det Hele kunne danne $(n+1)$ Ligninger mellem $x,y,z,\alpha,\beta,\gamma\ldots.$ For nu ret at tydeliggjøre, hvilken Uendelighed af Flader der kan dannes af en given Frembringer, ville vi for et Øieblik tænke os, at vi for de n første af de $(n+1)$ Ligninger have en aldeles bestemt Functionsform, medens den $(n+1)$te Lignings Form er ubestemt; her er da for denne $(n+1)$te Ligning en Uendelighed af Former mulig; for hver enkelt Form, vi vælge, faae vi en Flade, altsaa en Uendelighed af Flader. Betænke vi derhos, at vi jo kunne lade hver enkelt af de $(n+1)$ Ligninger skiftevis have en ubestemt Form, da bliver det klart, hvilken Uendelighed af Flader det ene System af $(n+1)$ Ligninger ved Eliminationen af de n Constanter kan give Anledning til. I denne uendelige Vrimmel er fremdeles en systematisk Deling mulig. For at foretage en saadan Deling, kunne vi særlig tage Hensyn til 1) Frembringeren og 2) Loven for Frembringelsen. Alle de Flader, som have samme Frembringer, kunne da henføres til samme Klasse; hver enkelt Klasse kan endvidere deles i Familier, der alle have Klassens Særkjende, nemlig den samme Art Curver til Frembringere; Familiens Særkjende sees af den Lov, hvorefter Generatrix bevæger sig. Indenfor hver Familie fremkommer derimod forskjellige Arter, idet den samme Lov optræder i forskjellige Modificationer, og indenfor samme Art deforskjellige Individer, idet bestemte Størrelser indsættes i den særskilte Form. Cylindrene udgjøre saaledes for sig en stor Familie, der hører under Classen af retlinede Flader; Familiens Særkjende er, at alle Frembringerstillinger ere parallele. Frembringeren kan nemlig under sin Be- vægelse vedblive at være parallel med sig selv, og dog frembringe meget forskjellige Arter af Cylindre: Omdreiningscylindre, hyperbolske Cylindre o. s. v. Omdreiningscylindre t. Ex. afgive en Uendelighed af Individer, thi for hver forskjellig Radius, vi give den paa Frembringeren perpendiculære Cirkel, faaes et forskjelligt Individ. (Kategorierne: Classe, Familie, Art, Individ anføres her efter en Forelæsning af Hr. Prof. A. Steen).}

% --- p. 87 ---
\setcounter{footnote}{0}%
\opage{87}Hvorledes Differentialcoefficienter af forskjellig Orden
i deres Anvendelse faae en forskjellig Betydning, kan sees
% the note 1) of p. 86 runs over onto the foot of this page: set in full at its mark
% spaced words Klasse, Familier, Arter, Individer stand inside footnote 1 (p. 86), left unchanged

% --- p. 88 ---
\setcounter{footnote}{0}%
\opage{88}af Exempler saavel fra Geometrien som fra Mechaniken.
Vælge vi fra Geometrien Røringsproblemet, og lade $f(x)=y$, $\Phi(\xi)=\eta$ være Ligningerne for to Curver, som
have Punktet $(x,y)$ fælleds, saa bliver Differensen af de til Punktet $(x,y)$ svarende consecutive Ordinater at udtrykke ved:

\[ f(x)-\Phi(x)+(f'(x)-\Phi'(x))\frac{h}{1}+\ldots(f^{n}(x)-\Phi^{n}(x))\frac{h^{n}}{[n]}\ldots, \]

\[ y-\eta+\left(\frac{dy}{dx}-\frac{d\eta}{d\xi}\right)\frac{h}{1}+\ldots\left(\frac{d^{n}y}{dx^{n}}-\frac{d^{n}\eta}{d\xi^{n}}\right)\frac{h^{n}}{[n]}\ldots, \]

hvori $y=\eta, x=\xi$. Størrelsen af denne Differens afhænger,
idet $h$ har samme uendelig lille Værdi, af hvilket
Led der er det første, som ikke er $0$, d. v. s., af, hvilken
Orden Røringen er; er f. Ex. $y=\eta, \frac{dy}{dx}=\frac{d\eta}{d\xi}\ldots.\frac{d^{n}y}{dx^{n}}=\frac{d^{n}\eta}{d\xi^{n}}$, men $\frac{d^{n+1}y}{dx^{n+1}}\gtrless\frac{d^{n+1}\eta}{d\xi^{n+1}}$,
saa at Røringen
er af $n$te Orden; da er $\alpha)$ Differensen desto mindre, jo større $n$ er, og $\beta)$ skifter Fortegn med $h$, naar $n$ er lige, men ikke, naar
$n$ er ulige. Til $(\alpha)$ svarer en desto inderligere Røring, d. e. de to Curver have desto flere Elementer tilfælleds, jo større $n$ er; ved $(\beta)$ afgjøres det, hvorvidt Røringen er forbunden
med Skjæring, eller ikke. I Mechaniken afgiver Udtrykket
for Hastigheden: $v=\frac{ds}{dt}$, og for Kraften: $\Phi=\frac{d^{2}s}{dt^{2}}$ letfattelige Exempler.
% the note 1) of p. 86 runs over onto the foot of this page: set in full at its mark

% --- p. 89 ---
\setcounter{footnote}{0}%
\opage{89}Af Differentialcoefficienternes Anvendelse sees endvidere,
at Differentialerne selv umulig kunne være rene Nuller,
men maae som bestemte Forholdsmomenter være Størrelsesanlæg,
og følgelig i Principet selv Størrelser. Saaledes
viser Betingelsen før

Røringen „af nte Orden, at de to
Curver have følgende Punkter tilfælleds:

\begin{quote}
\emph{Abscisser:} \emph{Ordinater:}
\end{quote}
\[ x \quad y, (\eta) \]
\[ x+dx \quad y+dy, (\eta+d\eta) \]
\[ x+2dx \quad y+2dy+d^{2}y, (\eta+d\eta+d^{2}\eta) \]
\[ x+n dx \quad y+n dy+\frac{n(n-1)}{1.2}d^{2}y\ldots d^{n}y, \]
\[ (\eta+n d\eta+\ldots d^{n}\eta) \]

Røringen af $n$te Orden vil altsaa med andre Ord sige, at de to Curver have $n$ consecutive Elementer tilfælleds.
Mathematiken gjør selv Forskjel paa Røring af nte og Røring
af 1ste Orden, og lærer da med det Samme, at det
at have $n$ Elementer tilfælleds, er, for $n>1$, mere end at have eet Element fælleds, eller i Tegn: $n dx>dx$; men naar $n dx>dx$, uagtet $n$ er et endeligt Tal, kan $dx$ umulig være $0$.\footnote[1]{I den deskriptive Geometrie, for ogsaa at see Sagen fra denne Side, bevises ved Anvendelse af Keglesnittenes projective Egenskaber, at, naar en Sexkant indskrives i et Keglesnit, da ligge de 3 Punkter, hvori denne Sexkants modstaaende Sider skjære hinanden, to og to, i een og samme rette Linie. En Femkant er da i Consequens heraf at betragte som en Sexkant med \emph{een uendelig lille Side}; denne uendelig lille Side er imidlertid mere end et Punkt, thi den angiver Retningen af Tangenten til Keglesnittet (og et Punkt kan ikke angive nogen Retning; derfor ville ogsaa de 3 Punkter, hvori Femkantens modstaaende Sider samt Tangenten skjære hinanden, to og to, ligeledes ligge i een og samme rette Linie.}

At Differentialet som saadant ikke umiddelbart kan være $0$, sees, om muligt, endnu tydeligere, naar vi, hvad Be-

% --- p. 90 ---
\setcounter{footnote}{0}%
\opage{90}vægelsen angaaer, betragte Forholdet af Rummet og Tiden.
Medens nemlig $\frac{ds}{dt}=v$ betegner et Forhold af Rum og
Tid og derved en i sit Anlæg bestemt Hastighed, kan $\frac{0}{0}$, der
jo udtrykker en absolut Fornegtelse, saavel af Tiden som af
Rummet, og altsaa falder i den rene Abstraction udenfor
Tid og Rum, umulig angive Hastighedsanlæg for en Bevægelse,
der skal foregaae i Tid og Rum. Ved eensidig at
fastholde de abstracte Beregningsudtryk seer man feil af
Sagens Væsen, og hilder sig i en Modsigelse. I Ligningen $f(t+dt)=s+ds$ er $dt$ constant, medens $ds$ kan variere, efter som $t$ variere,
eftersom t varierer; hvor dialektisk denne constante Størrelse
imidlertid er, indlyser blandt Andet deraf, at man for
en varierende Hastighed sætter $v=\frac{ds}{dt}$, som om $v$ i Tidselementet $dt$ var constant. Er nu $dt=0$, saa er her hverken Hastighed eller Hastighedsanlæg; er $dt$ derimod ikke $0$, men et virkeligt Tidselement, saa vedbliver Hastigheden
endog i det valgte Tidselement at variere; her er da en
Varieren, som, idet den svarer til Anlæg og Mulighed, kun
ved at svæve imellem Noget og Intet, og saaledes holde sig
dialektisk, kan svare til sit Begreb.

\begin{center}γ) \emph{Differentiationsprinciper: den ved Forholdsconseqvenserne} bestemte Methode.\end{center}

Hvilke forskjellige Indtryk den af Leibniz indførte Differentialregning
gjorde paa hans Samtidige, sees blandt
Andet deraf, at medens Mathematikere som Jacob og Johan
Bernouilli anvendte den med Held paa de vanskeligste % PRINTED AS IS: Bernouilli (Rettelser: Bernoulli)
Problemer, vilde Leibnitzes Ven Huygens ikke forlade de
Gamles Methode, og Hollænderen Nieuwentiit optraadte
endog polemisk, idet han indvendte: 1) at Leibnitz ligesaalidt
som Barrow og Newton, hos hvem allerede de uendelige
Størrelser forekom, kunde forklare, hvorledes disse
Størrelser adskilte sig fra Intet eller Nul; 2) at det ikke
% --- p. 91 ---
\setcounter{footnote}{0}%
\opage{91}var klart, hvorledes Differentialerne af høiere Orden lode
sig adskille fra dem af første Orden, og $3)$ at Differentialmethoden
ikke lod sig anvende paa Exponentialfunctionen.\footnote[1]{Weissenborn anfr. Skr. S. 98.}
Den Vaklen og Usikkerhed, der gaaer igjennem Leibnitzes
Forklaringer og Forsvar, viser tydelig nok, at man var
kommen ind paa en ny Bane, en Vei, hvis Retning og
Beskaffenhed først efterhaanden kunde blive undersøgt og
nærmere prøvet.\footnote[2]{Hvad Leibnitz skriver til Oldenburg i et Brev fra Paris, 27 August 1676: \textit{Fateor tamen nondum me quicquid in hoc genere desiderari potest consecutum, quanquam maximi momenti esse sciam,} gjælder ikke blot for Datiden, men tilbeels ogsaa for den følgende Tid. Jvfr. Gerhardt: Die Entdeckung der Differentialrechnung durch Leibnitz. Halle 1848. S. 26. „\textit{M. Maclaurin, l'un des plus grands Géometres de ce siecle} (18de Aarh.) \textit{éléva le premier sa voix au milieu des acclamations qu'attiroient les utilités réelles du Calcul des Infiniment-Petits. Il osa avertir du danger qu'on couroit en s'y livrant avec tant de confiance. Il n'y a rien de concevable, dit il, dans la supposition d'un nombre Infiniment-Grand ou Infiniment-Petit; \& le passage du Fini à l'Infini est obscur et incomprehensible.}« \textit{M. Saverien. Histoire critique du Calcul des Infiniment-Petits (i »Application de la Géométrie ordinaire«)} S. 7—8. „Da... Einiges in Newtons Verfahren nicht hinreichend klar und deutlich war, so könnte es nicht fehlen, daß dasselbe mannigfache Angriffe zu erleiden hatte. Gegen dieselben trat Maclaurin auf („\textit{Treatise of Fluxions}“), um einerseits alle Zweifel, die noch in der Fluxionsrechnung übrig seien, zu zerstreuen, andererseits die Leibnizische Differentialmethode anzugreifen. Der letzteren machte er zwei Dinge zum Vorwurfe, einmal nämlich, daß das Uendlichkleine für den menschlichen Verstand unfaßbar, sodann daß Leibnitz selbst den Begriff des Uendlichkleinen für eine bloße „Fiction“ erklärt habe.“ Weissenborn. S. 58. Endnu et characteristisk Træk af ovennævnte \textit{Histoire du Calcul des Infiniments-Petits.} Efter en kort Beskrivelse af den Archimediske Grændsemethode vedbliver Forfatteren S. XIV saaledes: „\textit{Cette méthode fut long-tems en usage. Dans la vûe de la} simplifier, on crut qu'on pouvoit se passer de concevoir des figures circonscrites et inscrites dans des aires curvilignes, comme étant toûjours finies et déterminables. L'expédient qu'on trouva pour cela, fut de substituer à ces figures finies et déterminables des élémens indivisibles ou Infiniment-Petits. Le nombre de ces élémens étant supposé infini, leur somme est évidemment égale à l'aire curviligne. On doit cette idée à Cavallieri. Quelque séduisante qu'elle lui parut par sa simplicité, elle ne le rassura pas sur son entiere certitude. Ce nombre infini lui paroissoit géométriquement indéterminable. Il tâchoit donc d'écarter cette idée de nombre infini, lorsqu' il reconnut bien des difficultés qu'il ne put résoudre. C'en fut assez pour lui rendre suspecte sa nouvelle Géométrie. Afin de l'appuïer, Cavallieri joignit des démonstrations incontestables à celles qu'il avoit déduites de ces principes hypothetiques; mais ce mêlange de vérités \& de suppositions, bien loin de rassurer les esprits, les inquiéta. On vit alors pour la premiere fois (selon la juste remarque de M. Maclaurin, dans l'Introduction de son Traité des Fluxions) des disputes parmi les Géometres.“} Men medens det Uendeliges Analyse

% --- p. 92 ---
\setcounter{footnote}{0}%
\opage{92}saaledes vistnok i en Henseende giver os noget aldeles Nyt,
er dette Nye dog ikke desto mindre, naar de Overgangsstadier,
Videnskabens Historie udviser, tages i Betragtning,
en naturlig Conseqvens af de oprindelige Forudsætninger, en
Udfoldning af de i det Gamle værende Spirer: Størrelsens
Udvikling til Forholdsstørrelse er her et Afgjørende. For at
see, hvorpaa det i Differentiationsprinciperne Eiendommelige
væsentlig beroer, mærke man sig Fremgangsmaaden ved de
enkelte Functioner, dernæst Behandlingen af de sammensatte,
og, hvad Methoden i dens Heelhed angaaer, Brugbarheden
af Taylors Række som Functionssystemernes Grundformel.

Hvad Enkeltfunctionerne angaaer, bliver jo Differentialcoefficienten bestemt derved, at Fremstillingen ligefrem svarer til Definitionen; saaledes er:
\[ \frac{d(a+x)}{dx}=\lim\frac{a+(x+h)-(a+x)}{h}=\lim\frac{h}{h}=1;\quad\frac{d(ax)}{dx} \]
% the note 2) of p. 91 runs over onto the foot of this page: set in full at its mark

% --- p. 93 ---
\setcounter{footnote}{0}%
\opage{93}\[ =\lim\frac{a(x+h)-ax}{h}=\lim\frac{ah}{h}=a;\quad\frac{dx^{n}}{dx}=\lim\frac{(x+h)^{n}-x^{n}}{h} \]
\[ =nx^{n-1} \]
o. s. v. Det Eiendommelige sees altsaa af Forholdsstørrelsernes
indbyrdes Forhold.

For at differentiere en sammensat Function udvikler
man den i Række efter stigende Potenser af $h$, men da
denne Række egentlig faaer sin Betydning ved de Enkeltfunctioner,
hvoraf den sammensatte bestaaer, maae disse særlig
udvikles i en Række for sig, af disse igjen en til Functionens
Sammensætning svarende enkelt Række dannes, og
de ubestemte Coefficienters Methode anvendes. Saaledes finder
man $\frac{d(y_{1}+y_{2})}{dx}$, idet $Y=y_{1}+y_{2}$, d. v. s. $F(x)=f_{1}(x)+f_{2}(x)$, ved at udvikle $F(x+h)=f_{1}(x+h)+f_{2}(x+h)$ i Række; altsaa:
\[ F(x+h)=F(x)+F'(x)\frac{h}{1}+\ldots. \]
\[ =f_{1}(x)+f_{2}(x)+f_{1}'(x)\frac{h}{1}+f_{2}'(x)\frac{h}{1}+\ldots. \]
; da nu, ifølge de
ubestemte Coefficienters Methode, $F'(x)=f_{1}'(x)+f_{2}'(x)$,
saa er $\frac{d.Y}{dx}=\frac{d(y_{1}+y_{2})}{dx}=\frac{dy_{1}}{dx}+\frac{dy_{2}}{dx}$. Er Opgaven at
finde $\frac{d(y_{1}\cdot y_{2})}{dx}$, saa bliver $F(x+h)=f_{1}(x+h)f_{2}(x+h)$
\[ =F(x)+F'(x)\frac{h}{1}+\ldots=(f_{1}(x)+f_{1}'(x)\frac{h}{1}+\ldots)(f_{2}(x)+f_{2}'(x)\frac{h}{1}+\ldots)=f_{1}(x)\cdot f_{2}(x)+(f_{1}(x)\cdot f_{2}'(x)+f_{2}(x)f_{1}'(x))\frac{h}{1} \]
$+\ldots$, altsaa: $F'(x)=f_{1}(x)f_{2}'(x)+f_{2}(x)f_{1}'(x)$, det vil sige:
\[ \frac{d(y_{1}y_{2})}{dx}=y_{1}\frac{dy_{2}}{dx}+y_{2}\frac{dy_{1}}{dx}, \]
o. s. v.

Det kan i denne Sammenhæng naturligviis ikke være
Opgaven at anføre de enkelte „Differentiationsregler“; Begrebsudviklingen
dreier sig i det Væsentlige om Betingelserne
for Taylorsrækkens Anvendelse. For at erholde den enkelte

% --- p. 94 ---
\setcounter{footnote}{0}%
\opage{94}Differentialcoefficient $f'(x)$ alene, maa man jo lade Rækkens
øvrige Led bortfalde. Denne Leddenes Bortfalden faaer
nu vistnok en exact Begrundelse derved, at man i Ligningen:

\[ \lim\frac{f(x+h)-f(x)}{h}=f'(x)+f''(x)\frac{h}{1.2}+f'''(x)\frac{h^{2}}{1.2.3}+\ldots \]

antager $h=0$; men ved at gjøre $h$ til $0$, gjør man med det Samme $f'(x)$ til $\frac{0}{0}$, og de ovenfor paapegede Vanskeligheder
vende tilbage. Betragter man derimod $h$ ikke som $0$, men som $dx$, da faaer man, for $f(x+h)=y+k$, vel en virkelig Differentialcoesficient $f'(x)=\frac{dy}{dx}$ istedetfor en i % PRINTED AS IS: Differentialcoesficient (Rettelser: Differentialcoefficient)
sig modsigende Betegnelse; men saa synes der igjen at tilkomme
Summen af Rækkens øvrige Led en saadan Betydning,
at de, naar man fordrer et exact Udtryk, ikke kunne
bortkastes.

Denne Mislighed hæves ved nærmere Betragtning af
den omhandlede Række. Til $f(x)=(x-a)^{n}$ svarer $(x+h-a)^{n}=((x-a)+h)^{n}=(x-a)^{n}+n(x-a)^{n-1}\frac{h}{1}+n(n-1)(x-a)^{n-2}\frac{h^{2}}{1.2}\ldots.+n(n-1)\ldots(n-r+1)(x-a)^{n-r}\frac{h^{r}}{[r]}\ldots..$

Er nu $1)$ $n$ positiv heeel, da er Rækken afsluttet med $h^{n}$, og, forsaavidt vi endnu betragte den som Taylors Formel, $\frac{d^{n+p}(x-a)^{n}}{dx^{n+p}}=0$, idet $\frac{d^{n}(x-a)^{n}}{dx^{n}}=c$. % PRINTED AS IS: heeel (Rettelser: heel)
Er $2)$ $n$ negativ, da sees, at alle Differentialcoefficienterne
maae indeholde $(x-a)$ i Nævneren, og
altsaa, naar $x=a$, $\frac{d^{n}y}{dx^{n}}=\infty$.
Er $3)$ $n$ brudden, t. Ex. $n+1>r>n$, da er det klart, at $\frac{d^{r}(x-a)^{n}}{dx^{r}}=n(n-1)$

% --- p. 95 ---
\setcounter{footnote}{0}%
\opage{95}\[ \ldots.(n-r+1)(x-a)^{n-r}=\frac{n(n-1)\ldots.(n-r+1)}{(x-a)^{r-n}}, \]

og at altsaa $\frac{d^{r}(x-a)^{n}}{dx^{r}}$ maa indeholde $x-a$ i Nævneren.
Differentialcoefficienterne af høiere end $r$te Orden maae da saaledes, for $x=a$, alle blive $=\infty$.
Resultatet er nu
dette: Taylors Formel, skjøndt almindelig gjældende, kan for
specielle Værdier af $x$ blive ubrugelig; men, \emph{naar} Rækken er brugelig, er den for $\mathrm{Lim}\frac{f^{n+1}(x)}{(n+1)f^{n}(x)}h<1$
altid convergent, og destomere convergent, jo
mindre $h$ er, med andre Ord: jo mindre $h$ er, desto
færre Led behøves for at opnaae en vis Grad af
\emph{Nøiagtighed.}

At her i denne Opfattelse dog endnu er noget Utilfredsstillende,
har Hegel, og det ikke uden Grund, udhævet paa
følgende Maade: „Lagrange hat bekanntlich die ursprüngliche
Methode Newtons, die Methode der Reihen, wieder aufgenommen,
um die Schwierigkeiten, welche die Vorstellung
des Uendlich-Kleinen, so wie derjenigen, welche die Methode
der ersten und letzten Verhältnisse und Grenzen mit sich
führt, überhoben zu seyn. Es ist von seinem Functionen-Kalkul,
dessen sonstige Vorzüge in Rücksicht auf Präcision,
Abstraktion und Allgemeinheit anerkannt genug sind, als
hierher gehörig nur dieß anzuführen, daß er auf dem Fundamentalsatze
beruht, daß die Differenz ohne daß sie Null
werde, so klein angenommen werden könne, daß
jedes Glied der Reihe die Summe aller folgenden
an Größe übertreffe. --- Es wird auch in dieser Methode
von den Kategorien vom Zuwachs und von der Differenz
der Function angefangen, deren veränderliche Größe
den Zuwachs erhalte, womit die lästige Reihe hereinkommt,
von der ursprünglichen Function; so wie im Verfolg die
wegzulassenden Glieder der Reihe nur in der Rücksicht, daß
sie eine Summe konstituiren, in Betracht kommen, und der

% --- p. 96 ---
\setcounter{footnote}{0}%
\opage{96}Grund sie wegzulassen, in das Relative ihres Quantums
gesezt wird.“\footnote[1]{Wissensch. der Logik. 1ster Theil. S. 316—317.}

For at hæve dette Utilfredsstillende behøve di dog ikke % PRINTED AS IS: di (for vi)
med Hegel at fordre en saa qvalitativ Afslutning, at vi
derved i Grunden opgive Qvantiteten; det er, for at overbevise
sig om, at en absolut Nøiagtighed kan opnaaes, tilstrækkeligt
at tydeliggjøre sig den qvalitative Forandring, Rækken
undergaaer, naar $h$ forvandles til $dx$. I $\frac{f(x+dx)-f(x)}{dx}$
\[ =f'(x)+f''(x)\frac{dx}{1}+f'''(x)\frac{dx^{2}}{1.2}+\ldots. \]
maa det fra
Mathematikens Standpunkt udtrykkelig fremhæves, at $f''(x)dx$ ikke blot er meget mindre, men uendelig mange Gange mindre end $f'(x)$, at de med andre Ord ere „incomparable“.
Da nu det Samme gjentager sig i Sammenstillingen af $f'''(x)\frac{dx^{2}}{1.2}$ med $f''(x)dx$ o. s. v., saa indsees, at Rækkens
Led bortfalde, ikke fordi de ere saa ubetydelige, at den Unøiagtighed,
der begaaes ved deres Udeladelse er meget ringe;
men fordi de som incomparable Led, som Led, der
ifølge deres Beskaffenhed tilhøre en anden
Sphere, i Begrebet maa bortfalde.

I nogle Tilfælde, hvor det virkelig seer ud, som om
man ved at udelade et Led af høiere Orden begik en Feil,
viser det sig, at Feilen erstattes, sva at der dog alligevel  % PRINTED AS IS: sva (for saa)
ikke begaaes nogen Feil. En saadan „Feilens Compensation“
kan t. Ex. sees i følgende Tilfælde. For den almindelige
Parabel: $y^{2}=px$ bliver, ifølge Differentiationsmethoden,
Subtangenten bestemt ved $subtg=y\frac{dx}{dy}$; $(y+dy)^{2}=p(x+dx)$ og $dy=\frac{pdx}{2y}-\frac{dy^{2}}{2y}$; dette giver, idet $\frac{dy^{2}}{2y}$ bortkastes,

% --- p. 97 ---
\setcounter{footnote}{0}%
\opage{97}\[ \frac{dx}{dy}=\frac{2y}{p}. \]
Man erholder altsaa: $subtg.=y\frac{dx}{dy}=\frac{2y^{2}}{p}$, eller, da $y^{2}=px$, $subtg.=2x$. Ved denne Fremgangsmaade synes det vistnok at være en Feil, at Størrelsen $\frac{dy^{2}}{2y}$ blev udeladt; men ved at lade Differentialet af Parablens
Ordinat i det consecutive Punkt falde sammen med Differentialet
af Tangentens Ordinat, har man for Parablens
Vedkommende fra Begyndelsen af gjort $dy$ for stor. Ifølge
en af Apollonius opstillet, og altsaa af Differentialregningen
uafhængig Sætning, haves Subtangenten $=2x$, og, idet vi ved $u$ betegne Forskjellen mellem Differentialet af Tangentens og Parablens Ordinat, haves $\frac{2x}{y}=\frac{dx}{dy+u}$, hvoraf
sees, at $u=\frac{dy^{2}}{2y}$, og at altsaa, det Stykke, der skulde fra-brages $dy$, nøiagtig svarer til det Led, der i Medfør af
Differentiationens Princip blev bortkastet\footnote[1]{Weissenborn. Anf. Skr. S. 156—58.}

Væsentlig forstaaet bliver i Ligningen $dy=\frac{pdx}{2y}-\frac{dy^{2}}{2y}$ det sidste Led ikke bortkastet, fordi $dy^{2}$ antages at være
en endnu mindre Deel af en Linie end $dx$ eller $dy$, men
fordi $dy^{2}$ ikke er eensartet med en Linie, men med en Linies
Grændse, et Punkt, et forsvindende Moment. Betyder
nu Bortkastelsen af et Led, at dette, for at kunne bortkastes,
maa være forvandlet fra blivende Størrelse til forsvindende Moment, da maa jo den exacte Handling, seet
under Reflexionen, søge sin sidste retfærdiggjørende Grund i
en Dialektik.

% --- p. 98 ---
\setcounter{footnote}{0}%
\opage{98}
\begin{center}c) Integralet: Elementet og den endelige Størrelse.\end{center}

Hvad Differentialet egentlig er, kan først tilfulde belyses
ved dets Modsætning, nemlig Integralet. „I Modsætning
til Differentialet som det Negative er Integralet det
Positive. Er Differentialet et Punkt, saa er Integralet en
Linie; er Differentialet en Linie, saa er Integralet en Flade,
er Differentialet en Flade, saa er Integralet et Legeme.
Betyder Differentialet et Moment i Liniens uendelige Discretion,
saa betyder Integralet de discrete Momenters Opgaaen
i det endelige Continuum; betegner Differentialet det
uendelig lille Led, saa betegner Integralet Summen af slige
Led i uendelig stort Antal;“\footnote[1]{Philos. og Mathem. S. 94, S. 60.} bestemmes Differentialet ved
den afledede Function $f'(x)dx$, saa bestemmes igjen Integralet ved den oprindelige $f(x)$.

Idet den uendelige Summation
efterhaanden sees at falde sammen med den Opgave
at tilbageføre en afledet Function til den tilsvarende oprindelige\footnote[2]{Idet man definerer Integralet saaledes, at $y=\int f'(x)\,dx$ er den Størrelse, som ved at differentieres med Hensyn til x giver $f'(x)$, faaer man unegtelig ved Integrationen \emph{ganske nøiagtig} den søgte Function, nemlig $f(x)$; men denne paa Definitionen støttede Nøiagtighed staaer og falder med det Spørgsmaal, om Definitionen selv er abæqvat, d. e., om man nu ogsaa virkelig er berettiget til at sætte $\frac{dy}{dx}=\lim\frac{f(x+h)-f(x)}{h}=f'(x)$, eller om ikke Rækken for $\frac{dy}{dx}$ dog egentlig er: $f'(x)+f''(x)\frac{dx}{1.2}+f'''(x)\frac{dx^{2}}{1.2.3}+\ldots$ og om man saa ikke har begaaet en Unøiagtighed ved af denne Række at bortkaste alle efter $f'(x)$ følgende Led. Ved en given Definition, en vedtagen Betegnelse, kan en Tankeforvirring forebygges; men dermed er et Problem endnu ikke løst.}
vil Udviklingen af Integralets Begreb tillige vise
sig at falde sammen med Uendelighedsproblemets Løsning.

% --- p. 99 ---
\setcounter{footnote}{0}%
\opage{99}
\begin{center}α) Integrerende Summation: Moment og Element.\end{center}

Den mathematiske Sætning, at enhver endelig Størrelse
kan antages at bestaae af uendelig mange uendelig
smaae Dele, er i sin Almindelighed jo vistnok trivial, men
dog altid ligesaa betydningsfuld som den logiske Grundsætning,
at enhver Ting er, hvad den er. Det Triviale ligger
ikke i Abstractionen, men deri, at den tomme Abstraction
fastholdes som et Værende, hvorved det Dialektiske oversees,
og Begrebsudviklingen standser. Som et givet Hele bestaaer
den endelige Størrelse naturligviis af Dele; opfattes nu
dette Hele fra Eenhedens Side, sees Continuiteten; opfattes
det fra Delenes Side, sees Discretionen. Den Reflexionsfordobling,
ifølge hvilken det væsentlige Hele har sin Bestaaen
i de bestaaende Dele, kan Størrelsen som det rene
Qvantum imidlertid ikke fremstille; det i Qvantitetskategorien
Umiddelbare er endnu saa fremherskende, at den ene af
Reflexionens Sider forsvinder, naar den anden sættes: idet
Discretionen sættes, brydes Continuiteten; idet Continuiteten
sættes, forsvinder Discretionen. Adskillelsen af Discretion
og Continuitet har igjen foranlediget en Adskillelse af discrete
og continuerede Størrelser, som om de udgjorde tvende
for sig bestaaende Arter af Størrelser\footnote[1]{„Størrelser kunne enten være discrete (adskilte) eller continuerte (sammenhængende). De første ere blotte Samlinger af hvilkesomhelst eensartede fra hinanden adskilte Enkeltheder (f. Ex. en Men- neskehob, en Dreslok, en Steenbunke), saa at Størrelsens Grad eller Værdi maa angives ved et Tal. De sidste ere Heelheder, hvor ingen Adskillelse findes mellem hvilkesomhelst eensartede Dele (f. Ex. en Længde, en Vægt, et Tidsrum). To eensartede continuerte Størrelser kunne stedse sammenlignes i Henseende til deres Værdier. Det simpleste Tilfælde er det, hvor den ene Størrelse er nøiagtigen indeholdt et heelt Antal Gange, $p$, i den anden, hvor $p$ kan være 1, 2, 3... Tages den første som Enhed, kan den anden siges at være liig med Tallet p“... Ramus. Elementær Geometrie, Kbhvn 1850; Indledning.} At en saadan
Adskillelse blot er foreløbig, indlyser imidlertid deraf, at en
continueret Størrelse kun kan bestemmes, idet man sætter
en anden dermed eensartet som Eenhed, og nu undersøger,
hvormange Gange denne indeholdes i hiin; men ved en
saadan Undersøgelse bliver den continuerede Størrelse jo
netop seet under Discretionens Synspunkt. Adskillelsen af
discret og continueret Størrelse reducerer sig i Begrebet til
en Adskillelse af endelig og uendelig Discretion. ¹) Ved den

% --- p. 100 ---
\setcounter{footnote}{0}%
\opage{100}endelige Discretion bliver Continuiteten brudt, med den
uendelige Discretion vender den brudte Continuitet restitueret
tilbage; i den endelige Discretion ere Delene som tilværende
Led endnu særlige Størrelser\footnote[1]{Overgangen fra Differens- til Differentialcoefficienten sees, hvad Taylorsrækken angaaer, saaledes: $y_{0}, y_{1}, y_{2} \ldots y_{n}$ være $f(x), f(x+h), f(x+2h), \ldots f(x+nh)$; man har da: $y_{1}=y_{0}+\Delta y_{0}$, $y_{2}=y_{0}+2\Delta y_{0}+\Delta^{2}y_{0}$, $y_{3}=y_{0}+3\Delta y_{0}+3\Delta^{2}y_{0}+\Delta^{3}y_{0}, \ldots$ $y_{n}=y_{0}+\frac{n}{1}\Delta y_{0}+\frac{n(n-1)}{1.2}\Delta^{2}y_{0}+\frac{n(n-1)(n-2)}{1.2.3.}\Delta^{3}y_{0}$ $+\ldots$. altsaa $f(x+nh)=f(x)+\frac{n}{1}\Delta.f(x)+\frac{n(n-1)}{1.2}\Delta^{2}.f(x)$ $+\ldots.$; sættes $nh=i$, saa er $n=\frac{i}{h}$. $f(x+i)=f(x)+\frac{i}{1}\Delta\frac{f(x)}{h}+\frac{i(i-h)}{1.2}\Delta^{2}\frac{f(x)}{h^{2}}+\ldots$; forvandles nu $h$ til $dx$, erholder} i den uendelige Discretion

Adskillelsen af discret og continueret Størrelse er altsaa her
bleven forklaret saaledes, at den i Forklaringen ophæver sig selv.
En Samling af 20 Mennesker er en Hob, en Samling af 19,
af 18, af 17, af 16\dots{}. Mennesker ligeledes: Størrelsen: Hob
er, saaledes opfattet, en continueret Størrelse. Skal Kjendetegnet
paa en discret Størrelse være dette, at dens Værdi „maa angives
ved et Tal“, da ere saadanne Størrelser som: „en Længde, en
Vægt, et Tidsrum“, lutter discrete Størrelser; thi en Alen er
4 Qvarteer, et Pund er 32 Lod, en Minut er 60 Secunder.
Som enhver Deel af en Længde selv er en Længde, er ogsaa enhver
Deel af Hoben, naar vi holde os til det rene Størrelsesbegreb,
og ikke indblande uvedkommende Forestillinger, selv en Hob.
% the note 1) of p. 99 runs over onto the foot of this page: set in full at its mark

% --- p. 101 ---
\setcounter{footnote}{0}%
\opage{101}ere Delene forsvindende Led, Størrelsesmomenter, dialektiske
Mellembestemmelser af Noget og Intet\footnote[1]{Hermed stemmer det, at Functionens Begreb opfattes „als reale Existens des discret-continuirlichen Qvantums“. H. Schwarz. Anfrt. Skr. S. 22.} Som Moment
er Delen i en uendelig Overgang til Intet, uden dog at
blive til Intet; som Element er Delene i en uendelig Overgang % PRINTED AS IS: Delene (Rettelser: Delen)
til Noget, et Noget, hvoraf det Endelige frembringes.
Til den endelige Discretion, hvori Delene holde sig som
tilværende Led, svarer en endelig Summation: det Hele
fremkommer ved en Sammentællen af Delene; til den uendelige
Discretion, hvori Delene ere Momenter, svarer en
uendelig Summation, en Forvandling af Momentet til

Forholdet imellem Differential- og Integralregningen er i
det Endeliges Analyse allerede præfigureret ved Modsætningen
mellem den directe og den omvendte Differensregning. Differensregningen
bestaaer i „at finde Differenserne af opgivne Functioner
og derefter Differenserne af disses Differenser. Saaledes er $\Delta.x^{n}=x^{n+1}-x^{n}=x^{n}(x-1)$. Heraf sees, at hvis man tager
Differensen af Leddene i en Qvotientrække, faaes en ny Qvotientrække.
Til $1, x, x^{2}, x^{3}, x^{4} \ldots.$ svarer $x-1, x(x-1), x^{2}(x-1), x^{3}(x-1)\ldots.$ og $(x-1)^{2}, x(x-1)^{2}, x^{2}(x-1)^{2}$ er igjen dannet af Differenserne i den foregaaende Række.“
A. Steen. Reen Mathem. S. 338.

„Den Beregning, hvorved man finder en Function af dens
Differens, kaldes omvendt Differensregning, og i Modsætning
dertil bliver den egentlige Differensregning kaldt directe. Den
omvendte kaldes ogsaa Integration integrere i Modsætning til
differentiere). Af $\Delta.s_{n}=u_{n}$ faaes Integralet $s_{n}=\Sigma.u_{n}$, af
$\Delta.x^{n}=x^{n}(x-1)$ Integralet „$x^{n}=\Sigma.x^{n}(x-1)$; men da heri $x-1$ er en fælleds Factor for alle Led under $\Sigma$, som er uafhængig
af $n$ (constant), kan den sættes udenfor $\Sigma$, saa at $x^{n}=(x-1)\Sigma.x^{n}$ eller $\Sigma.x^{n}=\frac{x^{n}}{x-1}$. Anfrt. Skr. S. 340.
% NOTE: its mark (Størrelser) i) was not located: the note is set at the page end
\footnote[1]{man $f(x+i)=f(x)+\frac{d.f(x)}{dx}\frac{i}{1}+\frac{d^{2}f(x)}{dx^{2}}\frac{i^{2}}{1.2}+\ldots=f(x)+f'(x)\frac{i}{1}+f''(x)\frac{i^{2}}{1.2}+\ldots.$ (the opening of the p. 100 footnote "Overgangen ..."; the mark is on p. 100)}

% --- p. 102 ---
\setcounter{footnote}{0}%
\opage{102}Element, en intellectuel Handling, der bestemmer det Endelige
som Produkt af det Uendelige. Forholdet imellem Elementet
og den endelige Størrelse kan da for det Første oplyses
ved Summationen.

Er Størrelsen $x$ deelt saaledes, at hver Deel $\Delta x=\frac{x}{4}$,

bliver den hele Størrelse:
\[ 4\frac{x}{4}=\Sigma\Delta x=x \]
; tænkes $x$ derimod
deelt saaledes, at den enkelte Deel ikke er en endelig
Størrelse, men et til Element forvandlet Moment, kan $x$ som endelig Størrelse ikke fremkomme ved Addition: $dx+dx+\ldots.$, men kun ifølge Begrebet derved, at $x$ opfattes
som en uendelig Sum af $dx$, altsaa som $\int_{0}^{x}dx=x$. At Discretionsmomentet $dx$ ikke er $0$, men et Element, sees nu
ligefrem deraf, at en uendelig Sum af Nuller ikke kan give
Andet end $0$; men er dx ikke absolut Nul, saa er $\int_{0}^{x}$ heller
ikke absolut $\infty$. Betegnes $\int_{0}^{x}dx=x$ ved $\infty .0$, maa det

udtrykkelig fremhæves, at $\infty$ svarer til $\frac{x}{dx}$, idet $0$ svarer til $dx$. Forskjellen imellem $\Sigma_{0}^{x}\Delta x$ og $\int_{0}^{x}dx$ er ingenlunde den,
at kun det første Udtryk skulde angive noget Bestemt,
det sidste derimod opløse sig i tom Formalisme; men
Forskjellen er, at $x$ opfattes i første Udtryk som discret,
i sidste Udtryk derimod som en continueret Størrelse.

Denne Forskjel antager en endnu mere articuleret Form,
naar vi sammenligne Summerne\footnote[1]{Ved $\Sigma_{0}^{x}f'(x)\Delta x$ forstaae vi her Summen af alle de Værdier $f'(x)\Delta x$ gjennemløber, naar x gjennemløber samtlige Værdier mellem 0 og x med Differensen $\Delta x$.} $\Sigma_{0}^{x}f'(x)\Delta x$ og $\int_{0}^{x}f'^{(x)}dx$.

% --- p. 103 ---
\setcounter{footnote}{0}%
\opage{103}I den endelige Sum $\Sigma_{0}^{x}x\Delta x=(\Delta x+2\Delta x+3\Delta x+\ldots n\Delta x)\Delta x=(1+2+3+\ldots n)(\Delta x)^{2}$, viser Differensen $\Delta x=2\Delta x-\Delta x=3\Delta x-2\Delta x\ldots.$ den endelige Discretion,
et Spring, en endelig Afstand imellem de endelige
Led; i den uendelige Sum $\int_{0}^{x}xdx=(dx+2dx+3dx+\ldots ndx)dx=(1+2+3+\ldots n)dx^{2}$, seer man af Differentialet $dx=2dx-dx=3dx-2dx\ldots$, at Discretionen
er uendelig, at her altsaa, med andre Ord, fra $0$ til $x$ er en
continuerlig Overgang. I Ligningen $\Delta x=\frac{x}{n}$ er $n$ endelig; i Ligningen $dx=\frac{x}{n}$ er $n$ derimod uendelig; som Følge
heraf bliver da
\[ \Sigma_{0}^{x}x\Delta x=\frac{n(n+1)}{2}\frac{x^{2}}{n^{2}}=\frac{x^{2}}{2}+\frac{x^{2}}{2n} \]
, medens
\[ \int_{0}^{x}x\,dx \]
derimod giver
\[ \mathrm{Lim}\frac{(1+\frac{1}{n})}{2}x^{2}=\frac{x^{2}}{2} \]
. Medens
\[ \Sigma_{0}^{x}x^{2}\Delta x=[\Delta x^{2}+(2\Delta x)^{2}+(3\Delta x)^{3}+\ldots(n\Delta x)^{2}]\Delta x= \] % PRINTED AS IS: (3\Delta x)^{3} (for (3\Delta x)^{2})
\[ (1+2^{2}+3^{2}+\ldots n^{2})(\Delta x)^{3}=\frac{n(n+1)(2n+1)}{1.2.3}(\Delta x)^{3} \]
\[ =\frac{n(n+1)(2n+1)}{1.2.3}\frac{x^{3}}{n^{3}} \]
giver
\[ \frac{x^{3}}{3}+\frac{x^{3}}{2n}+\frac{x^{3}}{6n^{2}} \]
paa Grund
af Discretionen: giver
\[ \int_{0}^{x}x^{2}dx=[(dx)^{2}+(2dx)^{2}+\ldots.(ndx)^{2}]dx=(1+2^{2}+3^{2}+\ldots n^{2})dx^{3}=\frac{n(n+1)(2n+1)}{1.2.3} \]
\[ \frac{x^{3}}{n}=\mathrm{Lim}\frac{(1+\frac{1}{n})(2+\frac{1}{n})x^{3}}{1.2.3} \] % PRINTED AS IS: \frac{x^{3}}{n} (for \frac{x^{3}}{3})
derimod $\frac{x^{3}}{3}$ paa Grund af
Continuiteten. Til
\[ \Sigma_{0}^{x}x^{3}\Delta x=[\Delta x^{3}+(2\Delta x)^{3}+\ldots(n\Delta x)^{3}]\Delta x=(1+2^{3}+\ldots n^{3})(\Delta x)^{4}=(1+2+3+\ldots n)^{2}(\Delta x)^{4} \]
svarer
\[ \frac{n^{2}(1+n)^{2}}{2^{2}}\frac{x^{4}}{n^{4}}=\frac{x^{4}}{4}+\frac{x^{4}}{2n}+\frac{x^{4}}{4n^{2}} \]
, efterdi $n$ er

% --- p. 104 ---
\setcounter{footnote}{0}%
\opage{104}endelig; til
\[ \int_{0}^{x}x^{3}dx=[dx^{3}+(2dx^{3}+\ldots(ndx)^{3}]dx=(1+2^{3}+\ldots.n^{3})dx^{4}=(1+2+3+\ldots.n)^{2}dx^{4} \]

svarer derimod
\[ \frac{n^{2}(1+n)^{2}}{2^{2}}\frac{x^{4}}{n^{4}}=\mathrm{Lim}\frac{(\frac{1}{n}+1)^{2}}{4}x^{4}=\frac{x^{4}}{4} \]
; o. s. v.

I den uendelige Summation tydeliggjøres, hvad allerede
tidligere (II., $\gamma$) er paapeget, det Endeliges og det
Uendeliges dialektiske Eenhed. Den Sætning, at en
Størrelse bestaaer af et uendeligt Antal uendelig smaae Dele
har da i sig et ikke blot ligesaa rigt, men endog rigere Anlæg
til articulerende Bestemmelser, end den Sætning, at
den endelige Størrelse bestaaer af et endeligt Antal endelige
Dele. Elementerne $dx, xdx, x^{2}dx$ o. s. v. ere i Forhold til
de paagjældende Størrelser $x, \frac{x^{2}}{2}, \frac{x^{3}}{3}\ldots$ ligesaa eiendommelig
bestemte, som en hvilkensomhelst endelig Deel $\Delta x, n\Delta x\ldots$, i Forhold til det tilsvarende endelige Hele. Saalænge
Forstanden søger al Forskjel og al Bestemthed i det Endelige,
kan den vel behandle Størrelserne som led i en omfattende
Operationsmechanik, men den er ikke istand til at
fatte, hvad Størrelse egentlig er, ikke istand til at erkjende
Størrelsens Væsen. Til en Erkjendelse af Størrelsens Væsen
udfordres en Erkjendelse af tens Oprindelse, dens Vorden
og Forsvinden, dens Opløsning i sit Moment, dens
Fremgang af sit Element.

\begin{center}β) Integrationsprinciper: afledede Functioners Forhold til de oprindelige.\end{center}

Den uendelige Summation kan ikke ligefrem udføres;
for at afgive Analogier, som ovenstaaende, maa allerede den
endelige Summation undergaae en saa væsentlig Forvandling,
at den ikke længere beroer paa simpel Addition. Den uende%

% --- p. 105 ---
\setcounter{footnote}{0}%
\opage{105}lige Summation bestemmes ved Forholdet mellem Element
og Størrelse; men dette Forhold bestemmes igjen, idet det
fuldstændige Størrelsesbegreb falder sammen med Functionsbegrebet,
ved Forholdet mellem den afledede Function og
den tilsvarende oprindelige. Som Element er det paagjældende
Differential: $d f(x)=dy$ et Produkt af den afledede
Function, Differentialcoefficienten $f'(x)$, og det Moment, $dx=dx$, hvorved den Uafhængige varierer. Forsaavidt
Størrelsen er en Function af sig selv og Eenheden: $f(x)=x$, og $dx=df(x)=f'(x)dx=dy$, er her imellem Moment
og Element kun den formelle Forskjel, at dx som et Differentieringsudtryk
er Moment, som Integrationens reale Mulighed
derimod Element. I væsentlig bestemte Functioner har
Forskjellen væsentlig Betydning; saaledes er i $f(x+dx)=(x+dx)^{2}=y+dy$ Differentialet dx kun Moment, Differentialet $dy=df(x)=f'(x)dx$ derimod det egentlige
Element.

Fra Forskjellen mellem Moment og Element faaer Integrationen
først i Begrebet en alsidig Belysning. Opfattes $\int_{0}^{x}df(x)=\int_{0}^{x}f'(x)dx$ som Summen af de uendelig mange
Led, der fremkommer, idet $x$ i $f'(x)dx$ continuerlig gjennemløber
Værdierne fra $0$ til $x$, da vil denne Sum ifølge
Beskaffenheden af $f'(x)dx$, d. e. ifølge Elementets Beskaffenhed,
give, for $f(x)=x$, Summen: $dx+dx+\ldots=x$; for $f(x)=x^{2}$, Summen:
\[ \int_{0}^{x}d(x^{2})=\int_{0}^{x}f'(x)dx=\int_{0}^{x}2xdx \]

\[ =2\int_{0}^{x}xdx=2(1+2+3+\ldots.n)dx^{2}=\mathrm{Lim}\,2\frac{x^{2}}{n^{2}} \]
\[ (1+2+3+\ldots.n)=\mathrm{Lim}\,2\frac{(n+1)}{1.2}x^{2}=\mathrm{Lim}(1+\frac{1}{n}) \]
\[ x^{2}=x^{2} \]
, o. s. v. Er det nu heraf indlysende, hvorledes den
uendelige Summation forvandler sig til Integration, idet
Integrationen, derved at Eftertrykket falder paa Elementet $f'(x)dx$, og følgelig paa $f'(x)$, principielt løser sin Opgave

% --- p. 106 ---
\setcounter{footnote}{0}%
\opage{106}at gaae fra Differentialet tilbage til Integralet, fra den afledede
Function tilbage til den oprindelige: da vil det nu
ogsaa kunne vise sig, hvori Forskjellen mellem endelig og
uendelig Summation væsentlig er begrundet.

Tillægges Udtrykket $(x+\Delta x)^{2}=x^{2}+2x\Delta x+(\Delta x)^{2}$ geometrisk Betydning, da sees det, at alle Leddene ere
comparable, forsaavidt de alle betegne Arealer; men forandres
nu $\Delta x$ til $dx$, da giver $(x+dx)^{2}=x^{2}+2xdx+dx^{2}$

incomparable Led, idet $x^{2}$ er et Areal, $xdx$ Arealets Element og $dx^{2}$ Elementets Moment. Hvad Legemsformen
angaaer, vil $(x+\Delta x)^{3}=x^{3}+3x^{2}\Delta x+3x(\Delta x)^{2}+(\Delta x)^{3}$ give comparable Led, nemlig lutter Legemer; medens $(x+dx)^{3}=x^{3}+3x^{2}dx+3xdx^{2}+dx^{3}$ derimod giver incomparable
Led, da $x^{3}$ jo er et Legeme, $x^{2}dx$ en Flade, $xdx^{2}$ en Linie
og dx et Punkt. Stilles nu den Opgave at bestemme
Størrelsen af den Tilvært et Areal, $a^{2}$, faaer, ved at summere
de Tilvæxter, dette Areal erholder, idet $x$ fra den
endelige Værdi $a$, ved Gjentagelse af Tilvæxten $\Delta x$, har naaet Værdien $b$; da ligger det i Sagens Natur, at den
endelige Summation maa i hver Addend have alle Tilvæxtens

Led, saa at Summen bliver $\sum_{x=a}^{x=b}[2x\Delta x+(\Delta x)^{2}]$, ikke $\sum_{x=a}^{x=b}2x\Delta x$.

Er det derimod Opgaven at summere de Tilvæxter, Arealet
$a^{2}$ erholder, idet $x=a$ ved Tilvæxter af $dx=dx$ er voxet
indtil $x=b$; da maa den uendelige Summation udelukkende
holde sig til Arealets Element, saa at Summen bliver:
\[ \int_{x=a}^{x=b}2xdx \]
, ikke
\[ \int_{x=a}^{x=b}(2xdx+dx^{2}) \]
.\footnote[1]{$\int_{a}^{b}(2x\,dx+dx^{2})=2a\,n\,dx+n^{2}dx^{2}$, hvor $n=\frac{b-a}{dx}$, altsaa $\int_{a}^{b}(2xdx+dx^{2})=2a(b-a)+(b-a)^{2}=b^{2}-a^{2}$. Hvorimod $\int_{a}^{b}2xdx=2a n\,dx+n^{2}dx^{2}-n\,dx^{2}=b^{2}-a^{2}-(b-a)dx$. Ville vi nu see denne Differens under Synspunktet af „Feilenes Compensation“ (jvfr. III., b, $\gamma$); da have vi ved at udelade $dx^{2}$ af dy virkelig begaaet en Feil, men denne Feil bliver saa igjen compenseret derved, at vi i $\int_{a}^{b}2xdx=b^{2}-a^{2}$ udelade det negative Led ÷ (b — a) dx. Den væsentlige Grund, hvorfor (b — a) dx udelades er imidlertid den, at her ved dette Udtryk kun betegnes en Linie, ikke et Areal.} At Elementets Moment

% --- p. 107 ---
\setcounter{footnote}{0}%
\opage{107}bortfalder, er jo ligefrem begrundet deri, at et Moment,
der allerede er indeholdt i Elementet selv, ikke kan faae
særlig Betydning.\footnote[1]{Hvorledes en Tænker, som G. Berkeley f. Ex. har søgt efter en Begrebsbestemmelse, hvori Forholdet mellem Moment, Element og endelig Størrelse kunde fastholdes, sees blandt Andet af Følgende: Because there is no Number of Parts so great, but it is possible there may be a Line containing more, the Inchline is said to contain Parts more 'than any assignable Number; which is true, not of the \emph{Inch} taken absolutely, but only for the Things signified by it. But Men not retaining that Distinction in their Thoughts, slide into a belief that the small particular Line described on Paper contains in it self Parts innumerable. There is no such thing as the ten-thousandth Part of an Inch; but there is of a \emph{Mile} or \emph{Diamether} of the \emph{Earth}, which may be signified by that Inch. When therefore I delineate a Triangle on Paper, and take one side not above an Inch, for Example, in length to be the \emph{Radius}: This I consider as divided into ten thousand or an hundred thousand Parts, or more. For though the ten-thousandth Part of that Line considered in it self, is nothing at all, and consequently may be neglected without any Error or Inconveniency; yet these described Lines being only Marks standing for greater Quantities, whereof it may be the ten-thousandth Part is very considerable, it follows, that to prevent notable Errors in Practice, the \emph{Radius} must be taken of ten thousand Parts, or more. Principles of human knowledge. Lond. 1734. S. 147—148.}
% the note 1) of p. 106 runs over onto the foot of this page: set in full at its mark

% --- p. 108 ---
\setcounter{footnote}{0}%
\opage{108}Som Arealets Element er $xdx$ en dialektisk Mellembestemmelse
af Linie og Flade, som Elementets Moment er $dx^{2}$ en dialektisk Mellembestemmelse af Anlæg til en Linies
og til en Flades Anlæg. I den geometriske Opfattelse af $x^{3}$, hvortil svarer $\int_{0}^{x}3x^{3}dx$, ikke $\int_{0}^{x}(3x^{2}dx+3xdx^{2}+dx^{3})$, % PRINTED AS IS: 3x^{3}dx (for 3x^{2}dx)
sees, hvad $(x+dx)^{3}$ angaaer, ikke blot i $x^{3}$ et Legeme, i
$x^{2}dx$ et Anlæg til et Legeme: den egentlige Flade, i $xdx^{2}$
Anlæg til et legemligt Anlæg: den legemlige Linie, men
endog i $dx^{3}$ et Anlæg til det legemlige Anlægs Anlæg: det
legemlige Punkt.

Den gamle dialektiske Strid angaaende Punktets, Liniens,
Fladens og Legemets indbyrdes Forhold: om det
nemlig er Punktet, der er det Oprindelige, saa at Punktet
frembringer Linien, Linien Fladen, Fladen Legemet; eller
det maaskee omvendt er Legemet, der er det Oprindelige, saa
at Fladen ved Abstraction afledes af Legemet, Linien af
Fladen, Punktet af Linien\footnote[1]{I sin Metaphysik (Metaphys. III. 5) har Aristoteles behandlet det Spørgsmaal, om Tallene, Legemerne, Fladerne og Punkterne ere Væsenheder (οὐσίαι τινές), eller ikke. „I Sammenligning med de Ting, man vel meest maatte tillægge Væsenhed, saasom: Vand, Jord og Luft, hvoraf de sammensatte Legemer bestaae, ere Varme og Kulde og lignende Beskaffenheder (πάθη, Affectioner) jo ikke Væsenheder; kun Legemet, der har disse Affectioner, forbliver som et Værende, en Væsenhed (ὡς ὄν τι καὶ οὐσία τις οὖσα). Og dog synes Legemet at have mindre Væsenhed end Overfladen, denne mindre end Linien, og Linien mindre end Enheden og Punktet. Thi ved disse er Legemet begrændset, og det synes, som kunde disse vel bestaae uden Legemet, men Legemet derimod umuligt uden dem. Indrømmer man imidlertid, at Linierne og Punkterne have mere Væsenhed end Legemerne, kommer man i Forlegenhed; thi disse Grændsebestemmelser beroe jo øiensynlig paa Delinger af Legemet (διαιρέσεις ὄντα τοῦ σώματος) deels i Breden, deels i Dybden, deels i Længden. Desuden maa enten enhver Form indeholdes i Legemet, eller ingen: dersom altsaa Hermes ikke er i Stenen, og Tærningens Halvdeel ikke findes afgrænset i Tærningen selv, saa vil Overfladen heller ikke findes heri; thi dersom nogen Overflade fandtes i Tærningen, maatte ogsaa den findes, hvorved Tærningen halveres. Det Samme gjælder om Linien, Punktet og Enheden. Til de anførte Urimeligheder (ἄλογα) kommer endnu den, der angaaer Opstaaen og Forgaaen; thi snart blive tvende Overflader til een, naar nemlig to Legemer berøre hinanden, snart bliver der igjen, naar et Legeme deles, tvende Overflader, hvor der før ingen var: en tilværende Overflade kan altsaa forsvinde, og en forhen ikke tilværende fremkomme, alt eftersom Legmerne forenes, eller adskilles. Det synes da saaledes, som om Punkter, Linier og Flader, dersom de vare Væsenheder, maatte vexle mellem at existere og ikke at existere; men Punkter, Linier og Flader kunne hverken vorde eller forgaae, ihvorvel de snart ere, snart ikke ere.“ Denne hele Aporie, hvori Aristoteles ikke saameget oplyser Tingen, som sætter den paa Skruer, har her sit Knudepunkt, ikke i det Intellectuelles Forhold til det Materielle, men i det Spørgsmaal: hvad er det Oprindelige, det i sig Gyldige, enten det Abstracte eller det Concrete?}, --- dette Problem finder først
sin Løsning, naar Integrationens Principer alsidig opklares.

% --- p. 109 ---
\setcounter{footnote}{0}%
\opage{109}At det mathematiske Punkt ikke ved nogen Bevægelse
kan frembringe en Linie, eller den mathematiske Linie en
Flade, eller den mathematiske Flade et Legeme, sees deraf,
at Liniens Element $dx$, for at svare til det mathematiske
Punkt, maa være $0$; men naar $dx=0$ ikke kan integreres
til Linie, saa kan $xdx=x.0=0$ heller ikke integreres til
Flade, og $x^{2}dx=x^{2}.0=0$ da heller ikke integreres til Legeme.
I Modsætning til det mathematiske Punkt er altsaa
den endelige Legemsform det Oprindelige. Men denne Legemsform
maa dog, naar vi see paa det i væsentlig Forstand
Oprindelige, ikke destomindre opfattes som det Afledede,
thi Alt, hvad der beroer paa Vorden og Tilbliven,
er jo et Afledet. Den endelige Legemsform har som et Afledet, $x^{3}$, Tilbliven, idet den fremkommer ved Integration
af den legemlige Flade, $x^{2}dx$, denne ved Integration
% the note 1) of p. 108 runs over onto the foot of this page: set in full at its mark
% --- p. 110 ---
\setcounter{footnote}{0}%
\opage{110}af den legemlige Linie, $xdx^{2}$, denne ved Integration af det
legemlige Punkt, $dx^{3}$. Dette legemlige Punkt, Elementernes Element, er altsaa her det ikke Afledede, det i sig Oprindelige,
Legemlighedens Princip, ikke en tom Abstraction,
men et Totalitetens Anlæg\footnote[1]{Naar $dx^{3}$, fastholdt i reen mathematisk Form, ikke desto mindre tager sig ud som en tom Abstraction, hvori man ikke kan søge et Princip, da er Grunden denne, at det saaledes bestemte Punkt, uagtet det er den extensive Størrelses principielle Anlæg, ikke kan faae Realitet, uden med det Samme at betegne et Anlæg til intensiv Størrelse; men Læren om det Intensive falder, forsaavidt den intensive Størrelse altid er mere end Størrelse, udenfor det abstract Mathematiske.}, det Punkt, der \textit{in concreto}
svarer til sit Begreb.

Hvad de enkelte Functioner i Almindelighed angaaer,
kan man nu kjende Integrationens Rigtighed derpaa, at
man ved at differentiere det fundne Integral erholder det
oprindelige Differential, ligesom man, hvad Subtractionen
angaaer, kan kjende Differensen derpaa, at den adderet til
Subtrahenten giver Minuenden. Saaledes er t. Ex.
\[ \int x^{m}dx=\frac{x^{m+1}}{m+1}+c \]
efterdi
\[ d\left(\frac{x^{m+1}}{m+1}+c\right)=x^{m}dx \]
; her er i de
forskjellige Principer væsentlig kun eet Princip. Men, medens
man af en given oprindelig Function altid kan finde
den afledede, er det derimod ikke altid muligt, naar en
Function er given som Differentialcoefficient, at finde en
tilsvarende oprindelig. Vistnok kan man, for at anføre et

et Exempel, omdanne $\int\frac{x^{m}dx}{l(x)^{n}}$ til
\[ -\frac{x^{m+1}}{(n-1)(lx)^{n-1}}+\frac{m+1}{n-1}\int\frac{x^{m}dx}{lx^{n-1}}\ldots \]
, men under fortsat Anvendelse af denne
Formel standser man tilsidst ved $\int\frac{x^{m}dx}{l.x}$, som i Alminde-

% --- p. 111 ---
\setcounter{footnote}{0}%
\opage{111}lighed ($m \lessgtr -1$) ikke kan integreres, fordi der ikke gives
nogen bekjendt Function, hvis Differentialcoefficient er $\frac{x^{m}}{l.x}$\footnote[1]{Jvfr. Jürgensen. Begyndelsesgr. S. 75—77.}.

\begin{center}γ) Det bestemte Integral: Uendelighedsproblemets Løsning.\end{center}

Endskjøndt man ikke kan integrere enhver Function
under endelig Form, kan man dog til enhver Function finde
et Differential, og følgelig betragte enhver endelig Størrelse
som et Integral. Da Integralet nemlig er en Function af
deri indeholdte Variable, $\int_{0}^{x}\Phi(x)dx=f(x)$, kan man i Almindelighed
sætte: $f_{1}(x)=2, f_{2}(x)=3, f_{3}(x)=4\ldots$, og
bestemme $x$ saaledes, at enhver af disse Ligninger, selv om
man kun ved Tilnærmelse kan angive Talværdien af $x$, er
tilfredsstillet, og paa denne Maade erholde: $\int_{0}^{x}\Phi_{1}(x)dx=2$,
\[ \int_{0}^{x}\Phi_{2}(x)dx=3, \int_{0}^{x}\Phi_{3}(x)dx=4\ldots \]

Det articulerende Vexelforhold af Uendeligt med Endeligt
fremlyser af den Maade, hvorpaa det bestemte Integral
angives; thi ihvorvel den endelige Størrelse, hvad Værdien
angaaer, er sig selv lig, ad hvilken Vei vi end komme til
den, saa gjør det dog i Begrebet en uendelig Forskjel, om
en saadan Størrelse haves som umiddelbart given, eller
fremkommer ved Integration.

I Integrationen har Størrelsen sin Vorden; for at
denne Vorden kan ende med Tilværen, maa Integralet
„tages imellem bestemte Grændser.“ Fordringen paa bestemte Grændser betegnes umiddelbart derved, at der til
Differentialfunctionens almindelige Integral føies „en vilkaarlig
Constant.“ Forsaavidt en saadan Constant uden

% --- p. 112 ---
\setcounter{footnote}{0}%
\opage{112}videre tilføies, kan man, naar det kun gjælder om at vise
det heri Tilladelige, jo vistnok beraabe sig paa den Omstændighed,
at, efterdi $d.c=0$, naar $c$ er constant, og
følgelig $d(f(x)+c)=f'(x)dx$, vil man, naar Integrationen
foretages, erholde $\int f'(x)dx=f(x)+c$, og saaledes
faae den Constante med; men Hovedsagen er, at man indseer,
ikke blot det Tilladelige, men det Nødvendige i, at
Constanten tilføies.

I det Endeliges Analyse, hvor hver Størrelse, hvorledes
den end iøvrigt fremkommer, kan tages for sig, opfattes
som tilværende, og forsaavidt tillægges en vilkaarlig
Værdi, kan en særlig Indførelse af vilkaarlige Constanter
ikke spille nogen Hovedrolle;\footnote[1]{I Differensregningen, hvor den ved endelige Mellemrum varierende Størrelse banner Modsætningen til den constante, vil Integralet med Hensyn hertil allerede kræve Tilføining af en Constant. „Ethvert Integral indeholder en Ubestemthed, saalænge den ikke er hævet ved nærmere Betingelser. Dette ligger deri, at $\Delta.(s_{n}+c)=\Delta.s_{n}=u_{n}$, idet c er en vilkaarlig Constant, og altsaa kan man have $s_{n}+c=\Sigma.u_{n}$. Til ethvert Integral kan derfor føies en vilkaarlig Constant, hvis Bestemmelse afhænger af Betingelser, der nærmere maae angives. Derfor er $\Sigma.x^{n}=\frac{x^{n}}{x-1}+c$. Vil man ved $\Sigma.x^{n}$ forstaae Summen af Leddene $1, x, x^{2}:..x^{n-1}$, saa maa til $n=0$ svare $\Sigma.x^{0}=0$, altsaa: $0=\frac{1}{x-1}+c$, hvoraf $c=-\frac{1}{x-1}$, følgelig $\Sigma.x^{n}=\frac{x^{n}-1}{x-1}$, som er bekjendt, og stemmer med Formlen for Summen af en Qvotientrækkes Led.“ A. Steen. Reen Mathem. S. 340—341. Just fordi Differensen hæver Constanten, maa Integralet fordre den tilbage; just fordi $s_{n+1}+c-(s_{n}+c)=s_{n+1}-s_{n}$, maa $\Sigma u_{n}=\Sigma\Delta.s_{n}$ udtrykkelig sættes, ikke som $s_{n}$, men som $s_{n}+c$; thi c, der, idet $c-c=0$, ikke vedkommer $\Delta.s_{n}$, vedkommer derimod, efterdi $s_{n}+c$, for sig taget, tillige afhænger af c, det til $\Delta.s_{n}$ svarende Integral.} i det Uendeliges Analyse
danner den vilkaarlige Constant derimod en charakteristisk
Modsætning til den uendelig variable Størrelse. Udtrykkes

% --- p. 113 ---
\setcounter{footnote}{0}%
\opage{113}t. Ex. den constante Hastighed, $v$, ved $\frac{s}{t}=a$, saa bliver
naturligviis $s=vt=at$. Da $t$ selv i dette Udtryk opfattes som endelig vilkaarlig Størrelse, vilde Tilføiningen af
en vilkaarlig Constant ikke blot være overflødig, men endog
utilladelig; thi dersom $\frac{s}{t}=a$, kan deraf umulig følge

\[ s=at+c \]
med mindre man da ved at sætte $c=0$ gjør
Tilsætningen til et Spilfægteri. Anderledes, naar $s$ skal bestemmes ved $\int adt=at$; thi i denne Fremstilling er $t$ ikke
vilkaarlig, men uendelig variabel, har endnu ingen endelig
bestemt Værdi, er, med andre Ord, ikke i Tilværen, men
i Vorden. At sætte denne Vorden en Grændse ved at angive
et Slutningspunkt, uden tillige at angive et tilsvarende
Begyndelsespunkt, er ligesaa utilstrækkeligt, som det til Angivelsen
af, hvor langt En er reist, er utilstrækkeligt at anføre,
at han er reist til Kjøbenhavn, naar man har glemt
at anføre, hvorfra hans Reise gik ud; Vordens, Forandringens,
i vort Exempel: Bevægelsens Udgangspunkt angives
ved Constanten. For at bestemme $s$ ved $\int adt$, maa Udtrykket
blive: $s=\int adt=at+c$, thi ved at sætte
$t=0$, erholder man da i $c$ den Grændse, hvorfra Bevægelsen
skal regnes. Er $c$ endelig, saa er her en endelig
Afstand imellem et forudgivet Punkt og det for Bevægelsen
fastsatte Begyndelsespunkt; er $c=0$, falde disse to Punkter
sammen.

De for den bestemmende Integration nødvendige Grændsebestemmelser
danne fremdeles en charakteristisk Modsætning
til de Grændsebestemmelser, der fremkomme ved Differentiationen.
For $f(x)=x^{2}$ er $d.f(x)=d.x^{2}=f'(x)dx=2xdx$ en væsentlig Grændse; thi, geometrisk seet, er $xdx$ som Moment en Linie, og Linien er Fladens Grændse: den
Grændse, hvortil Differentiationen fører, er, som vi have
seet, i Forhold til det Begrændsede et Moment, et Anlæg,

% --- p. 114 ---
\setcounter{footnote}{0}%
\opage{114}et Element. Den Grændse, hvorved Integralet bestemmes,
er derimod, naar den da ikke fastholdes i den rene Abstraction
som $0$ eller $\infty$, aldrig et blot Moment eller en blot
Mulighed, men altid en tilværende endelig Størrelse, og i
denne sin Endelighed sideordnet med det derved bestemte
Integral. Saaledes idet $\int dx=x+c$, $\int 2x\,dx=x^{2}+c$, $\int 3x^{2}dx=x^{3}+c$, er $c$ en Linie, naar $x$ er
en Linie, et Areal, naar $x^{2}$ er et Areal, et Volumen, naar $x^{3}$ er et Volumen. Den tilføiede Constante, som, hvad
Værdien angaaer, er vilkaarlig, er altsaa hvad Størrelsens
Art og Beskaffenhed angaaer, ingenlunde vilkaarlig. At
Constanten maa være eensartet med Integralet, sees udtrykkelig,
naar den opfattes som en Function af Integralets
Variable, hørende med til Integralets Bestemthed, og derved
viser sig som væsentlig Deel af Integralet selv. Betegnes $\int 2x\,dx$ ved $\int f'(x)dx+c=f(x=b)-f(x=a)$
\[ =\int_{x=a}^{x=b}2xdx=b^{2}-a^{2} \]
; da har $c=-f(x=a)=-f(a)$
opgivet sin Umiddelbarhed, og er som en Functionens
Grændse forvandlet til en reflecteret Grændsebestemmelse.

Er nu Grændsebestemmelsen reflecteret i Functionen, saa
indsees tillige heraf, at man, istedetfor at sætte
\[ \int_{x=a}^{x=b}2\,x\,dx \]

\[ =\int 2\,b\,dx-\int 2a\,dx \]
, hvad her, eftersom $x$
er variabel, vilde
være aldeles meningsløst, først maa udføre Integrationen,
for derefter i Udtrykket: $\int 2x\,dx=x^{2}$ at kunne tillægge $x$

de vilkaarlige Værdier: $b$ og $a$. Forsaavidt den endelig
bestemte Størrelse da paa denne Maade fremstilles som et
bestemt Integral, er Vexelbestemmelsen af Endeligt og Uendeligt
articuleret i objectiv Reflexion. I og med enhver endelig
Størrelse bliver altsaa det Uendelige sat; thi enhver endelig

% --- p. 115 ---
\setcounter{footnote}{0}%
\opage{115}Størrelse er, naar den sees i Begrebet, en Determination
af det Uendelige, et definitivt Udtryk for $\infty. 0.$

Idet Integralet bestemmes saaledes, at de deri reflecterede
Grændser selv maae betragtes som Integraler, kan
Uendelighedsproblemet siges at være løst; thi medens Differentialet
viser det uendeliges Oprindelse af det Endelige,
viser Integralet omvendt det Endeliges Oprindelse af det
Uendelige. Kun det Uendeliges Analyse, saaledes som Mathematiken
gjennemfører den, kan altsaa lede til Problemets
Løsning; men Løsningen selv er, væsentlig seet, ikke
mathematisk: det er ikke en vanskelig Beregning, ikke en indviklet
Operationsmechanik, men en dialektisk Begrebsudvikling,
hvorpaa det i sidste Instans kommer an.

\begin{center}Slutning.\end{center}

Kan nu den i dette Skrift angivne Løsning af Uendelighedsproblemet
i nogen Grad bestaae sin videnskabelige
Prøve, da vil det heller ikke være vanskeligt at godtgjøre,
at en Dialektik, hvori Analysens exacte Elementer indgaae,
er for Philosophien af væsentlig formel, og just derved af
sand og reel Interesse.

Allerede i den ældre græske Philosophie opkomme, navnlig
hos Eleaterne, saadanne dialektiske Qvantitetsproblemer,
som, efterat de igjennem Aarhundreder ere bearbeidede selv
af de meest fremragende Tænkere, dog først i den nærværende
Tid, ved nemlig at sees under den høiere Analyses
Synspunkt, kunne finde deres fuldstændige Løsning. Idet
vi i denne Sammenhæng særlig fremhæve Delings- og
Bevægelsesproblemet, kunne vi slutte med at minde om,
hvorledes de forskjellige herhen hørende Overveielser alle tilsidst
indgaae under det ene omfattende, for Philosophien først
% --- p. 116 ---
\setcounter{footnote}{0}%
\opage{116}og sidst afgjørende Problem, nemlig: Idealitetspro d
blemet.

\begin{center}a) Delingsproblemet.\end{center}

De bekjendte Indsigelser, Eleaten Zeno har reist imod
Muligheden af, at en Fleerhed i det Værende skulde lade
sig tænke, har Zeller\footnote[1]{Die Philosophie der Griechen. 1ster Th. Tübingen 1856. S. 425—28.} sammendraget og fremstillet i følgende
Vendinger:

„Dersom det Værende var et Mangfoldigt, saa maatte
det paa een Gang være baade uendelig lidet og uendelig
stort. Uendelig lidet, thi da enhver Fleerhed er et Antal
af Eenheder, men kun det Udeelbare er en virkelig Eenhed,
saa maa enhver af de Mange enten selv være en udeelbar
Eenhed, eller bestaae af saadanne Eenheder. Men hvad
der er udeelbart, kan ikke have nogen Størrelse, thi alt, hvad
der har en Størrelse, er deelbart i det Uendelige. De enkelte
Dele, hvoraf det Mangfoldige bestaaer, have altsaa
ingen Størrelse. Der er altsaa Intet, der kan blive større
derved, at slige Dele tilføies det, og Intet, der kan blive
mindre derved, at de fradrages det. Men hvad der, naar
det føies til et Andet, ikke forstørrer det, og derved at det
fradrages, ikke formindsker det, det er Intet. Det Mangfoldige
er altsaa uendelig lidet, thi enhver af dets Bestanddele
er saa liden, at den er et Intet. Men paa den anden
Side maae disse Dele ogsaa være \emph{uendelig store.} Thi da
det, som ingen Størrelse har, ikke er, saa maae de Mange,
for at være, have en Størrelse, deres Dele maae altsaa
være adskilte fra hverandre. Ligeledes maa ogsaa det, der
fjerner dem fra hverandre, have en Størrelse, og det, som
igjen fjerner dette fra dem, ligeledes, og saa fremdeles i
det Uendelige, saa at vi ifølge heraf erholde uendelig mange
Størrelser eller en Uendelighed af Størrelser“. „Ved en
lignende Betragtningsmaade viser Zeno, at det Mangfol-

% --- p. 117 ---
\setcounter{footnote}{0}%
\opage{117}dige, hvad Tallet angaaer, maa være baade begrændset og
ubegrændset\dots{}. Ligesom ved det første Beviis Bestemmelsen
af den uendelige Størrelse, saaledes bliver her Bestemmelsen
af det uendelige Tal vundet derved, at Mangfoldigheden
opfattes som en Fleerhed af adskilte Størrelser, og at der
mellem hvert Par Adskilte indskydes et tredie Adskillende.
De Gamle pleiede desaarsag at betegne denne Deel af begge
Beviser som Beviset af Tvedelingen.“

Pointen er her opfattet saa skarpt og Knuden knyttet
saa fast, at hverken en Platos Dialektik eller en Aristoteles's
Analytik var istand til at løse Problemet. Hvad der hos
Grækerne maatte endog i almeenlogisk Henseende vanskeliggjøre
Løsningen af Uendelighedsproblemet, var den bekjendte
gjennem hele den græske Philosophie gaaende Adskillelse af
det Begrændsede (πέρας, πεπερασμένον) og det Ubegrændsede
(ἄπειρον), en Adskillelse, ifølge hvilken det Begrændsede var
Eet med det Fornuftige, det Ubegrændsede derimod med det
Ufornuftige, det Ubestemmelige, det i sig Utænkelige (ἄλογον).\footnote[1]{Den aristoteliske Synsmaade, der (jvfr. t. Ex. Phys. III. V. VI.) i det Hele taget løber ud paa, at det Uendelige ikke er i Virkeligheden, men kun efter Muligheden, er en Fornegtelse af det Uendeliges Realitet: Ὅτι δὲ εἰ μὴ ἔστιν ἄπειρον ἁπλῶς, πολλὰ ἀδύνατα συμβαίνει, δῆλον. τοῦ τε γὰρ χρόνου ἔσται τις ἀρχὴ καὶ τελευτή. καὶ τὰ μεγέθη οὐ διαιρετὰ εἰς μέγεθος, καὶ ἀριθμὸς οὐκ ἔσται ἄπειρος .... καὶ δῆλον ὅτι πῶς μὲν ἐστι, πῶς δ’ οὔ. Λέγεται δὲ τὸ εἶναι, τὸ μὲν δυνάμει, τὸ δὲ ἐντελεχείᾳ... Λείπεται οὖν δυνάμει εἶναι τὸ ἄπειρον. Phys. III. 6. Idet det Uendelige paa denne Maade holdes i den blotte Mulighed, altsaa i en Mulighed, der aldrig kan blive til Virkelighed, bliver det et Umuligt. Det Ubegrændsede, som hos Grækerne betegnede det Laveste, Materien, forvandledes i den middelalderlige Speculation til det Høieste, det Uendelige, som da igjen blev, om end ikke i orientalsk Forstand, et Udtryk for den fuldkomne, guddommelige Væren. Som Den, der har bragt den absolut uendelige Eenhed i Forhold til Tallet og det Mangfoldige, maae vi her nævne Cardinalen Nicolaus af Cusa († 1464); han har i sine tre Bøger \textit{De docta ignorantia} angivet følgende Synsmaade: „Den i Sandhed uende- lige Væren kan ikke have nogen anden Væren udenfor sig, der var istand til at danne en virkelig Modsætning, thi derved vilde den fuldkomne Væren blive begrændset; den uendelige Væren er den absolute Selvlighed (absoluta identitas), alle Formers Form, alle Muligheders Virkelighed. Da nu det absolut Største er Eet, saa kan det, der er istand til et Mere eller et Mindre kun være en Fleerhed, altsaa kun være i Tallet. Saaledes er den uendelige Eenhed som det absolut Mindste alle Tals Grund og Begyndelse, og som det absolut Største alle Tals Ende og Grændse.“ Maximum, quo majus esse nequit, simpliciter et absolute, cum majus sit, quam comprehendi per nos possit, quia est veritas infinita: non aliter quam incomprehensibiliter attingimus. Nam cum non sit de natura eorum, quæ excedens admittunt et excessum, super omne id est, quod per nos concipi potest: omnia enim quæcunque sensu, ratione aut intellectu apprehenduntur, inter se et ad invicem taliter differunt, quod nulla est æqualitas præcisa inter illa. Excedit igitur maxima æqualitas, quæ a nullo est alia aut diversa, omnem intellectum: quare maximum absolute, cum sit omne id quod esse potest, est penitus in actu, et sicut non potest majus esse, eadem ratione nec minus, cum sit omne id, quod esse potest. Minimum autem est, quo minus esse non potest. Et quoniam maximum est hujusmodi, manifestum est minimum maximo coincidere. Et hoc tibi clarius fit, si ad quantitatem maximum et minimum contrahis. Maxima enim quantitas est maxime magna. Minima quantitas est maxime parva. Absolve igitur a quantitate maximum et minimum, subtrahendo intellectualiter magnum et parvum: \& clare conspicis, maximum et minimum coincidere. D. Nicolai de Cusa Cardinalis opera. Basiliæ. De docta ignoratia. Lib. I. cap. IV. Denne Identitet af et Begrændset-Ubegrændset, et uendeligt og dog absolut Maximum, et uendeligt og dog absolut Minimum blev i det 16de Aarhundrede videre udviklet af G. Bruno.}

% --- p. 118 ---
\setcounter{footnote}{0}%
\opage{118}Med Newtons Fluxionsmethode og Leibnitzes Differen
tialregning gjør den nyere Tid forsaavidt et i Forhold til
Oldtiden betydeligt Fremskridt, som Hensynet til det Uendelige,
hvorledes man saa end for Resten opfatter det, fra nu
af viser sig stedse mere paatrængende. Idet Delingsproblemet
gik over i den høiere Analyse, begyndte det Dialektiske
at røre sig i den exacte Videnskab; hvor betegnende er ikke
% the note 1) of p. 117 runs over onto the foot of this page: set in full at its mark
% --- p. 119 ---
\setcounter{footnote}{0}%
\opage{119}i denne Henseende Saveriens af Maclaurin anførte Ord:
»\textit{on vit alors pour la premiere fois des disputes parmi
les Géometres}«! Mathematikeren driver de exacte Analyser
til de yderste Grændser, og seer saa til sin ubehagelige
Overraskelse Dialektiken bryde frem.\footnote[1]{Nous supposons que la soudivision d'un Etre peut se continuer a l'infini, \& dès que cette soudivision devient telle, que nous ne pouvons plus la saisir, nous regardons comme nul, le reste de cet Etre infiniment petit, quoique nous voulions qu' il y ait encore dans lui une division infinie, \& que nous égalons néanmoins à zero. Cela est admirable: nous voulons concevoir l'infini, \& connoître cet infini par le fini même. Cette contrarieté vient de ce que nous ne savons au juste ce que c'est que l'infini. Dire qu'une chose est infinie, c'est dire que nous n'y connoissons pas de terme; et dès-lors nous devons rejetter tout ce qui sembleroit supposer en nous cette connoissance. Saverien Histoire du Calc. p. XVII.—XIX. Selv Lacroix, der dog i sin \textit{Traité du Calcul différentiel et intégral. Tom. I. Paris 1810} ikke indlader sig videre paa, hvad han kalder \textit{le changement de métaphysique,} kan dog alligevel ikke afholde sig fra Sidehug til dette brillende Begreb, dette lille dialektiske Udtryk, \textit{cet être infiniment petit!} Saaledes i anf. Skr. S. 242. En partageant avec toute l'Europe, le respect attaché au nom et aux travaux de M. Lagrange, j'oserai néanmoins n'être pas de l'avis de cet homme si justement célèbre, sur les motifs qui paraissent le porter à introduire une nouvelle manière d'écrire les résultats du calcul différentiel; car je crois qu'il est facile de prouver que l'ancienne n'a point en elle même l'inconvénient de rappeler continuellement l'idée fausse des infiniment petits.} Det 17de og tildeels
det 18de Aarhundredes Philosophie var imidlertid, forsaavidt
den da ikke var aldeles skeptisk, saa fixeret i Dogmatisme, at
Bevidstheden om det Dialektiske, endog hos Tænkere som
Spinoza og Leibniz, blev aldeles tilbagetrængt. Intet Under
da, at Begreberne for en Spinoza chrystalliserede i exacte
Definitioner\footnote[2]{Jvfr. Benedict Spinoza, af Dr. H. Brøchner. Kbhn. 1857. S. 44—58 o. fl. Stdr.} intet Under, at en Leibniz med alt sit Genie

% --- p. 120 ---
\setcounter{footnote}{0}%
\opage{120}ikke var istand til at begribe som Philosoph, hvad han selv
opdagede som Mathematiker; intet Under, at en Wolf kunde
hylde den mathematiske Methode som Forbillede for den philosophiske,
medens en Berkeley\footnote[1]{A. Treatise concerning the Principles of Human Knowledge. London 1734. S. 146. He whose Understanding is prepossest with the Doctrine of \emph{abstract} general Ideas may be persuaded, that (whatever be thought of the Ideas of Sense) Extension in abstract is infinitely divisible. And one who thinks the Objects of Sense exist without the Mind, will perhaps in virtue therefore be brought to admit, that a Line but an Inch long may contain innumerable Parts really existing, though too small to be discerned. These Errors are grafted as well in the Minds of Geometricians, as of other Men, and have a like influence on their Reasonings .... S. 150: Of late the Speculations about Infinites have run so high, and grown to such strange Notions, as have occasioned no small Scruples and Disputes among the Geometers of the present Age. If it be said that several Theoremes undoubtedly true, are discovered by Methods in which Infinitesimals are made use of, which could never have been, if their Existence included a Contradiction in it. I answer, that upon a thorough Examination it will not be found, that in any Instance it is necessary to make use of or conceive infinitesimal Parts of finit Lines, or even Quantities less than the Minimum Sensibile: Nay, it will be evident this is never done, it being impossible. ... Particularly, \emph{Matter or the absolute Existence of Corporeal Objects,} hath been shewn to be that wherein the most avowed and pernicious Enemies of all Knowlegde, whetever humane or divine, have ever placed their chief Strength and Confidence. S. 152 o. flg.} netop i den Modsigelse, der
ligger i den uendelige Deelbarhed, meente at see en bundløs
Urimelighed, et Middel til at bortforklare Materiens Realitet.

Uagtet den dogmatiske og skeptiske Philosophie imod
Slutningen af det 18de Aarhundrede mere og mere fortrængtes
af den kritiske, var dog i Henseende til Problemets
Løsning saa godt som Intet vundet, da Tænkningen, trods dens
forandrede Stilling til Væren, endnu blev staaende ved det

% --- p. 121 ---
\setcounter{footnote}{0}%
\opage{121}formelle Contradictionsprincip. Iblandt Kants kosmologiske
Antinomier spiller naturligviis ogsaa Delingsproblemet sin
Rolle som uopløseligt; ligesom det, hvad disse Antinomier
angaaer, i det Hele er charakteristisk, at den Indsigt, Kant
som Mathematiker unegtelig har i Forholdet mellem det
Endelige og det Uendelige, aldeles ikke er ham nærværende,
og altsaa ikke kommer ham tilgode, naar han philosopherer.
Først i den ved Fichte forberedte, af Hegel udviklede Dialektik
begynder et nyt Lys at opgaae over det Uendelige\footnote[1]{See H. Schwarz. Anf. Skrift.}

Hvad nu Delingsproblemet særlig angaaer, maae vi
vel, for at fatte dets fuldstændige Løsning, helst opfatte det fra
to Hovedsider, nemlig den blot formelle og den substantielle

Side; thi Sammensmeltningen af disse to Synspunkter,
saaledes som den allerede findes hos Zeno, har vistnok ikke
bidraget lidet til at forvirre Betragtningen og vanskeliggjøre
Opgavens Løsning.

Zenos Forudsætning, at det Værende, dersom det
Mangfoldige er, maa have en Størrelse, eller, at „det, som
ingen Størrelse har, ikke er“, fordrer da for det Første en
Abstraheren fra det Substantielle, en Bortseen fra det Værendes
Realitet. At være til er at være Størrelse, og
Modsigelsen i det, at være til, en Modsigelse i det, at
være Størrelse: her er Problemet i dets abstracte Form.

Modsigelsen i Størrelsens Begreb er Deelbarheden.
Denne Modsigelse opløser sig i den Indsigt, at Deelbarheden
ved at fortsættes i det Uendelige aldeles ikke medfører
nogen Tilintetgjørelse for nogensomhelst Deel af den deelte
Størrelse. Ligesom $1$ ikke taber Noget ved at fremstilles
som $\frac{1}{2}+\frac{1}{2}$ eller som $2.\frac{1}{2}$, saaledes taber en Linie heller ikke

det Allerringeste af sin Størrelse som et Hele derved, at den
tænkes deelt i to ligestore Dele, enhver af disse Dele igjen
i to Dele o. s. fr. Lade vi Rækken: $1+\frac{1}{2}+\left(\frac{1}{2}\right)^{2}....$

% --- p. 122 ---
\setcounter{footnote}{0}%
\opage{122}$+\left(\frac{1}{2}\right)^{n}+\ldots$ repræsentere en i det uendelige gaaende
Dichotomie, da vil Udtrykket $\mathrm{Lim}\left(\frac{1}{2}\right)^{n}$ betegne et Element,
en uendelig lille Størrelse, men ikke det absolute Nul. Ved
en i det Uendelige fortsat Deling kan det absolute Nul ikke
fremkomme; skulde en fortsat Deling føre til Nul, maatte
den jo baade medføre en fortsat Tilintetgjørelse og tillige
standse ved en sidste Grændse, men begge Antagelser ere i
lige Grad stridende mod den uendelige Deelbarhed og mod
Størrelsens Væsen.\footnote[1]{At Nul ikke kan naaes ved fortsat Deling, følger ogsaa deraf, at enhver ny Deling forudsætter, at der sættes en ny Grændse; men den nye Grændse kan ikke sættes umiddelbart sammen med den allerede satte; er her ikke et, om end nok saa lille, Mellemrum mellem Grændserne, falde de sammen til een Grændse, og der skeer ingen ny Deling. Naar Grændserne rykke hinanden nærmere og nærmere, maa det Begrændsede unegtelig blive mindre og mindre; men de sig hinanden nærmende Grændser tilintetgjøre dog alligevel ikke Delene, de bevæge sig kun hen over Delene. Er saaledes $y=f(x)$ Ligningen for en Curve, og $z=\Phi(x)$ Udtrykket for et Areal begrændset af Abscisseaxen, en Deel af Curven og to Ordinater, hvoraf den fjerneste svarer til Abscissen $x$; da er $\Phi(x+h)-\Phi(x)$ mindre end det Rectangel, der svarer til den største Ordinat og Abscissetilvæxten $h$, men større end det Rectangel, der svarer til den mindste Ordinat og $h$. Altsaa: $h f(x) \lessgtr \Phi(x+h)-\Phi(x) \lessgtr h.f(x+h)$; $h f(x) \lessgtr \Phi'(x) h+\Phi''(x)\frac{h^{2}}{1.2}+\ldots \lessgtr h f(x)+h^{2} f'(x)+\ldots$, og $f(x) \lessgtr \Phi'(x)+\Phi''(x)\frac{h}{1.2}+\ldots \lessgtr f(x)+f'(x)\frac{h}{1}+\ldots$ Gaaer man nu til Grændsen ved at sætte $h=0$, d. v. s. lader man de to Ordinater nærme sig hinanden, indtil de falde sammen, blive ovenstaaende Grændser til Eet med den imellem dem liggende Størrelse, d. v. s. $\Phi'(x)=f(x)$. Men idet Grænd- serne falde sammen, og det af dem indesluttede Areal indsvinder til $0$, gaaer dog alligevel Intet tabt, ingen Deel af Størrelsen bliver tilintetgjort; kun den Forandring foregaaer, at de Arealets Dele, der før faldt indenfor de afskilte, nu falde udenfor de sammenfaldne Grændser. Forandring af Grændser tilveiebringer vistnok altid Forandring i Størrelser og Størrelsers Forhold; men ingen Grændseforandring, ingen Omflytning, Udvidelse, Nærmelse eller Sammenfalden af Størrelsesgrændser formaaer at gjøre nogensomhelst allerede tilværende Størrelse til Intet.}

% --- p. 123 ---
\setcounter{footnote}{0}%
\opage{123}Den Svimmelhed, der saa let griber En ved Forestillingen
om en i det uendelige sig fortsættende Deling,
vedkommer egentlig ikke Sagen; det er en Aandsforvirring,
som hidrører derfra, at man istedetfor at tænke, giver sig
til at phantasere, og vil phantaserende forestille sig det, der
ikke lader sig forestille. Den fortsatte Deling er, naar vi
holde Tanken fast, kun en Gjentagelse af den første Deling:
i den første Deling ligger hele Problemet. Tænkes en endelig
Linie $x$ deelt ved et eneste Punkt i to Linier, da er Liniens
Element $dx$ ved det angivne Punkt med det Samme
deelt i to Elementer: at dele en Linie uden at dele en af
dens uendelig smaae Dele, er umuligt. Kan nu Linien $x$ deles, uden at Noget gaaer tabt eller tilintetgjøres, da gjælder
det Samme om $dx$; men er $dx$ deelt, uden at Noget er
tilintetgjort, da er jo eo ipso $\frac{dx}{2},\ \frac{dx}{3}\ldots\frac{dx}{n}$ ogsaa deelt,
uden at nogen Deel er tilintetgjort.
At dele Linien $x$ er altsaa med andre Ord at dele Elementet $dx$, men at dele
Elementet $dx$ er igjen at dele et Element af Elementet, at
dele $\frac{dx}{n}$; af det deelte $\frac{dx}{n}$ vil da $\mathrm{Lim}\,\pm\frac{dx}{n}$ betegne de Mo-
menter, som paa begge Sider grændse op til Grændsen, og
saaledes ved uendelig Tilnærmelse gaae over i det $+0$ og $-0$, hvorfra det rene $0$ dialektisk adskiller sig.
% the note 1) of p. 122 runs over onto the foot of this page: set in full at its mark

% --- p. 124 ---
\setcounter{footnote}{0}%
\opage{124}Lægges derhos, efter at de formelle Vanskeligheder ere
hævede, særligt Eftertryk paa den tilværende Størrelses
Realitet, da viser Problemet sig fra den substantielle
Side. Zenos Ord: „Kun det Udeelbare er en virkelig
Eenhed, enhver af de Mange maae altsaa enten selv være
en udeelbar Eenhed eller bestaae af saadanne Eenheder“,
angaae egentlig ikke Størrelsen blot som Størrelse, men det
i Størrelsen Reelle, Substansen. Den virkelige Eenhed
lider da nu ikke nærmest af den Modsigelse at være deelbar
i det uendelige; men Modsigelsen er, at den egentlige Substans
skal være baade deelbar og udeelbar. Forsaavidt
Størrelsen er substantiel, er dens Eenhed udeelbar; forsaavidt
Substansen er Størrelse, er dens Eenhed deelbar:
denne Modsigelse stod klar for Leibniz, da han, for at
undgaae den, udfandt Monadelæren\footnote[1]{Log. et Metaph. p. 5. Necesse autem est dari substantias simplices (vel monades, nam monas non est nisi substantiá simplex), quia dantur composita. Jvfr. „En Fremstilling af Leibniz's Monadelære.“ C. F. Nielsen. Kbhvn 1856.} Thi hvad er vel
en Monade? En Enkeltsubstans, en substantiel Størrelse,
en Størrelse, der er mere end en Størrelse, en oprindelig
Eenhed, der ikke kan deles. For at adskille Monaderne fra
de mathematiske og physiske Eenheder valgte Leibniz at
kalde dem substantielle Eenheder, metaphysiske Punkter. »\textit{On
les pourroit appeller points métaphysiques\dots{}; ainsi les
points physiques ne sont indivisibles qu'en apparence:
les points mathématiques sont exacts, mais ce ne sont
que des modalités: il n'y a que les points métaphysiques
ou de substance, qui soient exacts et réels; et sans eux
il n'y auroit rien de réel, puisque sans les véritables
unités il n'y auroit point de multitude}«. \textit{Log. et Metph.
p. 53.}

I Monaderne mener altsaa Leibniz at have de existerende
Eenheder, som uagtet de ved den dem iboende Kraft

% --- p. 125 ---
\setcounter{footnote}{0}%
\opage{125}bevirke, at der danne sig physiske og deelbare Størrelser, dog i
deres substantielle Enkelthed ikke desto mindre ere absolut
udeelbare. Det Spørgsmaal, hvorledes Det, som i sig ikke
har nogen Størrelse, kan være Princip for virkelige Størrelser, hvorledes Det, som selv er uden Udstrækning, kan
frembringe det Udstrakte, bliver da besvaret derhen, at Monaderne
som Enkeltsubstanser ere individuelle, saa at enhver
Monade har ifølge sin individuelle Begrændsning en oprindelig
virksom \textit{(force active primitive)} og en oprindelig
lidende Kraft \textit{(force passive primitive)}. Som Materiens
Princip har den lidende Kraft en Tendens til Udstrækning,
og foraarsager derved, at den oprindelige Materie \textit{(materia
prima)}, der ikke bestaaer \textit{»in extensione«}, men \textit{»in extensionis
exigentia«}, bliver en udstrakt og derved uendelig
deelbar Størrelse.\footnote[1]{Il est vrai (selon mon système) qu'il n'y a point de portion de la matiere, ou il n'y ait une infinité de corps organiques .... Il n'y a point de partie de la matiére, qui ne soit divisée actuellement. Log. et Metaph. p. 39, 44 o. fl.}

Det i denne Forklaring Utilfredsstillende er, saalænge
det Dialektiske oversees, for iøinefaldende til, at Modsigelsen
ikke skulde have fundet sit tilsvarende Udtryk i den kritiske
Philosophie. Under „ „den rene Fornufts Antinomier“
fremstiller Kant da ogsaa som „de transcendentale Ideers
anden Strid“ Striden om det Udeelbare- Deelbare. Thesis:
\emph{Enhver sammensat Substans i Verden bestaaer
af enkelte Dele, og der existerer overalt Intet,
uden det Enkelte, eller det, som deraf er sammensat.}
Antithesis: \emph{Ingen sammensat Ting i Verden
bestaaer af enkelte Dele, og der existerer overalt
intet Enkelt i Verden.}

„Antager man“, saaledes bevises Thesis, „at de
sammensatte Substanser ikke bestode af enkelte Dele, saa
vilde der, naar alle Sammensætninger tænkes opløste, ikke

% --- p. 126 ---
\setcounter{footnote}{0}%
\opage{126}blive nogen sammensat Deel, og (da der ingen enkelte Dele
gives) heller ingen enkelt tilbage: man fik da tilsidst kun det
rene Intet, og der kunde følgelig ikke gives nogen Substans.“
Beviset for Antithesis hviler paa følgende Grundtanke:
„Den sammensatte Ting (Substans) maa være i Rummet;
men Rummet bestaaer ikke af enkelte Dele, men af
mindre Rum; altsaa maa enhver Deel af det Sammensatte
indtage et Rum. I alt Sammensat ere de egentlig første
Dele enkelte; altsaa indtager det Enkelte et Rum. Da nn
nu ethvert Reale, der indtager et Rum, omfatter en Mangfoldighed
af udenfor hinanden værende Dele, og følgelig er
sammensat, og det som et reelt Sammensat ikke af Accidenser
(thi Accidenserne maae, for at være udenfor hinanden, forudsætte
Substansen), men af Substanser: saa maatte det
Enkelte, forsaavidt det skulde være i Rummet, være et substantielt
Sammensat; hvilket er en Modsigelse.“\footnote[1]{At Kant i denne Antithetik blot opfatter det Enkelte som Forudsætning for og Modsætning til det Sammensatte, og forsaavidt kun seer den Leibniziske Enkeltsubstans fra dens ene Side, fremhæver han selv udtrykkelig i sin „Anmerkung zur zweiten Antinomie.“ (Krit. d. r. V. S. 356) „Ich rede übrigens hier nur von dem Einfachen, so fern es nothwendig im Zusammengesetzten gegeben ist, indem dieses darin, als in seine Bestandtheile aufgelöset werden kann. Die eigentliche Bedeutung des Wortes \emph{Monas} (nach Leibnizens Gebrauch) sollte wohl nur auf das Einfache gehen, welches \emph{unmittelbar} als einfache Substanz gegeben ist (z. B. im Selbstbewußtseyn) und nicht als Element des Zusammengesetzten, welches man besser den Atomus nennen könnte. Und da ich nur in Ansehung des Zusammengesetzten die einfachen Substanzen, als deren Elemente beweisen will, so könnte ich die Antithese der zweiten Antinomie die transcendentale \emph{Atomistik} nennen. Weil aber dieses Wort schon vorlängst zur Bezeichnung einer besonderen Erklärungsart körperlicher Erscheinungen (molecularum) gebraucht worden, und also empirische Begriffe voraussetzt, so mag er der dialektische Grundsatz der \emph{Monadologie} heißen.“}

% --- p. 127 ---
\setcounter{footnote}{0}%
\opage{127}Spørgsmaalet om det Enkeltes Forhold til det Sammensatte
viser, hvad enten man saa om det Enkelte bruger
Udtrykket „Atomus“ eller „ Monas“, oprindelig hen til
Spørgsmaalet om Materiens Forhold til Kraften, og derved
til Spørgsmaalet om den extensive Størrelses Forhold
til den intensive. Det er den rationelle Mechanik, som ved
Anvendelse af det uendeliges Analyse fra sit exacte Standpunkt
udbreder et høiere Lys over de omspurgte Grundforhold,
og derved giver en for Problemets Løsning uundværlig
Forudsætning; men Løsningen selv er dialektisk.

Ved at fastholde den kritiske Grændse lader den exacte
Videnskab overalt saadanne modsatte Begreber som: Materie
og Kraft, det Intensive og det Extensive, det Enkelte og
det Sammensatte, udelukke hinanden; hvad der er Materie
er altsaa ikke Kraft, det Extensive er ikke intensivt, det Enkelte
ikke sammensat. At Kraften væsentlig er Materie og
Materien Kraft, at det Extensive i sin Realitet dog alligevel
er intensivt, og det Intensive extensivt, at det Enkelte i
Grunden kan være et Sammensat, og det Sammensatte et
Enkelt: slige Sætninger maa den rationelle Mechanik paa
det Bestemteste afvise, just fordi den er en exact, ikke en
dialektisk Videnskab.

Men netop ved sine exacte Analyser kan den rationelle
Mechanik paa en ganske eiendommelig Maade siges at forberede
og fremkalde en videnskabelig Indsigt i det Dialektiske.
Den rationelle Mechanik bliver nemlig ikke staaende
ved den Kjendsgjerning, at de endelige Masser paa endelig
Maade tiltrække hinanden, men tyder Tiltrækningen saaledes,
at hver mindste Deel i hver Masse maa, idet den tiltrækkende
tiltrækkes, paa een Gang være baade Udgangspunkt
og Angrebspunkt \textit{(point d'application)} for en virkende Kraft.
Forsaavidt her nu ikke handles om dette eller hiint Virkelige,
der altid er disse eller hine særegne Betingelser underlagt,
men om Phænomenernes Væsen, det Virkeliges Princip,
bliver det Endelige conseqvent opløst i det uendelige
% --- p. 128 ---
\setcounter{footnote}{0}%
\opage{128}som i sin Grund, og dette Uendelige saa igjen opfattet som
den reale Mulighed, der altid paa givne Betingelser virkeliggjøres
i et tilsvarende endeligt Udtryk.\footnote[1]{Il est évident, ... que la division de la masse en élémens infiniment petits, et la supposition d'une densité de chaque élément, qui ne varie pas dans les corps homogènes, ou qui varie par degrés insensibles dans les corps hétérogènes, ne conviennent point aux corps naturels; mais cela n'empêche pas qu'on ne puisse faire usage des formules fondées sur cette considération ... Poisson Traité de Méchanique. Tom. I. Paris 1833. S. 175. At Poisson ved \textit{»corps naturels«} kun tænker paa Tingene i deres Endelighed, ikke paa Tingenes Væsen, sees ogsaa af et Udtryk som dette: \emph{Un infiniment petit} est une grandeur moindre que toute grandeur donnée de la même nature. Anf. Skr. S. 14.}

Indsvinder, for at henvise til et Exempel, det materielle
Volumen $\varrho v = M$ til et extensivt Element $dM$, saa indsvinder
Kraften (Vægten) $P$ med det Samme til et intensivt
Element $dP$; en endelig Kraft (Vægt) er altsaa kun
mulig ved Massens tilsvarende endelige Volumen. Betegner $p$ Vægten af en Volumenenhed, da bliver Vægten af et
Volumenelement, udtrykt ved polære Coordinater, $d.P = p r^{2}\,dr\,\sin\alpha\,d\alpha\,d\Theta$. Integreres nu denne Ligning, faaes
endelig Vægt gjennem endeligt Volumen. Er det t. Ex. en
homogen Kugle, saa at $p$ altsaa er constant, erholder man,
for $r=a$, ved Integration:
\[ P=p\int_{0}^{a}\int_{0}^{\pi}\int_{0}^{2\pi} r^{2}\,dr\,\sin\alpha\,d\alpha\,d\Theta \]
\[ =2p\pi\int_{0}^{a}\int_{0}^{\pi} r^{2}\,dr\,\sin\alpha\,d\alpha=4p\pi\int_{0}^{\alpha} r^{2}\,dr=\frac{4}{3}p\pi a^{3}. \]

Ved enhver endog den allermindste Variation af $a$, maa, idet $p$ er constant, en tilsvarende Variation finde
Sted i $P$, og omvendt; d. e. enhver endog den allerringeste
Forandring i det Materielle, det Extensive, er med det Samme
en Forandring i Kraften, det Intensive.

% --- p. 129 ---
\setcounter{footnote}{0}%
\opage{129}Af den herved antydede Analyse fremgaaer det da som
uimodsigeligt Resultat, at Materien i det Uendelige forudsætter
Kraften, ligesom Kraften igjen i det uendelige forudsætter
Materien, at Materien altsaa ligesaa lidt kan være
uden Kraft, som Kraften uden Materie, at Materiens og
Kraftens Tilværen væsentlig er een og samme Væren. Da
Kraften kun er mulig formedelst Materien, er det umuligt
at aflede Materien af Kraften; da Materien omvendt kun
er mulig formedelst Kraften, er det umuligt at aflede Kraften
af Materien; den \emph{dynamiske} Anskuelse, der gjør Kraften
til Princip, er ligesaa uholdbar, som den \emph{materialistiske,}
der gjør Materien til Princip: Principet er hverken Kraften
som Kraft eller Materien som Materie, men Materiens og
Kraftens dialektiske Eenhed.
3 Materiens og Kraftens dialektiske Eenhed sees tillige
den extensive Størrelses dialektiske Eenhed med den intensive. Idet den extensive Størrelse $M$ indsvinder til $dM$ og
den intensive $P$ til $dP$, ere de begge i deres Anlæg identiske;
idet $dM$ integreres til $M$ og $dP$ til $P$, ere de begge
i deres Virkelighed identiske. Her er i den virkelige Kugle,
hvad Relationen mellem dens Masse og Vægt angaaer, ikke
en intensiv Størrelse for sig, forbunden med en extensiv,
men en intensiv Størrelse, hvis Realitet bestaaer i det Extensive,
en extensiv Størrelse, hvis Realitet er sat i det Intensive.
„Man kan,“ hedder det om Tyngdepunktet,\footnote[1]{Naturlærens mechaniske Deel af H. C. Ørsted. S. 28.}

forestille sig Tingene, som om hele Legemets Vægt var
forenet i Tyngdepunktet, og alle de øvrige Dele ingen Vægt
havde.“ Ja, saaledes kan man ganske vist „ „forestille“
sig Tingene; men saaledes kan man ikke \emph{tænke} sig Tingene.

I den extensive Størrelses dialektiske Eenhed med den
intensive sees Eenheden af det Udeelbare med det Deelbare,
det Enkelte med det Sammensatte. Forsaavidt det Intensive sættes i Forestillingen mod det Extensive, vil hiint

% --- p. 130 ---
\setcounter{footnote}{0}%
\opage{130}naturligviis repræsentere det Udeelbare, dette det Deelbare,
hiint det Enkelte, dette det Sammensatte. At det just er
en saadan forestillet Modsætning, hvortil Leibniß og Kant,
hver paa sin Maade, holde sig, er let at paavise. Hvad Leibniz kræver af sine »points metaphysiques“, at de
nemlig skulle være exacte som de mathematiske Punkter, og
reelle som de physiske, passer udelukkende paa de i Virkeligheden
intensive Punkter; et saadant intensivt Punkt er nu
just, for at blive ved det valgte Exempel, en Kugles Tyngdepunkt.
Naar paa den anden Side Kant, efter at han i
„Thesis“ har gjort samme Fordring til det Enkelte som
Leibniz, opstiller til Modsætning den Paastand, at ingen
sammensat Ting bestaaer af enkelte Dele, at der i det Hele
ikke existerer noget Enkelt: da gjør han i denne sin „Antithesis“
en meget eensidig Brug af den extensive Størrelse.
Misforstaaelsen hos Leibniz og Kant beroer, trods den store
Forskjel, der ellers er i deres philosophiske Synsmaade,
hvad det i Delingsproblemet Substantielle angaaer, væsentlig
paa det Samme. Det Sammensatte maa, derom ere begge
enige, nødvendigviis forudsætte det absolut Enkelte, det
Deelbare det absolut Udeelbare; men denne Antagelse er
hverken overeensstemmende med rationel Mechanik eller med
Tingenes Natur.

I den intensive Eenhed bestaaer den reelle Enkelthed;
men da det Intensive, som ovenfor er viist, har sin Realitet
i det Extensive, og da enhver endelig Kraft er en Resultant
af Composanter, hvis materielle Angrebspunkter ved
at være ved Siden af hverandre afgive det Extensive: saa
vil Atomet selv, forsaavidt det indtager et Rum, have Dele,
og disse Dele udgjøre materielle Angrebspunkter for de
Composanter, hvis Resultant er Atomets endelige Kraft,
dets intensive Eenhed. Antages Atomet nu dog alligevel at
være et i sig Udeelbart, et absolut Enkelt: da er det ikke et
Enkelt, fordi det ingen Dele har, men fordi der tillægges
Delenes Sammenhængskraft en saa stor Intensitet, at ingen
% --- p. 131 ---
\setcounter{footnote}{0}%
\opage{131}Magt i Tilværelsen formaaer at skille dem ad; med andre
Ord: Atomet er da et absolut Udeelbart, forsaavidt dets
Deles Sammenføining er absolut, Atomet er et absolut
Enkelt, forsaavidt det er et absolut Sammensat; men et
 absolut Sammensat er en Modsigelse, og et absolut Atom
en Fiction. Den Forestilling, at et absolut Enkelt nødvendigviis
skulde ligge til Grund for det Sammensatte, er
da ogsaa en reen Indbildning; thi hvad er vel det absolut
Enkelte? Det, som ifølge sit Begreb ikke indgaaer, og ikke
kan indgaae i nogen endelig Sammensætning! At den rationelle Mechanik ligesaa exact kan decomponere Kraften i
det Uendelige, som dele Massen i det Uendelige, medens
det dog er det i Kraften Intensive, der giver Enkeltheden
Realitet, viser, hvad desuden ligger i Sagens Natur, at,
ligesom enhver endelig Sammensætning er relativ, maa
ogsaa enhver tilsvarende Enkelthed være relativ. Paa det
Enkeltes Relativitet beroer det Sammensattes Realitet; men
det Enkeltes Relativitet beroer igjen, naar den sees under
Størrelsens og navnlig under den extensive Størrelses
Synspunkt, paa dets uendelige Deelbarhed. Det Sammensatte
har altsaa, ifølge de naturlige Overgange, sin Realitet
i det Enkelte. derved, at det Enkelte igjen har sin Realitet i
det Sammensatte: den uendelige Dialektik af det Enkelte og
det Sammensatte er Delingsproblemets Løsning.

\begin{center}b) Bevægelsesproblemet.\end{center}

Hvor inderlig dette Problem hænger sammen med det
foregaaende, sees af Eleaternes herhen hørende Indsigelser.
Imod Muligheden af en Bevægelse i Rummet har Zeno
opstillet fire Beviser, der alle grunde sig paa Rummets
uendelige Deelbarhed. „Dersom Rummet eller den Materie,
som opfylder Rummet, er uendelig, saa at ingen Deel er den
mindste, da er Bevægelsen umulig, fordi den forudsætter, at et
Legeme skrider frem fra et Sted til et andet, og altsaa gjennemløber
et Rum. Men dette er umuligt, efterdi 9* ethvert Rum
% --- p. 132 ---
\setcounter{footnote}{0}%
\opage{132}er uendeligt, og følgelig ikke kan giennemtrænges i en given
endelig Tid. Eller kortere: Uendeligheden af det Rum, der
skal gjennemløbes, er i Strid med Tidens Endelighed, uden
hvilken ingen Bevægelse er mulig“\footnote[1]{Tennemann. Geschichte der Philos. 1ster Band. Leipzig 1798. S. 197.}

Den Omhu, hvormed Aristoteles i sine physiske Skrifter
søger at afvæbne Zenos Indvendinger, giver os ikke blot et
Vidnesbyrd om hans Skarpsindighed, men minder tillige om
den Braad, der er skjult i de zenoniske Grunde. Zenos første
Beviis, hvortil der allerede ovenfor (S. 35) er taget Hensyn,
lader sig fremstille saaledes: „Naar et Legeme bevæger
sig, saa gjennemløber det en Linie fra det ene Punkt til det
andet: det maa altsaa, førend det kan komme til Endepunktet,
først være kommet til Midtpunktet. Da nu ethvert
Rum er deleligt i det Uendelige, saa bestaaer enhver af Liniens
Dele ligesaavel af uendelig mange Dele, som den
hele Linie; Legemet kan under sin Bevægelse altsaa aldrig
naae til Liniens Midtpunkt, endsige til dens Endepunkt,
d. e. Bevægelsen kan slet ikke tænkes.“

Forsaavidt Nerven i Beviset sættes deri, at et endeligt
Rum, som ved at deles i det Uendelige bliver uendeligt,
skulde, for at gjennemløbes, fordre, ikke en endelig,
men en uendelig Tid, har Aristoteles træffende bemærket, at
Tiden ligesaavel er deelbar i det Uendelige, som Rummet,
og at Zeno netop har feilet deri, at han har anseet Tiden
for anderledes begrændset end Rummet\footnote[2]{Ζήνων οὖν παραλογιζόμενος, ἄλλως τὸ τοῦ χρόνου πεπερασμένον λαμβάνει, καὶ ἄλλως τὸ τοῦ σταδίου . . . Jvfr. Tennemann. Anfr. Skr. 3ter Band. S. 151.} Herved er imidlertid,
og det vistnok med fuld Ret, den Betænkelighed anført, at, „om ogsaa Zeno har feilet i at betragte Tiden paa
en anden Maade end Rummet, saa har Aristoteles ved at
blotte denne Bevisets formelle Svaghed dog i Realiteten ikke
rokket Beviset selv; den Vanskelighed, at en Bevægelse i

% --- p. 133 ---
\setcounter{footnote}{0}%
\opage{133}Rummet kun kan tænkes som en i det uendelige sig tabende
Overgang fra Punkt til Punkt, at det Legeme, der skulde
bevæge sig, altsaa først maatte blive færdigt med en Uendelighed
af uendelig smaae Bevægelser,\footnote[1]{„Weil aber diese Theile unendlich sind, so muß er sich eine endlose Zeit bewegen, ehe er sie durchläuft, er kommt also nie aus den letzten.“ Tennemann. Anfr. Skr. 3ter Band. S. 152.} før en endelig Bevægelse
kunde opnaaes, staaer endnu ved Magt, om man
endog tænker sig Tiden deelt paa samme Maade som Rummet.“
Denne Vanskelighed, der da heller ikke hæves ved en
kantisk Adskillelse „zwischen Ding an sich und Erscheinung“,
hører blandt de slaaende Exempler, der vise Nødvendigheden
af at lade den exacte Analyse afgive Begrebsudviklingens
Forudsætninger. Det Blændende i Zenos Beviis svarer til
en Illusion, der ligesaavel gaaer igjennem den kantiske\footnote[2]{„Denn die Bewegung des Räumlichen ist so gut als eine Theilung, weil sie nicht anders kann gedacht werden, als wie ein Fortrücken von einer Stelle zur andern. Und da nun Aristoteles das Wesen jeder stetigen Größe darin setzt, daß jeder Theil derselben wieder aus Theilen bestehet (nicht etwa von uns blos so vorgestellt wird), so muß das Bewegliche unendliche Theile durchlaufen, ehe es von einer Stelle zur andern kommen kann; das ist aber so viel als, es kann nie von einer Stelle zur andern übergehen. Und gerade das ist es, was Zeno wollte.“ Tennemann. Anfr. Værk. 3ter Band. S. 153.}
som igjennem den aristoteliske Tænkning, den Illusion nemlig,
at det uendelige skulde falde ved Siden af det Endelige,
være dets Negation uden at være dets Affirmation, at man
altsaa, forsaavidt man var i en Uendelighed, skulde befinde
sig udenfor Endeligheden. At denne Forudsætning er grundfalsk,
har det Uendeliges Analyse godtgjort, idet den har
godtgjort, at den endelige Størrelse, væsentlig opfattet som
et bestemt Integral, netop er en Determination af $0 . \infty$, at man altsaa ved at bevæge sig i den bestemte Uendelighed
bevæger sig midt i en Endelighed. Forsaavidt Endeligt maales
med Endeligt, er Bevægelsen en uafbrudt Succession, en

% --- p. 134 ---
\setcounter{footnote}{0}%
\opage{134}stadig Overgang fra et Mindre til et Større: først det
Mindre, saa det Større, først en Tomme, saa et Qvarteer,
først den halve Vei, saa den hele; men en saadan Succession,
der i sin Continuitet kun er mulig formedelst det Endeliges
dialektiske Eenhed med det Uendelige, gjælder ikke Forholdet
mellem Uendeligt og Endeligt indbyrdes: Bevægelsen fordrer
ikke først et Uendeligt, thi idet et Uendeligt sættes, er det
Endelige allerede forudsat\footnote[1]{„I hvilken Forstand det Endelige kan siges at være samtidigt med det Uendelige, oplyser den rationelle Mechanik fra mangfoldige Steder. „D'Alemberts Princip anvendes ikke blot paa continuerligt virkende Kræfter, comparable med Tyngden, men ogsaa paa de saakaldte \emph{instantane} Kræfter, som frembringe pludselige Forandringer i Hastighed, f. Ex. ved faste Legemers Sammenstød. Denne anden Art af Kræfter er strængt taget ikke forskjellig fra den første; thi i Naturen vil ingen Forandring i Hastighed foregaae pludseligen, men Hastigheden maa, om end i et nok saa lille og umærkeligt Tidsinterval continuerligen gjennemløbe alle Værdier fra den oprindelige, som var tilstede, idet Kraften begyndte at virke, indtil den Slutningshastighed, som opnaaes til samme Tid, som Kraften ophører. Maalet for disse Kræfter haves følgelig paa Grund af det samme Princip, hvorefter f. Ex. Tyngden maales, nemlig efter Formlen $g=\frac{v}{t}$, hvilket ogsaa gjælder for en variabel Kraft $\Phi$ i Henseende til dens Middelværdie i en vis Tid $t$, nemlg $\frac{\int_{0}^{t}\Phi dt}{t}=\frac{v}{t}$.“ C. Ramus. Analytisk Mechanik. Kbhvn. 1852. S. 211. Vi komme da ifølge denne Betragtningsmaade, til at skjelne imellem: $\alpha$) et absolut Nu, det abstract-logiske Dieblik, Grændsen mellem Fortid og Fremtid, et Nul, som ikke kan integreres, $\beta$) et relativt Nu, det mathematiske Dieblik, Tidselementet $dt$, som kan integreres, og $\gamma$) et empirisk Dieblik, Stødets virkelige Nu, den umærkelig lille Tidsdeel. I det absolute Nu er Intet, og flere Intet. I det relative Nu sættes det Uendelige som det Endeliges Begyndelse, dets Element og Anlæg, og derhos tillige som en Uendeligheds Begyndelse: det Endeliges Begyndelse, falder alt- saa sammen med, er i det mathematiske Dieblik samtidig med det Uendeliges Begyndelse. I det empiriske Dieblik sættes det Endelige som fuldendt; men i og med denne det Endeliges Fuldendelse er en Uendelighed selv fuldendt: det Endeliges Fuldendelse falder altsaa sammen med, er i hvert virkeligt Dieblik samtidig med en Uendeligheds Fuldendelse.} Skulde Bevægelsen, som Zeno

% --- p. 135 ---
\setcounter{footnote}{0}%
\opage{135}foregiver, udtømme det Uendelige, før den kunde naae det
Endelige: ja, det vilde medtage en evig Tid, det vilde være
et uendeligt Sinkeri! men, naar Bevægelsen ligesaa oprindelig
begynder med det Endelige, som med det Uendelige,
naar det Endelige i Bevægelsen altid er, ikke bagefter eller
forud for, men \emph{samtidigt} med det Uendelige, fordi det
ved at forudsætte det uendelige selv er det uendeliges Forudsætning:
da er Dæmningen brudt, da har Bevægelsen frit
Løb, og Zeno er afvæbnet.

Imod Zenos andet Beviis (Achilles) anbringer Aristoteles
i Grunden, paa en lille Ændring nær, den samme
Indvending, som han allerede har brugt imod det første.
Beviset, der skal godtgjøre, at Achilles ikke kan naae Skildpadden,
d. v. s., at et Legeme, der antages at bevæge sig
med stor Hastighed, ikke kan indhente et andet Legeme, der
i nogen Afstand tænkes at bevæge sig med en endog meget
ringe Hastighed, har ligesom det Foregaaende den Fordeel,
at det holder sig i det Punktuelle. Man sætte t. Ex., at
to Legemer, $A$ og $B$, bevæge sig paa samme Linie, saaledes
at $A$, der har en ti Gange større Hastighed end $B$, gjennemløber
den første Halvdeel af Linien, idet $B$ gjennemløber
Tiendedelen af den anden Halvdeel; før $A$ kan naae $B$, maa
det da ogsaa gjennemløbe denne Tiendedeel, men medens $A$ gjennemløber benævnte Tiendedeel, gjennemløber $B$ igjen en
Tiendedeel af den følgende Tiendedeel, o. s. fr. Uagtet $B.s$ Forspring hver Gang bliver mindre og mindre, faaer det
ved sin Bevægelse dog et lille Forspring; dette Raisonnement
kan nu fortsættes i det Uendelige, og i det uendelige
give os en Tiendedeel af en Tiendedeel af en Tiendedeel:
% the note 1) of p. 134 runs over onto the foot of this page: set in full at its mark
% --- p. 136 ---
\setcounter{footnote}{0}%
\opage{136}Phantasien virker, Forestillingen forvirres, og Tanken svimler.
Imod en saa punktuel Fremstilling ere almeenlogiske
Bemærkninger uden Virkning: at sige et Par passende Ord
om Continuitet og Discretion, om stadig Formindskelse
o. s. v., forslaaer ikke; skal Blændværket hæves, maa det
Punktlige sættes mod det Punktlige, og den exacte Analyse
anbringes. Vælge vi da den Tid, $A$ vil være om at
gjennemløbe Strækningen $l$, Afstanden imellem $A's$ og $B's$ Udgangspunkt, til Tidseenhed, og kalde $A's$ Udgangspunkt $u$, danne sig følgende Rækker:

I Tiden $0$ har $A$ en Afstand fra $u=0$, $B$ en Afstand
fra $u=l$; i Tiden $1$ er $A's$ Afstand fra $u=l$, $B's$ Afstand
fra $u=(1+\frac{1}{10})l$; i Tiden $1+\frac{1}{10}$ er Afstanden
for $A=(1+\frac{1}{10})l$, for $B=(1+\frac{1}{10}+(\frac{1}{10})^{2})l$, i Tiden $1+\frac{1}{10}+(\frac{1}{10})^{2}$ for $A=(1+\frac{1}{10}+(\frac{1}{10})^{2})l$, for $B=(1+\frac{1}{10}+(\frac{1}{10})^{2}+(\frac{1}{10})^{3})l$; i Tiden $1+\frac{1}{10}\ldots+(\frac{1}{10})^{n}$ for $A=(1+\frac{1}{10}\ldots+(\frac{1}{10})^{n})l$, for $B=(1+\frac{1}{10}\ldots+(\frac{1}{10})^{n}+(\frac{1}{10})^{n+1})l$. Spørges nu: naar vil $A$ naae $B$?
da følger det af sig selv, at dette først vil kunne skee, naar
følgende Ligning er tilstillet: $(1+\frac{1}{10}+(\frac{1}{10})^{2}\ldots\ldots+(\frac{1}{10})^{n})l=(1+\frac{1}{10}+(\frac{1}{10})^{2}\ldots\ldots+(\frac{1}{10})^{n}+(\frac{1}{10})^{n+1})l$, det vil sige, naar $(\frac{1}{10})^{n+1}=0$, altsaa, naar $n=\infty$.
Den hertil svarende Tid: $1+\frac{1}{10}+(\frac{1}{10})^{2}\ldots+\mathrm{Lim}(\frac{1}{10})^{n}=\frac{1}{1-\frac{1}{10}}=\frac{10}{9}=1\frac{1}{9}$ er, som man seer endelig, og $A$ naaer
altsaa $B$ i en endelig Tid. Er den valgte Tidseenhed t. Ex. $1$ Sec., vil $A$ behøve $1\frac{1}{9}$ Sec. for at naae $B$; de uendelig
mange mindre og mindre Mellemrum, der svare til de
uendelig mange mindre og mindre Tidsdele, ville da i Løbet
af $1\frac{1}{9}$ Sec., idet de som det Endeliges Elementer Punkt
for Punkt resultere af det Endelige, give det endelig bestemte
Resultat.

Erkjendelsen af den Maade, hvorpaa det Uendelige her
gaaer op i det Endelige, hviler paa den Sætning, at en
Række med utallige Led kan give en endelig Sum, og er

% --- p. 137 ---
\setcounter{footnote}{0}%
\opage{137}forsaavidt betinget af Indsigt i uendelige Rækkers Convergens.
At den in abstracto uopnaaelige Grændse just
bliver opnaaet derved, at en uendelig Række ikke blot i
Tanken, men i Virkeligheden føres til Ende, at en Uendelighed
virkelig gjennemløbes i og med det Endelige, er altsaa den Dialektik, hvormed Illusionen ophører.

Det tredie Beviis (den flyvende Piil) er holdt i en
mere almindelig Form, og lader sig fremsætte saaledes\footnote[1]{Tennemann. Anfr. Skr. 1ster Band. S. 198—199.}: „Skulde et Legeme bevæge sig, saa vilde deraf følge, at det
paa een Gang maatte være i Bevægelse og i Hvile; thi
ethvert sig bevægende Legeme maa i ethvert Dieblik være i
et Rum, der netop er ligestort med det selv. Men saaledes
er Legemet i Hvile; thi hvad der, om end kun for et Dieblik,
befinder sig i et tilsvarende lige Rum, er med Hensyn
paa dette Rum i Hvile. Tænke vi os en Piil, der med
største Hurtighed gjennemflyver et Rum, saa maa den i hvert
Dieblik være i et andet Rum, og forsaavidt den indtager
dette hvert Dieblik være i Hvile. En saadan flyvende Piil
maatte da tænkes at være paa een Gang baade i Hvile og
i Bevægelse, hvilket er en Modsigelse.“

Hvad Aristoteles i denne Anledning har udviklet angaaende de continuerlige Størrelser, Rummet og Tiden
og begges uendelige Deelbarhed\footnote[2]{Physic. \textit{VI, 1.}}, angaaende Forskjellen
mellem Tidsgrændsen, Nuet og Tidselementet\footnote[3]{Physic. \textit{VI, 3.} „Die Unmöglichkeit aber, daß in dem Augenblicke als solchem etwas sich bewege, oder, als Gegensatz der Bewegung, ruhe, ergiebt sich aus den Begriffen der Schnelligkeit und Langsamkeit, welcher Gegensatz als wesentlich jeder Bewegung zukommend bereits vorhin gesetzt worden war; sodann aber auch aus der Definition des Augenblicks selbst, in welchem, dafern er zugleich der Vergangenheit und der Zukunft angehört, so oft die Ruhe in Bewegung, und umgekehrt, übergeht, Ruhe, und Bewegung in Bezug auf dasselbe Ding zugleich statt finden müßte, was widersinnig wäre.“ C. H. Weisse. Aristoteles Physik. Leipzig 1829. Anmerkung zum 6ten Buche. S. 593.}, angaaende

% --- p. 138 ---
\setcounter{footnote}{0}%
\opage{138}Forskjellen mellem en Hvile, der udelukker Bevægelse, og
den Hvilens uendelige Forhindring, der udgjør Bevægelsen
selv\footnote[1]{Charakteristisk er især den Maade, hvorpaa 8de Kapitel (Phys. VI) ender: „Men da Alt, hvad der bevæger sig, bevæger sig i en Tid, og gaaer fra Noget over i Noget (ἐκ τινὸς εἰς τὶ μεταβάλλει), kan det sig Bevægende da for det Første ikke i nogen Henseende hvile, hvad den Tid, hvori det ganske (ἐν ᾧ χρόνῳ κινεῖται καθ᾽ αὐτό) og ikke deelviis (καὶ μὴ τῶν ἐν ἐκείνου τινί) bevæger sig, angaaer. Thi Hvilen bestaaer deri, at saavel Tingen selv, som enhver af dens Dele er en Tidlang i Eet og det Samme. Kun da nemlig tale vi om Hvile, naar det med Sandhed kan siges, at Tingen med samt dens Dele er i forskjellige Dieblikke efter hinanden (ἐν ἄλλῳ καὶ ἄλλῳ τῷ νῦν) paa eet og samme Sted. Men naar Hvilen bestaaer heri, saa formaaer det sig Forandrende ikke at være heelt og holdent i Noget, hvad den første Tid angaaer. Thi al Tid er deelbar: saa at man med Sandhed kunde sige, at det Foranderlige var i de forskjellige Tidsdele paa samme Sted, det selv saavel som dets Dele. Forholder det sig derimod ikke saaledes, er det sig Bevægende kun i et eneste Nu (ἐν ἑνὶ μόνῳ τῷ νῦν) her eller der: saa vil det ikke hvad nogensomhelst Tid, men kun hvad en Tids-Grændse, angaaer være i Noget. I hvert Nu vil det vistnok altid befinde sig i Noget, men ikke saaledes, at det hviler (ἐν δὲ τῷ νῦν ἔστι μὲν ἀεὶ κατά τι μένον οὐ μέντοι ἠρεμεῖ). I det rene Nu kan hverken Bevægelse eller Hvile finde Sted, men derimod Ikke-Bevægelse og Ikke-Hvile. Men i Tiden kan det sig Bevægende derimod ikke forholde sig paa samme Maade, som det Hvilende; thi da vilde jo deraf følge, at det sig i Rummet Bevægende tillige hvilede.“}, hører vistnok til det Betydeligste, der endnu er
præsteret i Naturphilosophien; ligesom det jo ogsaa maatte
følge af Sagens Natur, at den Kategorie, der hos Aristoteles
er Physikens Grund-Kategorie, nemlig Bevægelsen
(κίνησις), vilde hos en saa alsidig og udholdende Tænker
finde en alsidig og, saavidt muligt, udtømmende Behandling.

Men, ihvorvel den aristoteliske Skarpsindighed fra mange
Sider trænger ind i Sagen, har dog en tilfredsstillende
% the note 3) of p. 137 runs over onto the foot of this page: set in full at its mark

% --- p. 139 ---
\setcounter{footnote}{0}%
\opage{139}Knudens Løsning været umulig, og det af dobbelt Grund.
For det Første er Savnet af exacte Analyser paafaldende;
hvad her og der anføres, for at gjøre Forholdet imellem
Stedforandring og Tidsforandring, det Hurtigere og det
Langsommere o. desl. tydeligt, kan, da Reflexionen ikke
understøttes af tilsvarende Operation, ikke komme ud over
den logiske Almindelighed, og Relationerne derfor heller ikke
articuleres i den Grad, at en punklig Overeensstemmelse af % PRINTED AS IS: punklig (for punktlig)
Begreb og Virkelighed skulde kunne vise sig: Adskillelsen af
Tidselementet og det rene Nu er som en Raisonnementets
Udflugt, en subtil Vending, en forblommet Talemaade forbleven
i Dunkelhed\footnote[1]{„Es ist schwer sich einen bestimmten Begriff von dem, was Aristoteles eigentlich unter diesen Augenblicken gedacht hat, zu bilden. Keine Theile der Zeit und doch Grenzpuncte? Sind es nur Ruhepuncte für das Vorstellende, wenn es eine Reihe von Bewegungen durchläuft, und sie vorstellt. So scheint es. Wenn er aber behauptet, daß die Grenze nur dem angehöre, dessen Grenze es ist (τὰ μὲν γὰρ πέρατα ἐκείνου μόνον ἐστίν, οὐ ἐστι πέρατα. \textit{Physie. IV. c. 11}), so muß man annehmen, daß auch der Augenblick etwas in der Zeit sey. Aber dieses läßt sich selbst nicht mit dem Begriff des Aristoteles von der Zeit vereinigen, daß sie ein Ganzes, dessen Theil jedesmal ein Zeittheil und ins Unendliche in Zeittheile getheilt werde, eine Folgerung, die selbst dann noch gültig ist, wenn er unter den Augenblicken bloß imaginäre Abschnitte sich gedacht hätte. Tennemann. Gescht. d. Phil. 3ten Band. S. 154—155.}. % PRINTED AS IS: Physie. (for Physic.) in note 1

Den anden ligesaa væsentlige Hindring er Mangel af
gjennemgaaende Dialektik. Vel er Aristoteles formeget Tænker
og formeget Græker til, at han ikke overalt i Begrebernes
Analyse skulde støde paa det Dialektiske\footnote[2]{„Uebrigens sind die Ergebnisse dieses Capitels“ (Talen er om \textit{Physic. VI, 6)} „ein Beispiel, wie die scharf und kräftig durchgeführte aristotelische Verfahrweise der abstract-verständigen Sonderung und Scheidung zuletzt auf ächt speculative Begriffbestimmungen führt: wie denn dieser Gegensatz von Stetigkeit der Veränderung und Plötzlichkeit — (Platon soll, wie Simplicius be- richtet, den Ausdruck ἐξαίφνης für den Begriff des Augenblicks gebraucht haben) — des durch die Veränderung hervorgerufenen seienden Zustandes, eine Knotenlinie von zeitlich-räumlichen oder andern zwischen Qualität und Quantität oscillirenden Maßverhältnissen, ganz unstreitig eine ächt speculative und dialektische, über die gemeine Reflexion weit hinausgehende und an den wichtigsten Folgen reiche Bestimmung ist.“ C. H. Weisse. Anfr. Skr. S. 602.}; men Tingen er,

% --- p. 140 ---
\setcounter{footnote}{0}%
\opage{140}at Aristoteles, istedetfor at holde fast paa Platos guddommelige
Tanke, at Philosophien videnskabelig er Dialektik, og
føre Tanken videre, tvertimod, hvad dette Punkt angaaer,
affalder fra Ideen, nedsætter Dialektiken til et Underordnet,
og maa desaarsag, hvor det kniber, hjælpe sig ved endelig
Raisonneren frem og tilbage. Intetsteds i den hele aristoteliske
Physik er vel Savnet af en exact Analyses Bistand
mere føleligt end netop i de Afsnit, hvor Aristoteles benytter
„Dieblikket“, og sigter paa den zenoniske Sphinx; og
intetsteds er Savnet af en gjennemført Dialektik mere føleligt
end i den Reflexion, hvormed den dogmatiske Tænker undviger
Modsigelsen, og saa gaaer „den flyvende Piil“ imøde.

Endskjøndt Bevægelsen ingenlunde er, saaledes som Zeno
har fremstillet Sagen, den Modsigelse, at det sig Bevægende
skulde paa een Gang være i en Bevægelse udelukkende Hvile
og en Hvile udelukkende Bevægelse, en Modsigelse, der i sin
Uopløselighed vilde være et umiskjendeligt Tegn paa Nonsens;
saa godtgjør dog Zenos Beviis, at der i enhver
Bevægelse alligevel er, ihvormeget Aristoteles end søger at
skjule det og snoe sig fra det, en Modsigelse, der i det
Uendelige sættes og opløses; at Kategorien: Bevægelse,
med andre Ord, ligesom Kategorierne: Vorden\footnote[1]{„Vorden er nemlig den Relation, som det Absolute selv er. Derfor alternerer det mellem Væren og Ikke-Væren, og er selv den fortsatte Række af begges Succession; thi hvad som vorder eller bliver til, hører deels til det Værende, forsaavidt som det allerede er Noget, deels til det Ikke-Værende, forsaavidt som det har en Fremtid, der endnu ikke er opfyldt; og begge Momenter danne en fortsat Række, fordi deres Afvexling er continueerlig, thi først naar Betragtningen standser i et eller andet bestemt Punkt, viser det Vordende sig som en Adskillelse af Momenterne, idet først da den ene Deel af det Vordende indbefattes under det Værende, og den anden under det Ikke-Værende.“ J. L. Heiberg. Perseus. Nr. 2. S. 23.} og For-

% the note 2) of p. 139 runs over onto the foot of this page: set in full at its mark

% --- p. 141 ---
\setcounter{footnote}{0}%
\opage{141}andring, er uendelig dialektisk. For at denne Dialektik
imidlertid ikke ved at løsrives fra det Punktuelle skal faae
Udseende af et Tankespil, hvori den virkelige Bevægelse forvandles
til tom Ordbevægelse, ville vi sikkre Udviklingen
ved exacte Forudsætninger, og desaarsag først tage den
aristoteliske Behandling af Zenos fjerde Beviis med.

Det fjerde Beviis, hvorved Zeno angriber Bevægelsen,
er, dersom dets overleverede Form ellers er authentisk, en
saa iøinefaldende Sophisme, eller maaskee rettere: Paralogisme,
at Aristoteles ikke saameget behøvede at gjendrive,
som at afvise det, hvad han da ogsaa har gjort\footnote[1]{„Men det fjerde er det om de lige Masser, der bevæge sig frem i Banen mod hinanden fra modsatte Sider, den ene fra Banens Ende, den anden fra dens Midte: hvoraf skal følge, at den dobbelte Tid bliver lig den halve. Men Feilslutningen (ὁ παραλογισμός) ligger i den Fordring, at det Ene i Relation til det Hvilende skulle i lige Tid med lige Hastighed gjennemløbe lige Strækning. Dette er falsk (τοῦτο δὲ ψεῦδος).“ Physic. \textit{VI, 9.}} Tendensen
i Beviset er ikke desto mindre afgjørende nok. Lod
det sig virkelig bevise, at enhver Hastighedsforskjel forsvinder,
naar Bevægelsen tænkes, da var Bevægelsens Utænkelighed
med det samme beviist. Den herunder hørende Modsigelse
er Aristoteles neppe bleven opmærksom paa: han forudsætter
uden videre en Forskjel af Hastigt og Langsomt, idet
han hævder Bevægelsen, og forudsætter saa igjen Bevægelsen,
idet han hævder Forskjellen mellem Hastigt og Langsomt\footnote[2]{Jvfr. \textit{Phys. VI, 1, 2, 3 og VI, 9.}}

Ved et Indblik i den exacte Analyse bliver det først
muligt at see, hvad der ligger i Problemet. Lad t. Ex.
Bevægelsen paa to Baner, $A$ og $B$, foregaae saaledes, at

% the note 1) of p. 140 runs over onto the foot of this page: set in full at its mark

% --- p. 142 ---
\setcounter{footnote}{0}%
\opage{142}medens der paa $A$ tilbagelægges een Miil i Timen, tilbagelægges
der paa $B$ otte Miil i samme Tidseenhed. Et ringe
Kjendskab til Differential- og Integralregningens første Sætninger
er da nok for at forstaae, at naar paa Banen $A$ i endelig
Tid Hastigheden
\[ v=\frac{s}{t}=1 \]
, saa ogsaa Dieblikshastigheden
\[ \frac{ds}{dt}=1 \]
, og at paa Banen $B$, hvor Hastigheden

\[ v_{1}=\frac{s_{1}}{t}=8 \]
i endelig Tid, maa Dieblikshastigheden ligeledes
være
\[ v_{1}=\frac{ds_{1}}{dt}=8 \]
. Men saa let som det nu er, efterat
det uendeliges Analyse eengang er indført, at finde sig tilrette
med den herhen hørende Operationsmechanik —
det gaaer uden Hovedbrud ved en Smule Routine, som af sig
selv —, saa umuligt har det været, og er det endnu for den
med Analysen ubekjendte Reflexion at magte Opgaven og
værge sig mod den Sophistik, som ved at benytte „Dieblikket“
til at udjævne enhver Forskjel i Hastighed vil afskaffe
Hastigheden selv, og gjøre Bevægelsen meningsløs.

I den uendelig lille Tidsdeel er, naar man holder
sig til et almindeligt Raisonnement, det Stykke Vei, der
gjennemløbes paa $B$, ligesaavel uendelig lille, som det, der
gjennemløbes paa $A$, $ds_{1}$ altsaa ligesaa vel uendelig lille
som $ds$. „Men er $ds<ds_{1}$, hvordan kan $ds_{1}$ da være
uendelig lille? Hvorledes er det muligt, hvorledes er det
tænkeligt, at Noget skulde kunne være endnu mindre end
det, som allerede er uendelig lille? Hvorledes er det muligt,
hvorledes er det tænkeligt, at Noget skulde være større end
det, som allerede er uendelig stort? Mathematikeren siger
det! Ja, men Mathematikeren mener det ikke alvorligt med
det uendelige, han sætter i Tanken enten en tom Negation
eller dog kun det Endelige. Mathematikeren betegner det
Uendelige; men han tænker ikke Uendelighedens Tanke: det
er ham i det Hele ikke saameget om Begreber som om Be%

% --- p. 143 ---
\setcounter{footnote}{0}%
\opage{143}tegnelser at giøre, han tænker kun paa at faae sine Betegnelser
saaledes indrettede, at Beregningen kan slaae til.“

Er Reflexionen først heret ind i den fixe Forestilling,
at det uendelige absolut maa være ligestort med sig selv:
vil Modsætningen mellem Hastighed i endelig Tid og Dieblikshastighed
let blive saa haard, at den kan optræde som
kantisk Antinomie. Thesis: Der \emph{maa nødvendigviis}
gives forskjellige Øieblikshastigheder, efterdi
der i Virkeligheden gives forskjellige Hastigheder.
Antithesis: Da enhver Forskjel i den rene
Dieblikshastighed er utænkelig, ere de saakaldte
virkelige Hastighedsforskjelligheder umulige.
For at bevise Thesis behøver man kun at gaae
ud fra den endelige Forskjel mellem
\[ \frac{s}{t} \]
og
\[ \frac{s_{1}}{t} \]
, og
derfra slutte til en tilsvarende Forskjel mellem
\[ \frac{ds}{dt} \]
og
\[ \frac{ds_{1}}{dt} \]
. For at bevise Antithesis begynde man med at sætte
\[ ds=ds_{1} \]
altsaa
\[ \frac{ds}{dt}=\frac{ds_{1}}{dt} \]
, og uddrage saa heraf Slntningen:
\[ \frac{s}{t}=\frac{s_{1}}{t} \]
. At den ene uendelig lille Størrelse kan være, ikke
blot større, men endog uendelig mange Gange større end den
anden, den ene uendelig store Størrelse, ikke blot mindre,
men endog uendelig mange Gange mindre end den anden,
er virkelig et Paradox saa haardt og anstødeligt, at den
saakaldte sunde Meneskeforstand snart vilde blive færdig med
det, dersom ikke den exacte Analyse fra saa mangfoldige
Sider fandt sig opfordret til at indskjærpe det, og derved
kom til at opretholde dets Anseelse. Men „den sunde Forstand“,
der er saa enig med sig selv i at forkaste den Synsmaade,
at der skulde kunne være Forskjel i det uendelig
Lille og det uendelig Store, og at det uendelige paa Grund
heraf ikke skulde kunne indgaae under Begrebet: Størrelse,

% --- p. 144 ---
\setcounter{footnote}{0}%
\opage{144}er tillige enig med sig selv i at forkaste den Paastand, at
der ikke skulde være en virkelig Forskjel imellem langsommere
og hurtigere Bevægelse. At den raisonnerende Forstand ved
saaledes at være enig med sig selv, netop er i Splid med
sig selv, har den i sin umiddelbare Suffisance naturligviis
ingen Forestilling om; men at Tingen ikke desto mindre
forholder sig saaledes, er let at vise. Forskjellen mellem
langsom og hurtig Bevægelse maa aabenbart beroe derpaa,
at den første har mere, den anden mindre Hvile i sig, at
den første altsaa er Hvilen nærmere, end den anden; en
Bevægelse med uendelig lille Hastighed er jo netop Hvilen
saa uendelig nær, at den kan betragtes som Hvile. Ikke
destomindre er Modsætningen af Hvile og Bevægelse dog
qualitativ: hvad der hviler, bevæger sig ikke, hvad der bevæger
sig, hviler ikke. Bortkaste vi nu de uendelige Størrelser,
kan Hastighedsforskjellen kun forklares deraf, at den rene
Bevægelse endelig afbrydes af Hvilen. At Legemet paa $A$ i Tidseenheden $t$ kun gjennemløber $s$, medens Legemet paa $B$ gjennemløber $s_{1}=8s$, maa altsaa forstaaes saaledes, at
hiint Legeme kun bevæger sig $i \frac{1}{8} t$, men derimod hviler $i \frac{7}{8} t$. Saalænge Tidseenheden, hvori en saadan Afvexling foregaaer,
er 1 Time, er Hvilens Afvexling med Bevægelsen
iøinefaldende; er Tidseenheden 1 Secund, er Afvexlingen
mindre paafaldende, og tager sig ud som en stødviis Bevægelse;
i en Tidseenhed af $\frac{1}{100}$ Secund vil Afvexlingen slet
ikke mærkes, men Phænomenet vise sig som en jævn Bevægelse.
At der paa $A$ kun gjennemløbes $1$ Miil, medens
der paa $B$ gjennemløbes $8$ Miil, vil da, naar Hvilens
Afvexling med Bevægelsen skeer i hver Hundredeel af Secundet
eller i Tiden $\tau$, tage sig ud for Anskuelsen, som om
Legemet paa $A$, trods dets ringere Hastighed, ligesaavel var
i jævn og uafbrudt Bevægelse som Legemet paa $B$; medens
Forklaringen dog er, at Bevægelsen paa $A$ idelig afbrydes

% --- p. 145 ---
\setcounter{footnote}{0}%
\opage{145}af Hvilen, idet her til hver Bevægelse i Tidsdelen $\frac{1}{8}\tau$ svarer en Hvile i Tidsdelen $\frac{7}{8}\tau$. Skal Bevægelsen paa $A$ være continuerlig, maae Tiden og Rummet være deelbare
i det Uendelige; i dette Punkt har Aristoteles seet rigtigt.
Men hvad Aristoteles ikke har seet og ikke kunnet see, fordi
han manglede det Støttepunkt, som kun det Uendeliges
Analyse kan give, er, at her ved denne Tidens og Rummets
uendelige Deelbarhed netop danne sig uendelig smaae
Størrelser med en saadan indbyrdes Forskjel, at de indgaae
i bestemte Endelighedsforhold med hinanden.

Zenos Indvending, at det Legeme, som skulde være i
Bevægelse, maatte paa een Gang være baade i Hvile og
Bevægelse, har altsaa det Sande i sig, at enhver Bevægelse
forudsætter en Hvile, som uophørlig er i Begreb med at
indtræde, og som idelig ophæves. Det er i denne Sammenhæng
indlysende, at Hegel har Ret, naar han anfører
„Ophævelsen“ blandt de dialektiske Kategorier. Vistnok er
Hvilen ogsaa en ophævet Bevægelse, men denne Ophævelse
er forsaavidt udialektisk, som Hvilen, endelig seet, ikke har
den ophævede Bevægelse i sig, men udenfor sig; i Bevægelsen
derimod er Hvilens Ophævelse, der skeer i hvert Tidselement,
dialektisk, thi den ophævede Hvile er med i Bevægelsen selv,
saa at Forskjellen mellem større og mindre Hastighed netop
beroer paa den større eller mindre Energie, hvormed Hvilen
ophæves: det Extensive gaaer ogsaa her op i det Intensive.
Med Hvilens dialektiske Ophævelse ere da Zenos Indvendinger
imod Bevægelsen tillige ophævede.

Vilde Nogen imidlertid gaae videre, og fordre Problemet
løst saaledes, at Bevægelsen i Rummet blev opfattet
og fastholdt i et Begreb, der udelukkede Anskuelsen: maatte
vel en saadan Fordring siges at beroe paa en reen Misforstaaelse.
Rum og Tid ere, for at minde om Kant, ikke
rene Forstandsbegreber, men aprioriske Anskuelser. At ville
% --- p. 146 ---
\setcounter{footnote}{0}%
\opage{146}forvandle Anskuelsen til et abstract Forstandsbegreb, at ville
opfatte det, som kun er for Anskuelsen, i et Begreb, der
tilintetgjør Anskuelsen, er naturligviis meningsløst; saadanne
Spørgsmaal, som: kan Rummet tænkes? kan Tiden tænkes?
kan Bevægelsen tænkes? maae da i denne Forstand besvares
ligefrem med Nei! Havde Zeno uden videre opstillet det som Fordring,
at Den, der vilde tænke Bevægelsen, skulde tænke den uden
Anskuelse, og udenfor Anskuelsen: da havde den i en saadan
Fordring liggende Modsigelse vist været for iøinefaldende til
at afgive et gjennem Aarhundreder gaaende Problem. Nei,
Zeno forudsætter, at Rum, Tid og Bevægelse ere givne
som Anskuelser, og fordrer saa, at disse Anskuelser tillige
skulle lade sig opfatte i tilsvarende Begreber. Saalænge
Begreb og Auskuelse ere i Strid, saa at det, der fremstiller % PRINTED AS IS: Auskuelse (for Anskuelse)
sig for Anskuelsen som en Virkelighed, viser sig for Tanken
som en Modsigelse og derved som en Umulighed, staaer
Problemet uopløst; naar det derimod fra alle Sider er viist,
at det i Bevægelsen, der umiddelbart fremstiller sig for Anskuelsen,
tillige, forsaavidt det sees under Reflexionens Synspunkt,
tilfredsstiller Tanken, har Problemet fundet sin Løsning.

\begin{center}c) Idealitetsproblemet.\end{center}

Bevægelsens Dialektik er en Dialektik af Begreb og Anskuelse,
og just derved en Dialektik af det uendelige og det Endelige. At
opløse det Endelige i det Uendelige saaledes, at det igjen fremgaaer
deraf, er med andre Ord at paavise det Endeliges Idealitet.
Da saadanne Videnskaber som Astronomie og Physik, hvis
„ideale Princip“, for at bruge et Udtryk af J. L. Heiberg\footnote[1]{See „\emph{Perseus}“, Journal for den speculative Idee. Nr. 1. Til Læseren.},
„er Mathematik“, nødvendigviis maae afgive Grundlaget,
naar den kosmologiske Idee med sit reelle Indhold\footnote[2]{Jvfr. Phil. Propæd. S. 184.} skal
udvikles og fremstilles, maa Philosophien ved mangesidig

% --- p. 147 ---
\setcounter{footnote}{0}%
\opage{147}Berøring med disse Videnskaber paa mangfoldige Maader
see sig opfordret til at gjøre Brug af den høiere Analyse.
Ideelæren kan nemlig ikke vinde reelt Indhold derved, at
Abstractionens tomme Former udfyldes, for ikke at sige
udstoppes med Exempler og Læresætninger, laante fra Astronomie
og Physik; derimod kunne selv de meest concrete
Phænomener gjennemklares af Idealitetens Lys, naar Ueudeligheden % PRINTED AS IS: Ueudeligheden (for Uendeligheden)
faaer en saadan Magt med det Endelige, at en
Stoffets Forbrænding kan foregaae i Dialektikens Ildluft
(πῦρ ἀεὶ ζῶον ἁπτόμενον μέτρα καὶ ἀποσβεννύμενον μέτρα).

En Idealisme, der egentlig kun bekymrer sig om det i
endelig og materiel Forstand Virkelige, forsaavidt den overalt
lægger an paa at opløse dette Virkelige som Configuration
af det almene Begreb, som et i den rene Idee forsvindende
Moment, har i Grunden ingen væsentlig Trang
til at indlade sig med de exacte Analyser. I Hegels Logik
staaer da ogsaa Behandlingen af det Mathematiske, henviist
til et Par Anmærkninger, i en temmelig løs Forbindelse
med Qvantitetskategoriernes dialektiske Udvikling; at den
hegelske Methode conseqvent maa føre bort fra det Mathematiske,
idet den fører bort fra de physiske Virkelighedsanalyser,
viser hans Naturphilosophie, og navnlig hans
Forklaring af Faldloven.\footnote[1]{Jvfr. Phil. Propæd. S. 162—163.}

Den Philosophie derimod, som, for at erkjende Tingenes
Væsen og hærde Ideens Gyldighed, maa binde an
med Kjendsgjerninger, og ved at optrævle Virkelighedens
Knuder søge de ideelle Traade, har uafviselig Trang ikke
blot til exacte Analyser, men til Realbegreber, der oprindelig
tilhøre Videnskaber, „hvis ideale Princip er Mathematik.“
Men er Trangen væsentlig, da er ogsaa Fordringen
uafviselig; hvorvidt det Mathematiske skal tages med i
Philosophien, og de exacte Begreber inddrages i Dialektikens
Strøm, vil altsaa ikke beroe paa den Philosopherendes Vil-

% --- p. 148 ---
\setcounter{footnote}{0}%
\opage{148}kaarlighed, men paa Erkjendelsens Gang, paa Totalitetens
Anlæg, paa Sagudviklingens nødvendige Methode\footnote[1]{Jvfr. Phil. Propæd. S. 111—117.}.

„Videnskabens Historie“, siger Ramus\footnote[2]{Differential- og Integralregning. Fortale.}, „fremstiller
det Factum, at de største Opdagelser i Physik og Astronomie,
som saa kraftigen have indgrebet i Menneskehedens
baade aandelige og sociale Liv, have været deels betingede
af deels paa det Nøiagtigste forbundne med Opdagelserne af
den mathematiske Analyses meest abstracte Sandheder. Er
den høiere Mathematiks Studium langt og vanskeligt, saa
afgive igjen dets Resultater en desto skjønnere Løn. Frugtesløse
ere derimod de Philosophers Bestræbelser, som ville
afkorte Veien ved mystiske og metaphysiske Speculationer.“
Men uagtet Mathematiken nu ganske vist er imod „de Philosophers
Bestræbelser“, som ved „mystiske Speculationer“
ville gjøre saavel de exacte Analyser som Erfaringsvidenskabens
Jagttagelser overflødige: er Mathematiken paa dens
nuværende Trin dog alligevel den Videnskab, som, navnlig
fra een Hovedside, fortrinlig egner sig til at corrigere
„Speculationerne“, og fremhjælpe den videnskabelige Metaphysik.
Vistnok finder den chemiske Physiks Atomlære
paa en vis Maade sin Stadfæstelse og Betingelserne for
sin conseqvente Udvikling i Mathematiken; men Forestillingen
om evige, i sig forhærdede Atomer (\textit{indivisibilia
absoluta}) lader sig umuligt bringe i Samklang med en
Analyse, der deler Stoffet i det Uendelige, og forvandler
Stofdelene til Functionspunkter. Den Sætning, at ethvert
Element væsentlig er et Moment, er en Hovedsætning, som
beviser, i hvilken Grad de exacte Analyser maae kunne bidrage
til at give Virkelighedsidealismen det rette aprioriske Grundlag,
naar de Begreber, der oprindelig udvikles i Mathematiken,
opfattes og bearbeides af Dialektiken.

% --- p. (back matter) ---
% ===== BACK MATTER: Indhold (PDF 163-164) and Rettelser (PDF 165), read at the image; no folios printed =====
\clearpage
\begin{center}{\large Indhold.}\end{center}
{\small\noindent\hfill Side.\\
\par\medskip\centerline{I.}\par\noindent
\par\medskip\centerline{Functionsbegrebet.}\par\noindent
a) Functionsbegrebets oprindelige Reqvisiter \dotfill 4--12.\\
$\alpha$) En constant og en variabel Størrelse. $\beta$) To variable\\
Størrelser, en afhængig og en uafhængig. $\gamma$) Grundforhold\\
mellem de to Variable.\\
b) Functionen og Loven \dotfill 12--20.\\
$\alpha$) Tilfældet. $\beta$) Reglen. $\gamma$) Loven.\\
c) Functionsbegrebets eiendommelige Fremstilling \dotfill 20--27.\\
$\alpha$) Sproget. $\beta$) Tanken. $\gamma$) Sproget og Tanken.
\par\medskip\centerline{II.}\par\noindent
\par\medskip\centerline{Uendelighedsproblemet.}\par\noindent
a) Det Endelige i dets umiddelbare Forskjel fra det\\
Uendelige \dotfill 28--33.\\
$\alpha$) Det Uendelige udenfor det Endelige. $\beta$) Det Uendelige\\
indenfor det Endelige. $\gamma$) Det Endeliges og\\
det Uendeliges dialektiske Eenhed.\\
b) Det Endeliges Overgang i det Uendelige: Tilnærmelsen. \dotfill 34--38.\\
$\alpha$) Et Opnaaeligt: Tilnærmelse af Endeligt til Endeligt.\\
$\beta$) Et Uopnaaeligt: Tilnærmelse af Endeligt\\
til Uendeligt. $\gamma$) Den uendelige Tilnærmelse.\\
c) Tilnærmelsen i punktuel Udvikling \dotfill 38--56.\\
$\alpha$) Mathematisk-logiske Overgange. $\alpha\alpha$) Analysens og\\
Synthesens exacte Anvendelse: Operation. $\beta\beta$) Analyse\\
paa synthetisk Grundlag: Transformation.\\
$\gamma\gamma$) Synthese paa analytisk Grundlag: Combination.\\
$\beta$) Functionssystemer: den indre Uendelighed.\\
$\alpha\alpha$) Functionssystemernes Grundformel: Reduction.\\
$\beta\beta$) Særskilte Functionssystemer: Deduction. $\gamma\gamma$) Tilvæxter\\
og Rester: Articulation.
\par\medskip\centerline{III.}\par\noindent
\par\medskip\centerline{Det Uendeliges Analyse.}\par\noindent
a) Functionsgrændsen: det Endeliges dialektiske Berøring\\
med det Uendelige \dotfill 57--76.\\
$\alpha$) Uendelige Rækkers Afbrydelse: Grændsen som uopnaaelig.\\
$\beta$) Rækkernes væsentlige Afslutning:\\
Grændsen som opnaaelig. $\gamma$) Uendelige Størrelser:\\
Begrændsningens Dialektik.\\
b) Differentialet: Momentets Dialektik \dotfill 78--97.\\
$\alpha$) Differential og Differentialcoefficient: Forholdsstørrelse.\\
$\beta$) Differentialofficienternes forskjellige Orden: System\\
af Forholdsstørrelser. $\gamma$) Differentiationsprincipper:\\
den ved Forholdsconseqvenserne bestemte Methode.\\
c) Integralet: Elementet og den endelige Størrelse \dotfill 98--115.\\
$\alpha$) Integrerende Summation: Moment og Element.\\
$\beta$) Integrationsprincipper: afledede Functioners\\
Forhold til de oprindelige. $\gamma$) Det bestemte Integral:\\
Uendelighedsproblemets Løsning.
\par\medskip\centerline{Slutning.}\par\noindent
a) Delingsproblemet \dotfill 116--131.\\
b) Bevægelsesproblemet \dotfill 131--146.\\
c) Idealitetsproblemet \dotfill 146--148\\
}
\clearpage
\begin{center}{\large Rettelser.}\end{center}
{\small\noindent
S. 38. Linie 9 og 14 fr. ov. Assymptote, læs: Asymptote.\\
- 40. --- 13 --- analythisk --- analytisk.\\
- 43. --- 23 fr. ned. retvinklene --- retvinklede.\\
- 45. --- 2 fr. ned. $A$ --- $A_1$ og\\
$\frac{1}{a}$ --- $\frac{1}{A_1}$.\\
- 46. Anm. $l(a)^{n-1}$ --- $(l.a)^{n-1}$.\\
- 50. Anm. Linie 1 --- impiriske --- empiriske.\\
- 55. Linie 1 fr. ov. Da nemlig --- Er nemlig.\\
- 59. --- 5 --- Reduction --- Reduction,\\
6 --- hianden --- hinanden.\\
- 67. --- 13 --- nte Led --- mte Led.\\
- 75. --- 12 --- $(a - x)\,l^{\frac{1}{a-x}}$ --- $(a - x)\,e^{\frac{1}{a-x}}$.\\
- 81. --- 1 --- bive --- blive.\\
- 90. --- 9 fr. ned. Benouilli --- Bernoulli.\\
- 94. --- 9 fr. ov. Differentialcoesficient, læs: Differentialcoefficient.\\
- 94. --- 7 fr. ned. heeel, læs: heel.\\
- 96. --- 9 og 14 fr. ov. $\frac{dx}{1}, \frac{dx^2}{1.2}$ --- $\frac{dx}{1.2}, \frac{dx^2}{1.2.3}$ og\\
10 --- $dx$ --- $\frac{dx}{1.2}$.\\
- 101. --- fr. ov. er Delene --- er Delen.\\
- 105. --- 5 og 4 fr. ned. --- Lim $\left(1+\frac{1}{n}\right)^{x^2}$\\
$= x^2$.\\
- 110. --- 3 --- $\int \frac{'x^m\,dx}{l(x)^n}$ --- $\int \frac{x^m\,dx}{(l.x)^n}$.\\
2 --- $\int \frac{'x^m\,dx}{l\,x^{n-1}}$ --- $\int \frac{'x^m\,dx}{(l.x)^{n-1}}$.\\
sidste --- $\int \frac{'x^m\,dx}{l.x}$ --- $\int \frac{'x^m\,dx}{l.x}$, for n positiv heel,\\
}

\end{document}
